\section{C. The descent equation and the congruence coordinates}

\begin{statement}{C.1}{Proved}{Descent, extension, and the exact fixed-point equation}
Let $\mathcal P$ have the meaning in the question. Descent along
$\varphi:F\twoheadrightarrow\widehat B$, followed by restriction to
$K_2\simeq F$, defines a continuous homomorphism
\[
 D:\mathcal P\longrightarrow\widehat{\mathrm{GT}}_0
\]
preserving the multiplier. In particular its restriction to
$\widehat{\mathrm{GT}}_0$ is an endomorphism. With the gauges specified
in the question,
\[
 C^\varphi=\operatorname{Fix}(D|_{\widehat{\mathrm{GT}}_0}).
\]
This is a closed subgroup containing the specified Galois image.
It is the literal pair of compatibilities of \S47, scan 240,
formula (30), after a gauge is chosen as in the proof.
The braid identity in Lead 9 is correct with the stated standard
normalization of the braid action.
\end{statement}
\begin{proof}
The kernel of $\varphi$ is the closed normal subgroup generated by
$x^2$ and $z^3$. Raising a topological generator of a procyclic group
to a unit power leaves that subgroup unchanged. Consequently every
pure peripheral automorphism preserves this kernel, and so does its
inverse. Descent therefore gives an automorphism, not merely an
endomorphism, of $\widehat B$.

There is just one kernel of a surjection $\widehat B\to S_3$:
the images of $s,r$ must be a transposition and a three-cycle,
and the six possibilities form one orbit under $\Aut(S_3)$.
Thus $K_2$ is characteristic. Every automorphism of $S_3$ is inner,
and its centre is trivial. We may consequently adjust the descended
automorphism by an inner automorphism of $\widehat B$ to make its
quotient action on $S_3$ the identity; the remaining ambiguity is
exactly an inner automorphism from $K_2$. Its restriction is a
well-defined outer automorphism of $K_2$. The three cusp classes
of this index-six covering are the three conjugacy classes of
lifts of $e^2$. Identity on $S_3$ fixes these classes individually,
and restriction has the same multiplier. Conjugation by any lift
of an element of $S_3$ commutes with this restriction modulo
$\Inn(K_2)$. This proves that the image lies in the symmetric
peripheral group. All operations descend from continuous operations
on compact automorphism groups, so the homomorphism is continuous.

For precision about extension, the conjugation homomorphism
$\widehat B\to\Aut(K_2)$ is injective. Its kernel centralizes $K_2$;
intersection with $K_2$ is its centre, which is trivial, and an
element in another coset would induce a trivial permutation of its
three cusp classes. Its image is the inverse image of the geometric
$S_3$ under $\Aut(K_2)\to\Out(K_2)$. Indeed any automorphism with
outer class in that $S_3$ differs from the chosen deck transformation
by an inner automorphism of $K_2$. Conjugating this inverse image by
an automorphism of $K_2$ centralizing $S_3$ gives the extension of a
gentle outer class. Changing the representative changes the extension
by an inner automorphism. Restriction and extension are inverse on
these outer classes. Hence $D(F)=F$ says exactly that descent and
extension agree as outer automorphisms of $\widehat B$.

Here is the promised normalization check. Put $e=sr$,
$x=(rs)^2$, $y=e^2$, and $\nu=(\lambda-1)/2$. For the cusp-normalized
extension $u_0$ write
\[
 u_0(e)=e^\lambda,\qquad u_0(r)=\xi r,\qquad
 u_0(s)=\eta s,\qquad \eta=\xi x^{-\nu}.
\]
Its restriction has $u_0(y)=y^\lambda$ and
$u_0(x)=\xi x^\lambda\xi^{-1}$. The normalized pair on $F$ is
obtained by the inner adjustment $\Int(\xi^{-1})$, and
$\xi=y^\nu f$. Its descent sends $s$ to $s$ and $e$ to
$\varphi(f)^{-1}e^\lambda\varphi(f)$. Conjugating this descent by
$\varphi(\xi)=e^\nu\varphi(f)$ gives agreement on $e$ automatically.
Agreement on $s$ is precisely
\[
 \eta=\varphi(\xi)s\varphi(\xi)^{-1}s^{-1},\qquad
 \xi=\bigl(\varphi(\xi)\sigma(\varphi(\xi))^{-1}\bigr)x^\nu,
 \quad \sigma=\Int(s).
\]
These are the equation and its equivalent normalized form in
\S47, scans 234 and 240, formulas (15) and (30).
If the outer actions agree, their conjugating element has the form
$b=e^t\varphi(\xi)$: the discrepancy centralizes $e^\lambda$, whose
centralizer is $\langle e\rangle$. The gauge
$\xi_a=y^a\xi x^{-a}$ changes $u_0$ to $\Int(e^{2a})u_0$
and the normalized pair to its inner translate by $x^a$.
Accordingly
\[
 b_a=e^{2a}bs^{-a}=e^{a+t}\varphi(\xi_a).
\]
Taking $a=-t$ gives the displayed literal equation. This proves
both directions of the fixed-point assertion, including its gauge.
The geometric maps defining descent are defined over $\mathbb Q$,
so the specified Galois image is fixed, by functoriality of the
arithmetic fundamental group with the same inner ambiguities.

Finally put $u=\sigma_1$, $v=\sigma_2$ in the profinite braid group
and $c_3=(uvu)^2=(uv)^3$. In the standard action
$u\mapsto u^\lambda$, $v\mapsto f(u^2,v^2)^{-1}v^\lambda f(u^2,v^2)$,
the first hexagon, with $m=(\lambda-1)/2$, is
\[
 u^\lambda f^{-1}v^\lambda f
       =f^{-1}uv\,u^{-2m}c_3^m.
\]
Multiplication on the right by $u^\lambda$ gives
$F(uvu)=f^{-1}uvu\,c_3^m=f^{-1}(uvu)^\lambda$.
At multiplier one the two-cycle equation identifies conjugation
by $uvu$ with interchange of $u^2,v^2$, and says that $f$ is the
corresponding nonabelian cocycle. The displayed equation of
\S47 identifies its primitive with its own image under $\varphi$,
with the inverse convention dictated by $y\mapsto f^{-1}yf$.
\end{proof}

\begin{statement}{C.2}{Open}{The harmonic comparison is not established by the factorization alone}
The explicit equivalence requested in Task C.2 between
$C^\varphi$, the harmonic and equianharmonic relations, and
Nakamura--Schneps relation (III$'$) is unresolved here.
C.25--C.26 prove the exact, fully normalized equivalence
with (GF$_0$) and (HF$_0$), including the unique parameter
extending the Kummer cocycle of $3$. They identify the
additional equations (V) and (IV$'_{\rm bis}$) needed
for the harmonic and equianharmonic equations.
The rational factorization is correct:
\[
 w=4\lambda(1-\lambda),\qquad
 j=64\frac{(4-w)^3}{w^2}
   =256\frac{(1-\lambda+\lambda^2)^3}{\lambda^2(1-\lambda)^2}.
\]
It does not identify the required based fundamental-group maps.
\end{statement}
\begin{proof}[Established comparison and precise missing step]
Substitution proves the identity. The opened Nakamura--Tsunogai note
\cite{PBC_NT2002} studies the four-punctured harmonic and
equianharmonic configurations and their immersion and covering maps.
The opened paper \cite{PBC_NS2000} defines an additional subgroup by
relations (III$'$) and (IV); its introduction does not identify this
subgroup with all of $\widehat{\mathrm{GT}}$. Neither assertion supplies
a substitution of the actual generators and tangential paths in
C.1 into both displayed relations. The harmonic primitive
substitutions are now supplied by C.25--C.26, with their
Kummer powers and central braid ambiguities resolved.
The remaining comparison requires the extra equations
(V) and (IV$'_{\rm bis}$), identification with the separately
defined Ihara cocycles, and the four-strand identities.
The rational factorization alone supplies none of these steps.
A further precise obstruction to one proposed geometric argument
is given in C.15 below.
\end{proof}

\begin{statement}{C.3}{Open}{The pentagon-to-descent implication}
It remains undecided here whether
$\widehat{\mathrm{GT}}\subset C^\varphi$.
\end{statement}
\begin{proof}[Attempt and obstruction]
C.1 makes this an equality between the restriction--filling operation
and the identity on pairs satisfying the pentagon. The pentagon gives
compatibility with the forgetful maps of the five-point configuration
space. The orbifold quotient $F\to C_2\amalg C_3$ is not one of those
forgetful maps. No verified comparison of these two operations has
been obtained. The next finite test must retain the two evaluations
$\iota(\xi)$ and $\varphi(\xi)$ and the torsion-frame data together;
the calculation in C.14 retains the two evaluations at level four
but has no obstruction there.
\end{proof}

\begin{statement}{C.4}{Proved}{The fixed subgroup of descent is proper}
There is a strict inclusion
\[
                 C^\varphi\varsubsetneqq\widehat{\mathrm{GT}}_0.
\]
C.41 strengthens this to index $2^{\aleph_0}$ and Haar measure zero,
and gives large odd prime-power quotients after killing its closed
normal closure. The stronger assertion
$C^\varphi\subset\ker\Delta$ remains unresolved here.
\end{statement}
\begin{proof}
C.23 constructs an actual full gentle element $e_U$ with
$\Delta(e_U)=1$ and $\Delta(D(e_U))=0$.
It therefore is not fixed by $D$. The construction refines
R24 E.58 to retain the second evaluation and both of its
torsion frames. It does not assert that every descent-fixed
element has parity zero.
\end{proof}

\begin{statement}{C.5}{Open}{Simultaneous pro-prime images of descent}
For every $\ell\geq5$, C.40 determines the full simultaneous
gentle image in $(F^{(\ell)},D(F)^{(\ell)})$: it is the whole
common-multiplier fibre product of the two normalized actual
local gentle groups, simultaneously over all these primes.
The corresponding image with $\widehat{\mathrm{GT}}$ in place
of the gentle group remains undetermined. This is the remaining
Open part of Task C.5. The strict full profinite inclusion
required in its last clause is proved in P.4.
\end{statement}
\begin{proof}[Established gentle image and remaining pentagon obstacle]
C.40 retains the ordinary local quotient and all six actual
deck-conjugate fillings in one quotient of the modular source.
Its seven free pro-prime factors have vanishing central
homology after the actual finite deck action. The entire
source-centralization is a pullback from a fixed finite
two-three quotient, so all three pure peripheral frames lift.
The finite refinement argument then gives the actual full
gentle element with both prescribed local coordinates.
This supplies the previously missing joint lifting step;
it does not treat the relative source as an absolute
pronilpotent quotient. A first coordinate in the image of
$\widehat{\mathrm{GT}}$ must satisfy the pentagon, but neither
this necessary condition nor the gentle lifting construction
determines the pair image of that subgroup. The remaining
attack must retain the pentagon in the full source lift.
\end{proof}

\begin{statement}{C.6}{Open}{Actual finite images and the least detector of parity}
The exact images of the full groups at the characteristic level-four
envelope, and the least finite quotient detecting $\Delta$, remain
undetermined. The necessary finite cycle-automorphism locus at
level four is computed completely in C.14. Every one of its elements
passes the descent equation in some gauge.
\end{statement}
\begin{proof}[Why the computation does not claim more]
A homomorphism of a finite characteristic quotient satisfying the
finite equations need not extend to a compatible system on every
finite characteristic quotient. Conversely every full gentle pair
has one of the computed residues. Thus C.14 proves a necessary
condition on actual images and proves that this necessary test
cannot distinguish Tasks C.3 and C.4. The parity uses pure lifts of
the order-two and order-three frames; their uniqueness holds in the
full group, not in the small matrix quotient. A finite detector must
retain enough of the source to determine those frame images. The
explicit calculation of that detector has not been completed.
\end{proof}

\begin{statement}{C.7}{Open}{The CM image is determined, with descent and manuscript comparisons still open}
C.20 determines the simultaneous adelic image of the two
normalized CM coordinates on the multiplier-one gentle group
and on its parity kernel. C.21 reduces the image on
$C^\varphi$ to exactly two possibilities. Choosing between
them, and the remaining manuscript-coordinate identifications
in Task C.7, remain unresolved. The field intersection in
Lead 25 is proved in C.13; its full image assertion, with
the modulo-two conditions retained, is proved in C.20.
The character and its vanishing on the Galois image are
the checked statements R24 E.65, with its supplied sign corrections.
\end{statement}
\begin{proof}[Established information and the gap]
The finite difference character is $\Delta$, and R24
supplies an actual full element with value one. C.20
combines the exact individual CM images with the field
intersection: the arithmetic image has index two, so
that one additional element proves full surjectivity.
This argument preserves the single global signs and
the modulo-two conditions; the field intersection
alone would not suffice without the individual image
calculation. The still missing step for $C^\varphi$
is whether its fixed-point equation forces the finite
difference character to vanish. The additional equality with the
manuscript's scalar of scan 11 requires the normalization in scan 12
and has not been inferred from the shared letter used for it.
\end{proof}

\begin{statement}{C.8}{Open}{Parity on descent and on the pentagon group}
The assertions $C^\varphi\subset\ker\Delta$ and
$\Delta|_{\widehat{\mathrm{GT}}}=0$ remain unresolved here.
C.27 refutes $C^\varphi=\ker\Delta$: an actual full
element has $\Delta=0$ and $\Delta D=1$, so it is not
fixed by $D$.
The assertions about closed AA subgroups in this task are settled
in C.16. The natural Galois maps to $\mathcal P$ and to
$\widehat{\mathrm{GT}}_0$ are not surjective.
\end{statement}
\begin{proof}[Established negative conclusion and unresolved implication]
P.4 proves $\widehat{\mathrm{GT}}\subsetneq
\widehat{\mathrm{GT}}_0$, while the natural Galois image is contained
in $\widehat{\mathrm{GT}}$. This proves non-surjectivity of both
natural maps, and hence refutes the indicated natural isomorphisms
in cote 143, scans 29--31, at the tripod. It makes no assertion
that their underlying abstract profinite groups cannot be isomorphic
by some unrelated map. Alternatively the R24 element with parity
one supplies the same non-surjectivity. The multiplier-one
four-strand covariance computed in R24 E.61 is identically one on
that element as well, so that particular projection of the pentagon
cannot prove parity zero. C.25--C.26 now give the exact
primitive comparison, but the additional equations (V),
(IV$'_{\rm bis}$), and (III$'$) still cannot be inferred
from fixedness. C.27 settles the proposed equality
negatively without deciding the remaining inclusion.
C.41 further determines the index and Haar measure of the fixed
subgroup and constructs characters vanishing on its closed normal
closure; these difference characters do not decide either remaining
parity assertion.
\end{proof}

\begin{statement}{C.9}{Open}{Character presentations beyond the quadratic rank}
C.30 determines the abstract group of continuous quadratic characters of
$\widehat{\mathrm{GT}}_0$ trivial on the Galois image: it has countably
infinite dimension, and the corresponding maximal exponent-two quotient
is topologically $\prod_{n\geq1}C_2$. It supplies an explicit infinite
independent family which vanishes even on the pentagon group. Thus
generation by $\Delta$ is false. C.32 further determines the entire
normal quotient of the actual metabelian image: it is a countable
product of $C_2$, for either the Galois or pentagon normal closure.
C.33 constructs a full arithmetic-trivial character of order
$13$, nonzero on multiplier one and invisible to that metabelian
quotient.
C.34--C.36 give independent odd-primary arithmetic-trivial
characters at every prime $p\geq13$, with unbounded explicit
rank lower bounds.
C.37 supplies the remaining primes $5,7,11$, with rank
lower bounds $1,1,513$.  Thus every prime at least five occurs.
C.38 strengthens this to exactly countably infinite character
dimension at each prime $p\geq5$, with a simultaneous quotient
$\prod_{p\geq5}\prod_{k,j\geq1}\mathbb Z/p^k\mathbb Z$.
In particular arbitrarily large finite $p$-power orders occur.
C.35 proves that the order-$13$ character also vanishes on the
full pentagon group; C.39 supplies such pentagon-trivial
characters also at $5$ and $7$.
C.41 gives the same simultaneous odd prime-power product as a
quotient after killing the entire descent-fixed subgroup, together
with a quadratic factor, by actual difference characters.
C.31 proves that $\Delta,\Delta D,\Delta D^2$ are three
independent characters, already jointly onto $C_2^3$ on the
multiplier-one subgroup.
A complete explicit basis, the full profinite abelianization of the quotient by the closed normal Galois
image, and the finite-order arithmetic-trivial characters of
$\widehat{\mathrm{GT}}$ remain undetermined.
\end{statement}
\begin{proof}[Exact remaining scope]
Let $N_G$ be the closed normal closure of the specified Galois image
in the gentle group. The quadratic characters are
$\Hom_{\mathrm{cts}}((\widehat{\mathrm{GT}}_0/N_G)^{\mathrm{ab}},
\mathbb F_2)$, now determined abstractly by C.30. That result does not
specify higher two-power characters or all odd-order characters and
their relations in the full abelianization. C.33 supplies an actual
order-$13$ character, and C.34--C.36 supply odd-primary quotients
at every prime at least thirteen, while C.37 covers $5,7,11$.
C.38 proves unbounded prime-power orders and arbitrary finite
independent families at every prime at least five; it does not
identify all their relations or imply a $\mathbb Z_p$ quotient.
The prime-three character space remains undetermined.
C.32 shows that every such higher
or odd-order character must use the image of the kernel of
the metabelian action: every finite image factoring through the
metabelian arithmetic-normal quotient is elementary abelian of
exponent two. These results do not supply the requested level-24 and
level-60 frame calculations. C.30's beta-defect characters
restrict to zero on the whole pentagon group, so they do not resolve
the corresponding question for that group. C.27 separately supplies
the source-compatible descended character $\Delta D$. C.31 completes
the next actual central-lifting calculation and, with C.28, proves
independence of $\Delta,\Delta D,\Delta D^2$. Independence of all
further iterates remains open. An explicit basis for all characters
requires a comparison between these families and the remaining
finite quotients.
\end{proof}

\begin{statement}{C.10}{Open}{Resolvent maps and the exceptional index-prime actions}
The index-prime classification and all three signatures are
proved in C.17, and both degree-five resolvents, their exact
domain and their arithmetic field are constructed in C.18.
C.29 supplies the degree-seven and degree-eleven families,
their exact finite stabilizer domains, and the fields of
all their marked fillings. The comparisons with pentagon
fixedness and (III$'$), the values on the specified
E.58/D.46 lifts, and the common fixed subgroup remain open.
The literal first
sentence of Lead 27(a), with no restriction $p\geq5$, is false:
$\PSL_2(\mathbb F_2)=S_3$ has an index-two subgroup and
$\PSL_2(\mathbb F_3)=A_4$ has an index-three subgroup.
\end{statement}
\begin{proof}[Counterexamples and missing descent data]
The subgroups are $A_3\subset S_3$ and the normal Klein four-group
in $A_4$. For the exceptional actions at $7,11$, restriction to
an open subgroup and filling chosen factors only defines an outer
action on the resulting tripod when the original action preserves
the subgroup and the retained factor classes. The domain must
include these stabilizer conditions. Establishing that a proposed
quadratic multiplier condition imposes all of them requires the
actual arithmetic dessin and its field of definition. The abstract
index and signature calculation alone does not give that field,
and does not give fixedness under every pentagon pair. C.18
and C.29 now supply the geometric and domain data at all
three primes. Their values on the E.58 and D.46 elements,
their arbitrary pentagon fixedness, and the common fixed
subgroup remain unfinished. The next attack is to restrict the retained
full frames to the explicit $A_4$ cover before filling,
rather than infer fixedness from the degree or multiplier.
\end{proof}

\begin{statement}{C.11}{Open}{Latt\`es compatibility does not yet characterize the frame image}
The equivalences in Task C.11 between commuting with the two
Latt\`es families and lying in the corresponding arithmetic frame
images remain unresolved individually. Their conjunction with
the assertion that commuting with both forces zero parity is
refuted in C.22: individual arithmeticity of either coordinate
holds for every positive gentle element. No deduction of
$\Delta=0$ for all elements commuting with both families is made.
\end{statement}
\begin{proof}[What construction supplies and what it does not]
For a CM elliptic curve and an endomorphism $[\alpha]$ defined over
its CM field, $[\alpha]$ commutes with the finite unit group.
It therefore descends to a rational self-map of the quotient
$E/\mu_6$ or $E/\mu_4$. Multiplication is unramified on $E$, so
all ramification of the descended map is over the orbifold cone
points. This gives the stated algebraic Latt\`es construction,
and its being defined over the CM field gives Galois equivariance
of composition on dessins with the chosen quotient marking.
It does not identify every informal subdivision of a triangulation
with that composition, nor show that commuting with the whole
family imposes exactly the CM coordinate condition. The latter
requires a reconstruction of the based torsion tower and its
comparison character, which is the remaining implication.
\end{proof}

\begin{statement}{C.12}{Proved}{The full gentle action preserves the congruence tower}
Use the canonical braid action of $F\in\widehat{\mathrm{GT}}_0$
and its quotient action $\Phi_F$ on $\widehat B$.
Every projective congruence kernel $\widehat\Gamma(n)$ is invariant,
and the conjugacy class of $\widehat\Gamma_0(n)$ is invariant.
There is a continuous homomorphism
\[
 F\longmapsto\left[\mathcal V_F\right]\in
       \PGL_2(\widehat{\mathbb Z}),\qquad
 \mathcal V_F=\begin{pmatrix}\lambda(F)&t(F)\\0&1\end{pmatrix},
\]
which implements the congruence action. Its indicated representative
has determinant exactly $\lambda(F)$; the determinant of its
projective class is only defined modulo squares. Its kernel is
explicitly
\[
 \{(1,f):f(U^2,L^2)=I\text{ in }
                \SL_2(\widehat{\mathbb Z})\},
 \quad
 U=\begin{pmatrix}1&1\\0&1\end{pmatrix},\quad
 L=\begin{pmatrix}1&0\\-1&1\end{pmatrix}.
\]
This gives a membership criterion for the congruence ghosts;
it does not claim a presentation or generators for that kernel.
\end{statement}
\begin{proof}
The two-cycle and three-cycle relations give an actual automorphism
of the profinite braid group on $u,v$, as in C.1. Evaluate its
images under $u\mapsto U$, $v\mapsto L$ over each domain
$\mathbb Z_p$. Write $A=U^\lambda$ and $B=f^{-1}L^\lambda f$.
Then $B$ has trace two and determinant one, and $ABA=BAB$.
Solving the four entries of this relation gives either $B=A$,
or, for a uniquely determined $a\in\mathbb Z_p$,
\[
 B=\begin{pmatrix}1-a&\lambda a^2\\
                  -1/\lambda&1+a\end{pmatrix}.
\]
For completeness, write $B=\left(\begin{smallmatrix}b&c\\d&2-b\end{smallmatrix}\right)$.
The lower-left equation gives $d(1+\lambda d)=0$.
In the branch $d=0$, the other equations and determinant give
$b=1,c=\lambda$, hence $B=A$. In the other branch
$d=-1/\lambda$ and the determinant gives
$c=\lambda(b-1)^2$; put $a=1-b$.
The first branch is impossible: precomposition by the actual braid
automorphism preserves the nonabelian image of the two braid
generators, whereas $A=B$ would make it abelian. Thus
\[
 \mathcal V_F=\begin{pmatrix}\lambda&\lambda a\\0&1\end{pmatrix}
\]
implements both images of $u,v$. The entries assemble over all
primes to give a matrix over $\widehat{\mathbb Z}$. A matrix
commuting with both $U$ and $L$ is scalar, by its four entries.
Consequently the normalized upper-triangular implementer is unique,
and uniqueness shows $\mathcal V_{FG}=\mathcal V_F\mathcal V_G$.
It also shows continuity, either from the entry formula or compactness.
Passing to the central quotient gives the asserted covariance on
$\widehat B$ and hence invariance of each congruence kernel.

At level $n$, $\Gamma_0(n)$ is the stabilizer of a primitive line
in $(\mathbb Z/n)^2$. Conjugation by a matrix in $\GL_2(\mathbb Z/n)$
changes it to the stabilizer of another primitive line. Such a line
has a primitive generator which extends to an $\SL_2$ basis, by
Bezout over $\mathbb Z/n$, so $\SL_2(\mathbb Z/n)$ is transitive
on the primitive lines. Surjectivity of the congruence map then
proves preservation of the full preimage's conjugacy class.

If the projective implementer is scalar, its lower-right entry
forces $\lambda=1$ and $t=0$. For multiplier one the formula of
C.1 for $uvu$ gives, with $S=ULU$,
\[
 f(U^2,L^2)=S\mathcal V_F S^{-1}\mathcal V_F^{-1}
       =\begin{pmatrix}1&-t\\-t&1+t^2\end{pmatrix}.
\]
This matrix is the identity exactly when $t=0$, proving the kernel
criterion. All of these statements concern the canonical genuine
action, not an unspecified representative of an outer class.
\end{proof}

\begin{statement}{C.13}{Proved}{The two maximal abelian CM fields meet in one extra quadratic extension}
Inside a fixed algebraic closure of $\mathbb Q$,
\[
 \mathbb Q(i)^{\mathrm{ab}}\cap
 \mathbb Q(\sqrt{-3})^{\mathrm{ab}}
          =\mathbb Q^{\mathrm{ab}}(\sqrt[4]{3}).
\]
This establishes the field-theoretic assertion in Lead 25.
\end{statement}
\begin{proof}
Call the intersection $L$. Both fields on the left are Galois over
$\mathbb Q$, and contain $\mathbb Q^{\mathrm{ab}}$. The group
$G=\Gal(L/\mathbb Q)$ has two distinct abelian subgroups $A,B$ of
index two, fixing the two different quadratic fields. Their
intersection $C$ is central: it commutes with $A$ and with $B$,
which together generate $G$. Thus $G/C=C_2^2$. Choose representatives
$a,b$ of a basis. Every commutator is generated by the central
$[a,b]$, and $[a,b]^2=[a^2,b]=1$ since $a^2\in C$.
The closed derived subgroup therefore has order at most two.
Its fixed field is $\mathbb Q^{\mathrm{ab}}$: the latter is already
contained in $L$, and any larger abelian subextension would by
definition be contained in $\mathbb Q^{\mathrm{ab}}$.

Put $t=\sqrt[4]{3}$, choosing the positive real root. The splitting
field $\mathbb Q(i,t)$ has dihedral group of order eight. Over
$\mathbb Q(i)$ it is cyclic of order four. Over
$\mathbb Q(\sqrt{-3})$ its group is the Klein four-group: in the
presentation $r(t)=it$, $r(i)=i$, $s(t)=t$, $s(i)=-i$, the elements
fixing $it^2=\sqrt{-3}$ are $1,r^2,rs,r^3s$.
Thus $t$ belongs to both maximal abelian extensions. It does not
belong to $\mathbb Q^{\mathrm{ab}}$, since its Galois closure has
nonabelian dihedral group. Hence $L/\mathbb Q^{\mathrm{ab}}$ has
degree exactly two and is the stated extension.
\end{proof}

\begin{statement}{C.14}{Proved}{The complete necessary level-four cycle test has no descent defect}
Let $H$ be the characteristic quotient of $F$ obtained from all
surjections to $S_4$. Then $|H|=442368$, $|H'|=110592$, and
$\exp(H)=12$. In the multiplier modulus $24$, there are exactly
$96$ pairs satisfying the two cycle equations and inducing an
automorphism of $H$: twelve for each
$\lambda=1,5,7,11,13,17,19,23$. Every one satisfies the
level-four image of the equation in C.1 in at least one allowed
gauge. These are necessary finite solutions, with no assertion of
full profinite extendability.
\end{statement}
\begin{proof}
The independent executable \texttt{code/pbc\_level4.py} lists all
$216$ generating pairs of $S_4$, divides them into their nine
simultaneous conjugacy orbits, and enumerates the subgroup of
$S_4^9$ generated by the nine coordinate pairs. This is exactly
$H$: all automorphisms of $S_4$ are inner, so taking one representative
per orbit leaves the intersection of the kernels unchanged.
Breadth-first enumeration records every element as a word in the
two displayed generators and verifies both substitutions
$(x,y)\mapsto(y,x)$ and $(x,y)\mapsto(y,(xy)^{-1})$ on every
Cayley edge. It records the abelian quotient $C_2^2$ and its kernel,
then checks the cycle equations on every element of that kernel.
For the automorphism test, the transformed nine coordinate pairs
must represent nine distinct generating-pair orbits. This is
sufficient and necessary: it permutes the factors defining the
intersection of kernels. It is stronger than checking only that
each target pair generates $S_4$.

The program also constructs $\PSL_2(\mathbb Z/4)$ as matrices modulo
a single sign and identifies it with $S_4$ through its conjugation
on the four Sylow-three subgroups. It evaluates the two maps
$x\mapsto(rs)^2,y\mapsto(sr)^2$ and $x\mapsto s,y\mapsto sr$
simultaneously, verifies both on every Cayley edge, and tests
$\xi_a=y^{\nu+a}fx^{-a}$ for every $a\bmod12$. The output records
all residues, all passing gauges, and the defect for every tested
gauge. The saved output contains the asserted counts and no failures.
This finite exhaustive calculation proves precisely the stated
finite assertion. It does not solve the inverse-limit extension
problem recorded in C.6.
\end{proof}

\begin{statement}{C.15}{Proved}{The harmonic orbifold cover is not a modular level-three-to-level-four cover}
The map $w=4\lambda(1-\lambda)$ is an orbifold covering of degree
two from the sphere of signature $(2,3,3)$ to that of signature
$(2,3,4)$ with the cone markings described below. It does not
identify a covering between the principal modular curves of levels
three and four over the same modular $j$-line.
\end{statement}
\begin{proof}
Put order-three cones at $\lambda=0,1$ and an order-two cone at
$\lambda=\infty$. On the $w$-sphere put cones of orders $3,2,4$
at $0,1,\infty$. The points $0,1$ map unramified to $0$;
$\lambda=1/2$ has ramification index two over $w=1$ and is an
ordinary upstairs point; infinity has ramification index two,
turning the downstairs order-four cone into an upstairs order-two
cone. The local indices prove the orbifold assertion.

A cover of principal modular curves over the same modular orbifold
would require inclusion of the corresponding subgroups of
$\PSL_2(\mathbb Z)$, up to conjugacy. But $T^3\in\Gamma(3)$
is not in $\Gamma(4)$, and $T^4\in\Gamma(4)$ is not in
$\Gamma(3)$, for $T=\left(\begin{smallmatrix}1&1\\0&1\end{smallmatrix}\right)$.
The groups are normal, so conjugacy changes nothing. Equivalently
all their cusp widths are three and four, and neither is a multiple
of the other. The abstract inclusion of tetrahedral into octahedral
rotation groups therefore does not supply the asserted modular
comparison of the two based CM frames. This corrects that step
of Lead 24 without refuting its still-unproved parity implication.
\end{proof}

\begin{statement}{C.16}{Known}{Closed AA subgroups have zero parity}
Every closed subgroup of $\widehat{\mathrm{GT}}$ with the property
AA of Hoshi--Mochizuki--Tsujimura is contained in a conjugate of
the specified Galois image. It is consequently contained in
$\ker\Delta$. Moreover that Galois image is the unique maximal
closed AA subgroup containing the chosen complex conjugation $c$.
\end{statement}
\begin{proof}
Theorem 5.17(ii) of \cite{PBC_HMT2025}, together with the containment
argument immediately preceding its maximality conclusion, states
that every closed AA subgroup is contained in a
$\widehat{\mathrm{GT}}$-conjugate of $G_{\mathbb Q}$ and that the
maximal ones are exactly these conjugates. The property here is
their Definition 5.12, with the configuration-space arithmeticity
conditions; it is not the much weaker assertion that one element
preserves a single tripod inertia class. Since $\Delta$ is a
character of the full gentle group it is conjugacy-invariant, and
R24 gives $\Delta(G_{\mathbb Q})=0$. The first conclusion follows.

Suppose $gG_{\mathbb Q}g^{-1}$ contains $c$. Then
$g^{-1}cg$ is a nontrivial involution of $G_{\mathbb Q}$. Such
involutions are conjugate to complex conjugation: their fixed
fields are real closures, and embeddings of real closures in a
fixed algebraic closure are conjugate. Thus choose $h\in G_{\mathbb Q}$
with $g^{-1}cg=hch^{-1}$. It follows that $gh$ centralizes $c$.
By P.1 its centralizer in $\widehat{\mathrm{GT}}$ is
$\langle c\rangle\subset G_{\mathbb Q}$. Hence $g\in G_{\mathbb Q}$,
so the conjugate subgroup is the original one. This proves
uniqueness, and explains exactly the role of the later theorem
in the AA clause of Task C.8.
\end{proof}

\subsection{Exceptional prime-index covers and the explicit degree-five resolvents}

\begin{statement}{C.17}{Proved}{The corrected prime-index classification and all three signatures}
For primes $p\geq5$, the group $\PSL_2(\mathbb F_p)$ has a
subgroup of index $p$ precisely for $p=5,7,11$. The corresponding
subgroups of the modular group have genus zero, one cusp, and
\[
 (e_2,e_3)=(1,2),(3,1),(3,2),
\]
respectively. Their profinite completions are consequently
\[
 C_2\amalg C_3\amalg C_3,\qquad
 C_2^{\amalg3}\amalg C_3,\qquad
 C_2^{\amalg3}\amalg C_3^{\amalg2}.
\]
Retaining any one order-two factor and any one order-three
factor, and killing the other factors, maps the unique cusp
to the standard cusp of $\widehat B$.
The restriction $p\geq5$ is necessary, as noted in C.10.
\end{statement}
\begin{proof}
An index-$p$ subgroup has order $(p^2-1)/2$, prime to $p$.
The tame finite subgroups of $\PGL_2$ are cyclic, dihedral,
$A_4,S_4,A_5$; see \cite[Introduction and Proposition 1.1]{PBC_Beauville2010}.
This list also follows from the tame genus-zero
Riemann--Hurwitz formula: cyclic inertia and
$\sum(1-1/e_i)=2-2/|H|$ leave the cyclic case and the
triples $(2,2,n),(2,3,3),(2,3,4),(2,3,5)$.
A prime-to-$p$ element in $\PSL_2(\mathbb F_p)$ has order
dividing $(p-1)/2$ or $(p+1)/2$, by its two eigenvalues in
$\mathbb F_{p^2}$. Thus a cyclic or dihedral subgroup has
order at most $p+1$, smaller than $(p^2-1)/2$ for $p\geq5$.
The remaining orders $12,24,60$ force $p=5,7,11$.

Here are explicit subgroups establishing existence. In each
case use $S=\left[\begin{smallmatrix}0&-1\\1&0\end{smallmatrix}\right]$
and the indicated matrix $B_p$, with projective classes taken
modulo $p$:
\[
 B_5=\begin{pmatrix}4&2\\1&2\end{pmatrix},\quad
 B_7=\begin{pmatrix}2&4\\1&6\end{pmatrix},\quad
 B_{11}=\begin{pmatrix}3&4\\1&9\end{pmatrix}.
\]
Their determinants are one; $S,B_p,SB_p$ have orders
$(2,3,3),(2,3,4),(2,3,5)$, and the generated groups have
orders $12,24,60$. These assertions can be checked by the
finite closures in \texttt{pbc\_resolvents.py}, which computes
the groups directly from the matrices.
For $R=\left[\begin{smallmatrix}0&-1\\1&-1\end{smallmatrix}\right]$,
the exact coset-cycle data are
\[
\begin{array}{c|ccc}
 p&S&R&SR\\\hline
 5&1\,2^2&1^2\,3&5\\
 7&1^3\,2^2&1\,3^2&7\\
 11&1^3\,2^4&1^2\,3^3&11.
\end{array}
\]
All cosets and all permutation cycles are independently
enumerated in that program. The fixed points of $S,R$ give
$e_2,e_3$, and the sole $SR$-cycle gives the unique cusp.
The genus formula
$g=1+p/12-e_2/4-e_3/3-1/2$ gives zero in all three cases.
The resulting genus-zero orbifold presentation has these
finite cyclic factors and no free factor. Its cusp generator
is the inverse product of the finite-factor generators in
their boundary order. Killing all but one of each order
therefore sends it to the inverse product of the two retained
generators, up to conjugacy. This proves the cusp assertion
and the displayed profinite free-product descriptions.
\end{proof}

\begin{statement}{C.18}{Proved}{Both degree-five resolvents, their exact domain and their arithmetic field}
Put $K=\mathbb Q(\sqrt{-15})$ and
\[
 b_\pm=\frac{-5\pm3\sqrt{-15}}6,\qquad
 \mathcal D_5=\{F\in\widehat{\mathrm{GT}}_0:
                         (\tfrac{-15}{\lambda(F)})=1\}.
\]
There are two continuous multiplier-preserving homomorphisms
\[
              \mathfrak T^{(5)}_\pm:\mathcal D_5
                      \longrightarrow\widehat{\mathrm{GT}}_0.
\]
They are the two choices of retained order-three factor in
Lead 27(a). Both are identity on the specified image of $G_K$.
The associated marked orbifold filling has exact field of
definition $K$; the two fillings are interchanged by its
nontrivial automorphism over $\mathbb Q$.
No fixedness assertion for arbitrary pentagon elements is
included in this statement.
\end{statement}
\begin{proof}
The index-five modular cover is the explicit polynomial
\[
            j=27t^3(9t^2+15t+40),\qquad
 j-1728=27(t-1)(3t^2+4t+8)^2.
 \tag{C.18.1}
\]
Its derivative is $405t^2(3t^2+4t+8)$. Thus its only branch
values are $0,1728,\infty$, with passports
$(3,1,1),(2,2,1),(5)$ as required. Its transitive monodromy
on five sheets is generated by an order-three cycle and a
double transposition, with product a five-cycle. It is $A_5$:
its order is divisible by $2,3,5$, and the possible index-two
subgroup of $A_5$ is excluded because three-cycles generate
$A_5$. The standard modular quotient modulo five has the
same spherical presentation $(2,3,5)$, so this is the
$A_4$-stabilizer cover in C.17.

As an orbifold over the modular $(2,3,\infty)$ line, the
$t$-line has an order-two point at $1$, order-three points
at $b_+,b_-$, and its unique cusp at $\infty$. The other
points above $0$ and $1728$ have precisely the ramification
that removes their stabilizers. Filling the point $b_\mp$
and using the affine coordinate
\[
                  j'_\pm=1728\frac{t-b_\pm}{1-b_\pm}
 \tag{C.18.2}
\]
leaves the standard $(2,3,\infty)$ orbifold. Equations
(C.18.1)--(C.18.2) are verified symbolically in
\texttt{pbc\_resolvents.py}; in particular no picture of a
triangulation is needed to identify this filling map.

We verify the exact domain on the whole gentle group, not
only on Galois. By C.12 its action on
$\PSL_2(\mathbb F_5)=A_5$ is conjugation by a projective
matrix of determinant class $\lambda$. The $A_4$ subgroups
form one conjugacy class, so every gentle element preserves
the conjugacy class of the index-five subgroup $\Gamma_5$.
Its normalizer in the modular group is itself, since
$N_{A_5}(A_4)=A_4$. Hence restriction after an inner adjustment
defines an unambiguous outer action on $\widehat\Gamma_5$.

There are exactly two $A_4$ conjugacy classes of elements of
order three, each of size four. They distinguish its two
order-three orbifold points: the centralizer in $A_5$ of
such an element has order three, so no further double-coset
ambiguity remains. The projective normalizer of $A_4$ is
$S_4$, with quotient $C_2$ acting by exchange of these two
classes. Its quotient character is the determinant square
class, namely $(\lambda/5)$. Raising the generator to the
power $\lambda$ exchanges the classes exactly when
$\lambda\equiv-1\pmod3$. Thus the permutation of the two
geometric order-three points is
\[
                 (\lambda/5)(\lambda\bmod3)
                        =(\tfrac{-15}{\lambda}).
 \tag{C.18.3}
\]
Here $(\lambda\bmod3)$ is regarded as $\pm1$.
The finite program enumerates the projective normalizer of
order $24$, its square-determinant subgroup of order $12$,
both conjugacy classes, and all eight unit residues modulo
fifteen; it verifies (C.18.3) on every residue.
Consequently the domain preserving either chosen factor is
exactly $\mathcal D_5$. The unique order-two point and cusp
are preserved automatically.

For $F\in\mathcal D_5$, restrict its modular extension to
$\widehat\Gamma_5$, then quotient by the closed normal
subgroup generated by the other order-three factor. This
subgroup is invariant under the restriction and its inverse,
so the result is an automorphism of $\widehat B$.
It preserves the two torsion classes and the cusp, with
multiplier $\lambda(F)$. Restriction to the characteristic
$K_2$ gives the indicated gentle pair, by the extension and
restriction argument in C.1. Normalizer self-equality removes
the ambiguity in the first restriction, and all subsequent
inner ambiguities map to inner automorphisms. The operations
therefore compose as homomorphisms. Each is continuous in
the compact automorphism topology, proving the assertion
about continuous homomorphisms and multipliers.

Finally (C.18.1) is defined over $\mathbb Q$, whereas either
marked filling (C.18.2) is defined over $K$ and fixes its
three specified orbifold points there. Functoriality of the
arithmetic fundamental group for this covering, filling,
and coordinate isomorphism shows that the induced operator
on $G_K$ is the original arithmetic outer action. The final
$K_2$ restriction has the same property. Thus
$\mathfrak T^{(5)}_\pm(g)=g$ for every $g\in G_K$.
Since the two roots $b_\pm$ are exchanged by the nontrivial
element of $\operatorname{Gal}(K/\mathbb Q)$, neither chosen
filling is defined over a smaller field. There is no hidden
descent by a deck automorphism swapping them: the degree-five
cover has trivial deck group, since $N_{A_5}(A_4)=A_4$.
On the specified
Galois image its domain is exactly $G_K$ by (C.18.3).
The comparison with the pentagon, its values on the specified
E.58/D.46 constructions, and the common fixed subgroup are
separate tasks and remain open.
\end{proof}

\begin{statement}{C.19}{Proved}{Zero parity does not imply the pentagon}
The full element lifting $\varphi_0$ in B.27 can be chosen
with $\Delta=0$. Consequently
\[
       \ker\Delta\not\subset\widehat{\mathrm{GT}}.
\]
More generally the metabelian image of
$\widehat{\mathrm{GT}}_0\cap\ker\Delta$ contains the
entire cubic sign family $\prod_{p\geq5}C_2$ of R28.
This does not decide the opposite inclusion
$\widehat{\mathrm{GT}}\subset\ker\Delta$.
\end{statement}
\begin{proof}
Retain the precise central extensions and peripheral
frames constructed in B.27. Let
\[
 q:\Gamma\longrightarrow
 Q_{12}=\mathrm{SL}_2(\mathbb Z/12)/\{\pm I\}
\]
be the modular quotient used to read the R24 frames.
Its order is $576$, and its image of $F$ has order
$96$ and is isomorphic to
$C_2^2\times\mathrm{SL}_2(\mathbb F_3)$.
One way to see the latter description is Chinese
remaindering: the mod-four kernel above identity
modulo two is $C_2^3$, with $-I$ a direct factor;
quotienting its diagonal with $-I$ modulo three
leaves the displayed product. Thus
\[
             q(F')\simeq Q_8,\qquad
             q(F'')\simeq C_2\subset Z(Q_{12}).
 \tag{C.19.1}
\]
Direct matrix enumeration independently checks the
orders $576,96,8,2$ and the last centrality in
\texttt{pbc\_negative\_shadow.py}.

Put $N=F''$ as in B.27. By (C.19.1), $q$ factors
through $E=\Gamma/[N,\Gamma]$, and, since its
central kernel image is a two-group, through $E_2$.
The Cartesian description in B.27.4 is
\[
 E_2=\overline\Gamma\mathbin\times_{\Gamma_0}E_{0,2},
 \qquad V=M'_5.
\]
The kernel of $E_2\to E_{0,2}$ is $V$, a pro-five
group. Every homomorphism from it to $Q_{12}$ is
trivial because five does not divide $576$.
Consequently $q$ factors further through $E_{0,2}$.
The automorphism on this factor in B.27 is identity,
and its two retained frames are $(b_s,1)$ and
$(b_r,1)$. Their images in $Q_{12}$ are therefore
both identity; the third retained frame is already
identity. The source lifting argument preserves all
these specified finite frame coordinates. Its final
inner correction fixing $\tau$ lies in the kernel
of $E$, so it preserves them as well.

For the resulting full gentle element, the unique
pure frames $\beta,A$ of R24 consequently satisfy
$q(\beta)=q(A)=1$. The coordinates $T,W$ that define
$\delta$ at multiplier one are both identity.
Thus $\Delta=\delta=0$. Its metabelian image is
still $\varphi_0$, so B.2 and B.9 exclude the
pentagon. This proves the stated noncontainment.

For a cubic sign at any one coefficient prime
$p\geq5$, the same argument uses a pro-$p$ kernel
and the identical prime-to-$p$ quotient $Q_{12}$.
Choose each one-prime lift in $\ker\Delta$.
Their finite products realize all finite sign
tuples, still with zero parity. Compactness of
$\widehat{\mathrm{GT}}_0\cap\ker\Delta$ and
closedness of its metabelian image then give
every infinite tuple, as in B.27.
\end{proof}

\begin{statement}{C.20}{Proved}{The full adelic CM image of the multiplier-one gentle group}
Let $\mathcal H=\ker(\lambda:\widehat{\mathrm{GT}}_0
\to\widehat{\mathbb Z}^{\times})$, put
$\mathcal O_3=\mathbb Z[\omega]$, $\mathcal O_4=\mathbb Z[i]$,
and define compact groups
\[
 H_j=\{u\in\widehat{\mathcal O}_j^{\times}:
                u\equiv1\pmod2,\ u\bar u=1\}/\{\pm1\},
                    \qquad j=3,4.
 \tag{C.20.1}
\]
In each factor the quotient is by the single global sign;
it is not a separate quotient by signs at every prime.
The actual adelic frame map is a continuous surjective
homomorphism
\[
           (T,W):\mathcal H\longrightarrow H_3\times H_4.
 \tag{C.20.2}
\]
Its image on $\mathcal H\cap\ker\Delta$ is exactly
the image of $G_{\mathbb Q^{\rm ab}}$, the subgroup of
index two defined by equality of the two restrictions to
$\mathbb Q^{\rm ab}(\sqrt[4]{3})$.
This proves the adelic image assertion of Lead 25 with
its necessary modulo-two normalization retained.
\end{statement}
\begin{proof}
We first construct the adelic map and its full target.
Use the modular quotient
$q:\Gamma\to\mathrm{SL}_2(\widehat{\mathbb Z})/\{\pm I\}$
and the normalized extension $\Phi_F$ of a gentle element.
C.12 supplies the unique intertwiner
\[
 V_F=\begin{pmatrix}\lambda(F)&t_F\\0&1\end{pmatrix},
 \qquad q\Phi_F=\operatorname{Ad}(V_F)q.
\]
Uniqueness gives $V_{FG}=V_FV_G$. For $F\in\mathcal H$
let $A_F,\beta_F\in F$ be the unique pure torsion frames
of R24 and put
\[
       T_F=q(A_F)^{-1}V_F,\qquad
       W_F=q(\beta_F)^{-1}V_F.
 \tag{C.20.3}
\]
The torsion-frame identities make these commute
projectively with the cubic and quartic rotations.
For the cubic rotation, its nonzero trace excludes
anticommutation of genuine matrix lifts. For the
quartic rotation, R24's signed identity modulo four
excludes anticommutation: $W_F\equiv I\pmod4$,
whereas $S$ and $-S$ are distinct modulo four.
The exact centralizer rings are
$\widehat{\mathcal O}_3$ and
$\widehat{\mathcal O}_4$, and determinant is their
quadratic norm. The frames are pure and the quotient
action modulo two is identity, so both coordinates
are one modulo two. Their determinants are one.
Thus (C.20.3) belongs to (C.20.1).

Uniqueness of the pure frames gives
$A_{FG}=\Phi_F(A_G)A_F$ and
$\beta_{FG}=\Phi_F(\beta_G)\beta_F$.
Substitution of the covariance and $V_{FG}=V_FV_G$
proves $T_{FG}=T_FT_G$, $W_{FG}=W_FW_G$
at every modulus, hence adelically. The pure frame
maps are continuous because conjugation gives a
continuous bijection from a compact group onto its
torsion conjugacy class with the prescribed pure
coset. This proves continuity. This is the same
actual covariance calculation as R24 E.65.12 and
E.65.15--16, now with its entire adelic target stated.

We next justify the arithmetic interpretation needed
to determine the image. An elliptic curve modulo
$\{\pm1\}$ is a projective line with four branch
points, one distinguished by the identity. Ordering
the remaining three gives the Legendre parameter;
changing the order is the geometric $S_3$ action.
This identifies the rigidified modular stack with
the quotient of the tripod by $S_3$.
Here the arithmetic normalization at the tangential
cusp is essential. The actual Tate degeneration has
torsion represented by $\zeta_n^a q^{b/n}$ and the
compatible exact sequence
\[
 0\longrightarrow\widehat{\mathbb Z}(1)
 \longrightarrow T(E_q)
 \longrightarrow\widehat{\mathbb Z}\longrightarrow0.
\]
The quotient has trivial Galois action, so its genuine
matrix is upper triangular with lower-right entry one;
the determinant is $\chi$. The Tate power-series
identities are defined over $\mathbb Z$, hence apply
after substitution into $\mathbb Q((\lambda))$;
see \cite[III, \S3, equations (23)--(26), pp.~102--103]
{PBC_MilneEC2021}. Since
$j(\lambda)=256(1-\lambda+\lambda^2)^3/
(\lambda^2(1-\lambda)^2)$ and
$j(q)=q^{-1}+744+\cdots$, its parameter is
$q=\lambda^2/256\,(1+\lambda\mathbb Q[[\lambda]])$.
The formal unit and the chosen roots affect the shear,
but not the trivial action on the quotient
$\widehat{\mathbb Z}$ or its lower-right entry one.
The Legendre curve differs from this Tate curve by
one quadratic twist, whose action is the same sign
at every coefficient prime and disappears in the
rigidification. Thus the actual cusp representative
is precisely the unique matrix $V_g$ used above.
Covariance and determinant alone would have allowed
independent scalar signs at different primes; this
Tate normalization excludes those extra choices.
Transporting
its projective Tate local system to either elliptic
point by the same paths that define the pure frames
gives exactly $A^{-1}V$ and $\beta^{-1}V$.
The finite inertia ambiguity is removed by their
prescribed modulo-two action, since the elliptic
automorphisms modulo $\{\pm1\}$ act faithfully on
two-torsion. Consequently the two arithmetic
representations are, in the transported bases, those of
\[
       E_3:y^2=x^3+1,\qquad E_4:y^2=x^3-x,
 \tag{C.20.4}
\]
modulo the single global sign.
That is the global $\mu_2$ ambiguity inherited from
the elliptic moduli problem, not a choice of
primewise projective twists.

For an explicit check of the curve normalizations,
at Legendre parameter $1/2$ the substitution
$X=4x-2$, $Y=8y$ gives $Y^2=X^3-4X$, the
quadratic twist by two of $E_4$. At
$\lambda^2-\lambda+1=0$, translation gives
$y^2=x^3+(2\lambda-1)/9$; its constant is the
cube of $-(2\lambda-1)/3$, so it is a quadratic
twist of $E_3$ over $\mathbb Q(\sqrt{-3})$.
The descent normalizations are characterized by
their two-torsion action. A rational $j=0$ model
with this two-torsion field has cubic constant
a rational cube, hence is a quadratic twist of
$E_3$; a rational $j=1728$ model with rational
two-torsion has coefficient a negative rational
square, hence is a quadratic twist of $E_4$.
Quadratic twists disappear in the specified global
projective quotient. This checks the full arithmetic
comparison rather than just the mod-twelve bits.

Put $K_j=\mathbb Q\mathcal O_j$. Both quadratic
rings have class number one, for example by their
Euclidean norms. CM reciprocity identifies the
Weber-function torsion action with
$\widehat{\mathcal O}_j^{\times}/\mu_{w_j}$,
where $w_3=6$ and $w_4=4$; taking all torsion
levels gives $K_j^{\rm ab}$.
The exact finite-level input is
\cite[Corollaries 3.4 and 3.7]{PBC_LozanoRobledo2022};
the adelic action is also
\cite[Theorem 3.10]{PBC_MilneCM2007}.
The projective Tate field contains these Weber values
($x^3$ in the cubic case, $x^2$ in the quartic case).
Conversely its Tate action over $K_j$ is linear over
the commutative ring $\widehat{\mathcal O}_j$, so
its torsion field is abelian over $K_j$. Thus its
projective field is exactly $K_j^{\rm ab}$.

To determine the image as an actual subgroup of the
specified matrices, observe that reduction gives
isomorphisms
\[
       \mu_{w_j}/\{\pm1\}
                 \xrightarrow{\sim}(\mathcal O_j/2)^\times.
 \tag{C.20.5}
\]
For $j=3$ these groups have order three; for $j=4$
they have order two. Hence the natural map
\[
 \{u\in\widehat{\mathcal O}_j^\times:
               u\equiv1\pmod2\}/\{\pm1\}
       \longrightarrow
               \widehat{\mathcal O}_j^\times/\mu_{w_j}
 \tag{C.20.6}
\]
is an isomorphism. The curves (C.20.4) have all their
two-torsion over $K_j$, so their projective images
lie on the left of (C.20.6). CM reciprocity makes
their projection onto the right surjective.
They therefore fill the whole left side.
The Weil pairing identifies determinant with the
cyclotomic character; its kernel is exactly $H_j$.
It follows that the image of $G_{\mathbb Q^{\rm ab}}$
in each coordinate is exactly $H_j$, not merely an
open subgroup of it.

By C.13 the two fields have intersection
$K_3^{\rm ab}\cap K_4^{\rm ab}
=\mathbb Q^{\rm ab}(\sqrt[4]{3})$.
Galois theory therefore makes the joint arithmetic
image the fibre product over its order-two Galois
group over $\mathbb Q^{\rm ab}$. Denote this
index-two subgroup of $H_3\times H_4$ by $H_{\rm ar}$.
R24 identifies its defining character with
$\epsilon_3(T)+\epsilon_2(W)$, exactly $\Delta$
on $\mathcal H$, retaining the global sign. Its
element E.58 lies outside $H_{\rm ar}$.

The image of (C.20.2) is a subgroup containing
$H_{\rm ar}$ and an element of its other coset.
It is therefore the entire product. Since the
character defining $H_{\rm ar}$ pulls back to
$\Delta$, the image of $\mathcal H\cap\ker\Delta$
is contained in $H_{\rm ar}$ and contains its
arithmetic realization. It is exactly $H_{\rm ar}$,
as claimed. No separate lifting of an arbitrary
adelic frame pair is needed for this last argument.
\end{proof}

\begin{statement}{C.21}{Proved}{The CM image dichotomy and uniqueness of its nonarithmetic character}
The positive CM image of either
$\widehat{\mathrm{GT}}$ or $C^\varphi$ has exactly
two possibilities: $H_{\rm ar}$ or $H_3\times H_4$.
It is the first precisely when $\Delta$ vanishes
on that group, and the second otherwise.
Every continuous finite-order character of
$\widehat{\mathrm{GT}}_0$ that factors through the
adelic CM coordinates and vanishes on the Galois
image is either zero or $\Delta$.
This includes every such character obtained by
reducing these same coordinates to another finite
congruence modulus. It does not classify characters
that use additional source information.
\end{statement}
\begin{proof}
Both groups in the first assertion contain the
specified Galois image. Their positive CM images
therefore contain $H_{\rm ar}$. By C.20 they are
subgroups of its index-two overgroup
$H_3\times H_4$, so there is no intermediate
possibility. The quotient character is precisely
$\Delta$. To test vanishing on the whole group,
an element of any multiplier can be divided by
a Galois element of that multiplier, which stays
in the group and leaves its $\Delta$ value unchanged.
Thus the positive test is sufficient.

A continuous character of the positive CM product
that vanishes on the arithmetic image factors
through $(H_3\times H_4)/H_{\rm ar}=C_2$.
Its only possible nontrivial value is the unique
quadratic character. For the entire gentle group,
the same division by an equal-multiplier Galois
element identifies this character with the unique
extension $\Delta$ in R24. Reduction to a finite
modulus is a quotient of the adelic coordinate map,
so cannot introduce another such character. This
does not assert that an arbitrary continuous
character of the full gentle group factors through
these two coordinates.
\end{proof}

\begin{statement}{C.22}{Refuted}{Separate arithmeticity of the CM coordinates does not force zero parity}
The assertion that separate arithmetic realizability of $T$
and $W$ forces $\Delta=0$ is false, even for actual full
gentle elements. In particular the following three claims
of Lead 27(b) cannot all hold: commutation with the
Eisenstein family is equivalent to separate arithmeticity
of $T$; commutation with the Gaussian family is equivalent
to separate arithmeticity of $W$; and commutation with
both families forces zero parity.
This statement does not decide either commutation
equivalence in isolation.
\end{statement}
\begin{proof}
Take the full multiplier-one element $e$ of R24 E.58.
It has $\Delta(e)=1$. By C.20 the individual arithmetic
maps onto $H_3$ and $H_4$ are surjective. Thus there
exist $g_3,g_4\in G_{\mathbb Q^{\rm ab}}$ with
\[
                  T(g_3)=T(e),\qquad W(g_4)=W(e).
\]
Both coordinates are separately arithmetic at every
level and adelically. They cannot be realized by one
common Galois element, since the joint arithmetic
image is $H_{\rm ar}$ and $\Delta(e)=1$.
This is the promised counterexample to the separate
arithmeticity assertion. If the two stated individual
equivalences both held, this same $e$ would commute
with both families; the third claim would then give
$\Delta(e)=0$, a contradiction. A valid replacement
must retain either a stronger geometric compatibility
or the requirement of a single common arithmetic
realization. The two independent coordinate tests
alone lose exactly this information.
\end{proof}

\begin{statement}{C.23}{Proved}{An actual full gentle element has different parity from its descent}
There is an actual multiplier-one element $e_U$ of
$\widehat{\mathrm{GT}}_0$ such that
\[
             \Delta(e_U)=1,\qquad \Delta(D(e_U))=0.
 \tag{C.23.1}
\]
Consequently $C^\varphi$ is a proper closed subgroup, and
\[
 \eta:=\Delta+\Delta\circ D
 \quad\hbox{is a nonzero continuous character with}\quad
 G_{\mathbb Q},\ C^\varphi\subset\ker\eta.
 \tag{C.23.2}
\]
In particular $\Delta\circ D\ne\Delta$.
This does not decide whether $\Delta\circ D$ vanishes
identically or whether $\Delta$ vanishes on $C^\varphi$.
\end{statement}
\begin{proof}
We retain the finite evaluation of the input and the
finite evaluations of its descent in the same actual
source. Write $\Gamma=C_2\amalg C_3$, $\tau=sr$,
$F=\ker(\Gamma\to S_3)$, and use temporarily the
R24 marking
\[
       x=\tau^{-2},\qquad y=sxs,\qquad z=(xy)^{-1}.
\]
Let $q:\Gamma\to Q_{12}=\mathrm{SL}_2(\mathbb Z/12)/\{\pm I\}$
be the marked modular quotient, and write
\[
 U_0=\begin{pmatrix}1&1\\0&1\end{pmatrix},\quad
 S=\begin{pmatrix}0&1\\-1&0\end{pmatrix},\quad
 R=\begin{pmatrix}0&-1\\1&-1\end{pmatrix}.
\]
Here $q(s)=[S]$, $q(r)=[R]$, $SR=U_0^{-1}$,
and the canonical signed restriction of $q$ to $F$
has generators $U_0^2,SU_0^2S^{-1}$.
The descent map $\varphi$ of C.1, in this marking, is
\[
            x\longmapsto\tau^{-1},\quad
            y\longmapsto s,\quad z\longmapsto r.
 \tag{C.23.3}
\]
Thus its signed matrix lift on the free group has
triple $(U_0,S,R)$, whose product is $I$.

Conjugation by $r$ on $F$ is the exact cycle
$\rho:(x,y,z)\mapsto(y,z,x)$, and conjugation by
$s$ interchanges $x,y$. The six outer classes of
$\Gamma/F$ are therefore represented by
$\rho^i$ and $\rho^i\theta$, where
$\theta(x,y)=(y,x)$ and $0\leq i<3$.
For each representative let $\kappa_\sigma:F\to
\mathrm{SL}_2(\mathbb Z/12)$ be the signed map of
(C.23.3) precomposed with $\sigma$. These are the
six ordered pairs of distinct entries of $(U_0,S,R)$.
Set
\[
 H=F'F^3,\qquad
 U=H\cap\bigcap_{\sigma\in S_3}\ker\kappa_\sigma.
 \tag{C.23.4}
\]
This is an explicitly specified open normal subgroup
of $\Gamma$. To check normality, composition of two
chosen representatives differs from another by an
inner automorphism of $F$; such a change does not
change the kernel of a homomorphism. The six kernels
are consequently permuted by conjugation in $\Gamma$.

The essential point is that
\[
                       [5I]\in q(U).
 \tag{C.23.5}
\]
Here is a short word certificate. With
$[a,b]=aba^{-1}b^{-1}$, put
\[
 w=[x^4,(xy^2)^4]^2\,[x^4,y^4]^2.
 \tag{C.23.6}
\]
It belongs to $F'\subset H$. In
$\mathrm{SL}_2(\mathbb Z/4)$ the expression
$[a^4,b^4]^2$ is identically $I$: modulo two the
fourth powers lie in the cyclic subgroup $A_3$ of
$\mathrm{SL}_2(\mathbb F_2)=S_3$, so their
commutator lies in the principal kernel, which has
exponent two. Thus every evaluation of $w$ is $I$
modulo four.
Modulo three, direct multiplication gives the following
complete table. Each sign denotes the scalar matrix
with that sign. In the first column the pair is
$(U_0^2,SU_0^2S^{-1})$; in the remaining columns it
is the indicated ordered pair.
\[
\begin{array}{c|c|rrrrrr}
 &q|_F&(U_0,S)&(U_0,R)&(S,U_0)&(S,R)&(R,U_0)&(R,S)\\ \hline
 {[x^4,(xy^2)^4]^2}&+&+&-&+&+&-&+\\
 {[x^4,y^4]^2}&-&+&-&+&+&-&+\\ \hline
 w&-&+&+&+&+&+&+
\end{array}
\]
Chinese remaindering proves $q(w)=[5I]$ and
$\kappa_\sigma(w)=I$ for all six $\sigma$,
which proves (C.23.5). The independent executed
certificate \texttt{pbc\_descent\_ghost\_word.py}
checks this table and all $48^2=2304$ instances
of the modulo-four identity. Its JSON retains the
actual matrices and both constituent blocks.

We now apply the relative centralizer construction of
R24 E.58 to this smaller $U$, keeping every old and
new finite coordinate. Since $U\subset H$,
\[
               q(U)\subset q(H)=C_2^2\times Q_8
\]
is a two-group. Hence $q$ factors through
$\Gamma_2=\Gamma/O^2(U)$. Put
$D_U=\Gamma/U$ and
$J=\operatorname{im}(\Gamma\to Q_{12}\times D_U)$.
The element $([5I],1)$ lies in $J$ by (C.23.5)
and centralizes the image of $r$.
The kernel of $\Gamma_2\to J$ is pro-two.
Coprime complement conjugacy, with compactness,
therefore lifts this specified element to
\[
 c\in C_{\Gamma_2}(r)\cap U^{(2)},\qquad
                    q(c)=[5I].
 \tag{C.23.7}
\]
Indeed a preliminary lift conjugates the two
order-three complements over $\langle r\rangle$;
a correction in that pro-two kernel makes it
centralize the same generator $r$.

The exact relative product
\[
 \Gamma/R_{\rm nil}(U)
       =\prod_{p,D_U}\Gamma/O^p(U)
\]
now carries the identity automorphism with the
specified peripheral conjugators
\[
 b_s=b_\tau=1,\qquad
 b_r=(c\text{ at }2,\ 1\text{ at every other prime}).
 \tag{C.23.8}
\]
All coordinates have the same value one in $D_U$.
The full relative lifting theorem retains these
actual conjugators. For completeness, the central
refinement argument of P.3 applies to this arbitrary
open normal $U$ without change: in a finite joint
source take the nilpotent residual of the image of
$U$; its quotient is nilpotent, and every auxiliary
central extension has nilpotent image of $U$ as
well. It therefore factors through
$\Gamma/R_{\rm nil}(U)$. The same count and
compactness give the full peripheral automorphism
with the specified frames. This is precisely the
relative form proved in R24 E.58, rather than an
assumption that a finite residue alone lifts.

Adjusting by the inverse of the retained
$\tau$-conjugator makes the full lift $\Phi$
fix $\tau$ and changes none of the retained data.
Its unique pure frames $A,\beta$ satisfy
\[
       q(A)=[5I],\qquad q(\beta)=1,
       \qquad f=s\Phi(s)\in R_{\rm nil}(U).
 \tag{C.23.9}
\]
The two torsion relations give the full two-cycle
and three-cycle equations, and restriction of
$\Phi$ to $F$ is an automorphism. Thus this is
an actual full gentle element $e_U$. Its cusp
intertwiner is $I$, so its normalized CM frames
are $([5I],1)$ and $\Delta(e_U)=1$.

It remains to compute the descent, including its
frames. Return to the question's marking
$x_{\rm old}=y^{-1}$, $y_{\rm old}=x^{-1}$.
The normalized representative is obtained from the
one just constructed by conjugating by $f$;
its normalized conjugator is $f_{\rm old}=f^{-1}$,
read in the old letters. In particular it belongs
to $U$. The same abstract map $\varphi$ is being
used throughout; in the old marking its cycle is
induced by $r^{-1}$, whereas the temporary cycle
above is induced by $r$. Changes of an inner
representative do not affect invariance of $U$.
For any positive gentle pair in that marking
put
\[
                  h=\rho(f_{\rm old})f_{\rm old}.
\]
The cycle equations give
$\alpha(xy)=h^{-1}(xy)h$. The subgroup $U$ is
invariant under this geometric outer $S_3$ action,
so $h\in U$ as well. By the identity member of
(C.23.4),
\[
       q\varphi(f_{\rm old})=q\varphi(h)=1.
 \tag{C.23.10}
\]
Write $v=\varphi(f_{\rm old})$ and
$u=\varphi(h)$. Before cusp normalization, the
descended automorphism fixes $s$, sends
$\tau$ to $v^{-1}\tau v$, and sends
$r$ to $u^{-1}ru$. Conjugating by $v$ fixes
$\tau$ and gives torsion frames
\[
                         \beta_D=v,
                 \qquad A_D=vu^{-1}.
\]
They are pure, since their images in $Q_{12}$
are identity and reduction modulo two is the
specified $S_3$ quotient. They are consequently
the unique pure frames. The descended congruence
action is identity, so its uniquely normalized
upper triangular intertwiner is identity too.
Any further cusp gauge multiplies the frames and
intertwiner by the same matrix and leaves $T,W$
unchanged. Therefore $T(D(e_U))=W(D(e_U))=1$
at the parity level, proving (C.23.1).

Finally C.1 makes $D$ a continuous homomorphism
fixing the specified Galois image. Thus (C.23.2)
is a continuous character, vanishes on Galois,
and vanishes at every fixed point of $D$.
It has value one on $e_U$, which proves all the
claimed strictness and nonvanishing statements.
\end{proof}

\begin{statement}{C.24}{Proved}{The descended parity has an explicit finite detector that does not detect the original parity}
On the multiplier-one gentle group, $\Delta\circ D$
is determined by the residue of the normalized $f$
in the explicitly marked finite quotient
\[
 E_\varphi=F\Big/\bigcap_{\sigma\in S_3}\ker\kappa_\sigma
 \tag{C.24.1}
\]
of C.23. The original character $\Delta$ is not
determined by that same residue. Moreover
$\Delta\circ D$ factors through the two adelic CM
coordinates of C.20 if and only if it is identically
zero. Thus either $C^\varphi\subset\ker\Delta$,
or the space of continuous quadratic characters
of $\widehat{\mathrm{GT}}_0$ trivial on Galois
has dimension at least two. The local argument below
initially leaves this alternative open; C.27 subsequently
proves $\Delta D\ne0$ and the two-character conclusion.
\end{statement}
\begin{proof}
Use the old marking and let $f$ be its normalized
conjugator. Put $v=\varphi(f)$ and
$u=\varphi(\rho(f)f)$. Both finite values in
$Q_{12}$ are determined by (C.24.1); the old
cycle is the inverse of the temporary cycle in
C.23, up to the harmless inner representative.
After conjugation by $v$, the descended modular
automorphism fixes $\tau$ and has torsion frames
$v$ and $vu^{-1}$ before making the latter pure.
Modulo $S_3$, $u$ lies in $A_3=\langle r\rangle$,
so it centralizes $r$ there. Before conjugation
by $v$, the descended map therefore fixes $r$
and $s$ modulo $S_3$, hence fixes $\tau$;
since $v\in A_3$, its conjugation
of the transposition $\tau$ is trivial only
when $v=1$ modulo $S_3$. Write $u=r^k$ modulo
$S_3$, with $k\in\{0,1,2\}$. The actual pure
frames are consequently
\[
                    \beta_D=v,\qquad
                    A_D=vu^{-1}r^k.
 \tag{C.24.2}
\]
Right multiplication by $r^k$ does not alter
the conjugation of $r$. Uniqueness follows from
$C_\Gamma(s)\cap K_2=C_\Gamma(r)\cap K_2=1$.
In particular their finite values retain the
signs of the actual pure frames.

C.12 supplies the normalized upper triangular
intertwiner $V_D$ for this descended action.
At multiplier one its reduction is $U_0^t$;
it is determined already by the conjugation
of $S$. Indeed two such reductions would have
quotient $U_0^a$ centralizing $[S]$, and direct
multiplication gives $a=0$ modulo twelve.
Thus (C.24.2) determines both signed finite
coordinates $A_D^{-1}V_D$ and $\beta_D^{-1}V_D$.
The canonical representative of a pure element
is the one with modulo-four reduction in the
four-element group of C.19. These representatives
therefore determine the two bits of R24 E.65
without a sign choice. This proves the asserted
finite factorization of $\Delta\circ D$.

The actual element $e_U$ of C.23 and the identity
have the same $f$-residue in (C.24.1), namely one,
and have opposite $\Delta$ values. This proves
the failure of factorization for the original
character. In particular a detector containing
only these six evaluations of $f$ cannot be a
detector of $\Delta$, even though it detects
$\Delta\circ D$.

If $\Delta\circ D$ factors through the CM
coordinates, C.21 makes it either zero or
$\Delta$: it is a quadratic character trivial
on Galois. C.23 excludes the latter value.
The converse for the zero character is immediate.
Finally if $\Delta\circ D=0$, every fixed point
of $D$ has $\Delta=0$. If it is nonzero, it
is different from the nonzero character $\Delta$
by C.23, so these two characters are linearly
independent over $\mathbb F_2$. This proves
the stated alternative, including its exact
unresolved scope.
\end{proof}

\begin{statement}{C.25}{Proved}{An exact comparison with the harmonic primitive equation}
Use the braid marking
\[
 x=\tau_1^2,\quad y=\tau_2^2,\quad
 \eta=\tau_1\tau_2\tau_1,\quad
 \tau=sr=\tau_1^{-1}
\]
after passing from $\widehat B_3$ to
$\Gamma=\widehat B_3/Z(\widehat B_3)$; thus
$\eta=s$ and $r=\tau_1\tau_2$ in $\Gamma$.
Let $g=\beta^{-1}$ be the unique harmonic primitive
of a full gentle pair, characterized by
\[
 f=\theta(g)^{-1}g,\qquad
 \theta(x,y)=(y,x),\qquad
                  g^{\rm ab}=(r_2,r_2).
 \tag{C.25.1}
\]
For every multiplier, membership in $C^\varphi$
is equivalent to the existence of $r_3\in\widehat{\mathbb Z}$
such that the following literal equation holds
in $\widehat B_3$:
\[
 g(\tau_1^2,\tau_2^2)
 =\eta^{2r_2-r_3}f(\tau_1,\eta)
                   \tau_1^{-2r_2+3r_3}.
 \tag{C.25.2}
\]
This is the equation denoted (GF$_0$) by
Nakamura--Tsunogai. The parameter $r_3$ is unique
when it exists. On $C^\varphi$ it is a continuous
cyclotomic one-cocycle extending the Kummer
cocycle of $3$ on the specified Galois image.
Here $r_2$ is the frame cocycle, extending the
Kummer cocycle of $2$.

With these same parameters, the usual harmonic
equation (I$'$) is equivalent to (GF$_1$), and
on $C^\varphi$ either one is equivalent to the
additional equation
\[
 f(\tau_1,\eta)
 =\eta^{r_3-2r_2}f(\tau_1^2,\eta)
                   \tau_1^{6r_2-3r_3}.
 \tag{C.25.3}
\]
This is Nakamura--Tsunogai's equation (V).
These assertions do not identify the frame
cocycles with Ihara's separately defined
extensions on every pentagon element, do not
remove (C.25.3), and do not imply the
Nakamura--Schneps relation (III$'$).
\end{statement}
\begin{proof}
The primary formula and its geometric interpretation
are opened in \cite[Theorem 2, (GF$_0$), and
\S4]{PBC_NT2002}. We now prove the full-group
comparison, without assuming that its Galois
proof extends to every gentle element.
The change from the question's marking is exactly
$x_{\rm new}=y_{\rm old}^{-1}$,
$y_{\rm new}=x_{\rm old}^{-1}$, as in C.23.
In particular the same abstract descent map is
\[
              \varphi(x)=\tau_1,\qquad
              \varphi(y)=\eta.
 \tag{C.25.4}
\]
Conjugation by $\eta$ interchanges $x,y$.
The unique pure conjugator $\beta$ of $s$ in
the canonical modular extension satisfies
$f=s\Phi(s)=\theta(\beta)\beta^{-1}$,
so its inverse is precisely the primitive $g$
in (C.25.1). Abelianization of that equation
makes the two coordinates of $g$ equal.

Write $v=\varphi(f)=f(\tau_1,\eta)$.
The canonical extension $\Phi$ of the input
pair and the directly descended automorphism
$d$ have
\[
 \Phi(\tau)=d(\tau)=\tau^\lambda,\qquad
 \Phi(s)=\beta s\beta^{-1},\qquad
                  d(s)=v^{-1}sv.
 \tag{C.25.5}
\]
C.1 shows that $D(F)=F$ is exactly equality
of these outer extensions. Since the centralizer
of $\tau^\lambda$ is $\langle\tau\rangle$,
this is equivalent to
\[
                 \Phi=\operatorname{Inn}(\tau^t)d
                       \quad(t\in\widehat{\mathbb Z}).
 \tag{C.25.6}
\]
Comparing the actual conjugators of $s$ and
using $C_\Gamma(s)=\langle s\rangle$ then gives
\[
               g=s^a v\tau^{-t},\qquad a\in\mathbb Z/2.
 \tag{C.25.7}
\]
Conversely (C.25.7) makes (C.25.5) agree after
that inner adjustment on both generators,
so proves (C.25.6).

There is no extra numerical freedom hidden in
(C.25.7). Identify $\Gamma^{\rm ab}$ with
$\mathbb Z/6$ by $[\tau_1]=1$; then
$[s]=3$, $[\tau]=5$, and
$[g]=4r_2$. Since $f$ is a commutator word,
$v\in\Gamma'$. Thus (C.25.7) implies
\[
                  t+3a\equiv4r_2\pmod6.
 \tag{C.25.8}
\]
It follows first that $t+2r_2$ belongs to
$3\widehat{\mathbb Z}$, and therefore
\[
                r_3=(t+2r_2)/3
 \tag{C.25.9}
\]
is well-defined. Modulo two, (C.25.8) gives
$a\equiv t\equiv r_3\equiv2r_2-r_3$.
Substitution into (C.25.7) is exactly
(C.25.2) in $\Gamma$. Conversely, that equation
immediately supplies (C.25.7) with
$t=-2r_2+3r_3$.

To pass from $\Gamma$ back to the literal
braid group, the possible discrepancy is a
power of its central generator
$(\tau_1\tau_2)^3=\eta^2$.
The braid abelianization sends each $\tau_i$
to one, sends that central generator to six,
and sends either side of (C.25.2) to $4r_2$:
\[
             3(2r_2-r_3)+(-2r_2+3r_3)=4r_2.
\]
Consequently the discrepancy is zero. The
restriction of this central quotient to
$\langle\tau_1^2,\tau_2^2\rangle$ is the
usual isomorphism onto $K_2$, so the left-hand
side is the same actual primitive throughout.
This proves both directions at every multiplier.

Uniqueness of $t$ in (C.25.6) follows because
the centre of $\Gamma$ is trivial: two values
would give an inner automorphism centralizing
all of $\Gamma$, and $\langle\tau\rangle$
is torsion-free. The graph of $t$ over the
compact group $C^\varphi$ is closed, proving
continuity. For $F,G\in C^\varphi$,
composition of (C.25.6) gives
\[
       t(FG)=t(F)+\lambda(F)t(G),
\]
since $d_F(\tau)=\tau^{\lambda(F)}$.
The frame multiplication law similarly gives
$r_2(FG)=r_2(F)+\lambda(F)r_2(G)$.
Formula (C.25.9) proves the claimed cocycle
law for $r_3$. On Galois the frame cocycle
$r_2$ equals the usual Kummer cocycle by
\cite[\S5, before Proposition 5.1]{PBC_NT2002};
its established (GF$_0$) formula and uniqueness
identify our $r_3$ with the Kummer cocycle of
$3$ there.

For the remaining explicit substitution, put
\[
       g_1=f(\tau_1^2,\eta)\tau_1^{4r_2}.
\]
It lies in $\langle\tau_1^2,\tau_2^2\rangle$:
in the subgroup generated by $\tau_1^2$ and
$\eta$, quotient by that free subgroup is
procyclic, and a commutator word maps to one.
Conjugating by $\eta$ interchanges $\tau_1$
and $\tau_2$ and fixes $\eta$. Therefore
\[
 \theta(g_1)^{-1}g_1
 =\tau_2^{-4r_2}f(\tau_2^2,\eta)^{-1}
                      f(\tau_1^2,\eta)\tau_1^{4r_2}.
 \tag{C.25.10}
\]
The harmonic equation (I$'$) says that this
is $f$. Uniqueness of the pure primitive
makes it equivalent to $g_1=g$, namely
(GF$_1$). Equating (GF$_1$) and (C.25.2),
then cancelling the displayed powers, gives
(C.25.3) exactly. This also proves the
converse implication on $C^\varphi$.

Finally \cite[Proposition 5.1 and the preceding
paragraph]{PBC_NT2002} explicitly distinguishes
its frame extension $\widetilde\rho_2$ from
Ihara's extension and derives their equality
under a further harmonic condition. We retain
that distinction, and do not replace the
newly obtained $r_3$ by another extension
without a proof. The equianharmonic and
four-strand comparisons still require their
additional equations.
\end{proof}

\begin{statement}{C.26}{Proved}{The same fixed-point equation is the equianharmonic primitive equation}
Keep C.25's marking and put
\[
 \nu=(\lambda-1)/2,\quad
 \varepsilon\in\{1,-1\},\quad\lambda\equiv\varepsilon\pmod3,
 \quad \nu_\varepsilon=(\lambda-\varepsilon)/2,
 \quad\xi_+=\tau_1\tau_2,\quad\xi_-=\tau_2\tau_1.
\]
Let $h$ be the unique equianharmonic primitive with
\[
 f=
 \begin{cases}
 y^{-\nu}\rho(h)^{-1}h,&\varepsilon=1,\\
 y^{-\nu}\rho(h)^{-1}y^{-1}h,&\varepsilon=-1.
 \end{cases}
 \tag{C.26.1}
\]
Then $h^{\rm ab}=(-\nu_\varepsilon/3,\nu_\varepsilon/3)$.
Membership in $C^\varphi$ is equivalent to the existence
of $r_3\in\widehat{\mathbb Z}$, with
$\lambda-\varepsilon-6r_3\in4\widehat{\mathbb Z}$,
satisfying the literal braid equation
\[
 h(\tau_1^2,\tau_2^2)
 =\xi_\varepsilon^{r_2+(\lambda-\varepsilon-6r_3)/4}
       f(\tau_1,\xi_\varepsilon)
       \tau_1^{3r_3-2r_2-\nu_\varepsilon}.
 \tag{C.26.2}
\]
This is (HF$_0$). The parameter is precisely the
unique cocycle of C.25. In particular its parity on
$C^\varphi$ is forced by
\[
                      r_3\equiv\nu_\varepsilon\pmod2.
 \tag{C.26.3}
\]
On $C^\varphi$, with these same parameters, the
usual equianharmonic equation (II$'$) is equivalent
to (HF$_1$), and either is equivalent to
\[
 f(\tau_1,\tau_1\tau_2)
 = (\tau_1\tau_2)^{-r_2}
         f(\tau_1^2,\tau_1\tau_2)\tau_1^{2r_2}.
 \tag{C.26.4}
\]
This is (IV$'_{\rm bis}$). Thus the exact extra
conditions for (I$'$) and (II$'$) on the descent-fixed
subgroup are (V) and (IV$'_{\rm bis}$), respectively.
No implication from fixedness alone to either extra
condition, or to (III$'$), is asserted.
\end{statement}
\begin{proof}
Write $a=\tau_1$, $b=\tau_2$, $r=ab$, and $s=aba$.
All intermediate conjugator calculations will be in
$\Gamma$; we lift them to the braid group explicitly.
The primitive (C.26.1) exists for every full gentle
pair. Indeed the canonical modular extension has
\[
                    \Phi(r)=f^{-1}r x^{-\nu}.
\]
Its torsion class is that of $r^\varepsilon$.
For $\varepsilon=1$, its unique pure conjugator
can be written $h^{-1}$. For $\varepsilon=-1$,
the corresponding conjugator can uniquely be
written $h^{-1}s$ with $h\in F$.
This uses the quotient $S_3$ to remove the
remaining $\langle r\rangle$ ambiguity. Since
$yr=sr^{-1}s$, in the two cases respectively
\[
 \Phi(r)=h^{-1}rh,\qquad
 \Phi(r)=h^{-1}sr^{-1}sh=h^{-1}yrh.
 \tag{C.26.5}
\]
Comparing with $f^{-1}rx^{-\nu}$ gives exactly
(C.26.1). Its abelianization uses
$\rho(A,B)=(-B,A-B)$ and gives
$(I-\rho)h^{\rm ab}=(0,\nu_\varepsilon)$,
which proves the stated abelianization and its
zero sum.

Put $u=f(a,r)$. The peripheral hexagon gives
\[
 \alpha(z)=x^\nu f(z,x)z^\lambda f(z,x)^{-1}x^{-\nu}.
 \tag{C.26.6}
\]
One direct check is that
$\alpha\rho=\operatorname{Inn}(f^{-1}y^{-\nu})\rho\alpha$,
while the three-cycle gives
$f^{-1}y^{-\nu}\rho(f)^{-1}=x^\nu f(z,x)z^\nu$.
Applying the first identity to $y$ and cancelling
the last power of $z$ proves (C.26.6).
Under the actual descent (C.25.4), $z$ maps to
$r$, so
\[
                   d(r)=a^\nu u^{-1}r^\varepsilon u a^{-\nu}.
 \tag{C.26.7}
\]
Here replacing a braid lift of $r$ by a central
multiple does not change evaluation of the
commutator word $f$.

Suppose first that $F\in C^\varphi$, and use
$t=-2r_2+3r_3$ from C.25, so
$\Phi=\operatorname{Inn}(a^{-t})d$.
If $\varepsilon=1$, comparison of (C.26.5)
and (C.26.7) gives
\[
                      h=r^k u a^{t-\nu}
                         \quad(k\in\mathbb Z/3).
 \tag{C.26.8}
\]
If $\varepsilon=-1$, the same comparison gives
$h=s r^k u a^{t-\nu}$. In $\Gamma$ one has
$as=r^{-1}$ and $\xi_-=a^{-1}ra$, hence this
last identity becomes
\[
              h=\xi_-^{k-1}f(a,\xi_-)
                                  a^{t-\nu-1}.
 \tag{C.26.9}
\]
In either case put $B=t-\nu_\varepsilon$.
Since $h$ and $u$ have zero image in
$\Gamma^{\rm ab}$, (C.26.8) or (C.26.9)
gives $2A+B=0$ modulo six for the indicated
power $A$ of $\xi_\varepsilon$.
Consequently $B$ is even and
$A\equiv-B/2$ modulo three.
But
\[
 -B/2=r_2+(\lambda-\varepsilon-6r_3)/4,
\]
which proves integrality, (C.26.3), and the
projective form of (C.26.2).
The braid abelianization of $h$ is zero,
and the right-hand side has abelianization
$2(-B/2)+B=0$. Its possible central discrepancy
therefore vanishes, exactly as in C.25.
This proves the literal braid identity.

Conversely assume (C.26.2), with the stated
integral exponents, and set $t=-2r_2+3r_3$.
For $\varepsilon=1$ it makes the conjugation
of $r$ in (C.26.5) agree with
$\operatorname{Inn}(a^{-t})d$.
For $\varepsilon=-1$, rewrite it using
$\xi_-=a^{-1}ra$ and $as=r^{-1}$;
it then makes the conjugation of $r^{-1}$
by $h^{-1}s$ agree with that same map.
Both automorphisms already send $\tau$ to
$\tau^\lambda$. Agreement on $r$ and $\tau$
is agreement on $\Gamma$, proving fixedness.
Uniqueness in C.25 identifies the parameter.

For precision about the remaining substitution,
write (HF$_1$) as $h_1=h$, where
\[
 B_1=3r_3-\nu_\varepsilon,\qquad
 h_1=\xi_\varepsilon^{-B_1/2}
                      f(a^2,\xi_\varepsilon)a^{B_1}.
 \tag{C.26.10}
\]
The exponent $B_1$ is even by (C.26.3), and
is divisible by three because
$\nu_\varepsilon$ is. Thus $h_1$ projects
into $F$ and has braid abelianization zero.
Substituting its projection into (C.26.1),
and using $rar^{-1}=b$, gives in both cases
\[
 f(x,y)=b^{-3r_3-\nu}f(b^2,r)^{-1}
                           r^\nu f(a^2,r)a^{3r_3-\nu}.
 \tag{C.26.11}
\]
For $\varepsilon=1$ use $r^\nu=1$ in
$\Gamma$. For $\varepsilon=-1$ use
$\xi_-=a^{-1}ra$ and $r^\nu=r^{-1}$;
the middle product is then $b^{-1}a^{-1}=r^{-1}$.
These are exactly the two required calculations,
not an identification of the two branches.
Equation (C.26.11) is the equianharmonic
relation (II$'$) of the opened primary note
\cite[Theorem 1]{PBC_NT2002}, also printed with
its corrected braid normalization in
\cite[\S5]{ROOT_NTY2010}.
Uniqueness of the primitive in (C.26.1) gives
both directions in $\Gamma$. Both sides of
(C.26.11) have braid abelianization zero,
as do $h$ and $h_1$, so the equivalence holds
in $\widehat B_3$ as written.

Finally comparison of (C.26.2) with
(C.26.10) cancels the outer powers and gives
\[
 f(a,\xi_\varepsilon)
       =\xi_\varepsilon^{-r_2}
                         f(a^2,\xi_\varepsilon)a^{2r_2}.
\]
For $\varepsilon=-1$ this is just the
conjugate by $a^{-1}$ of the equation with
$\xi_+=r$. It is therefore equivalent to
(C.26.4) in both branches. This proves the
claimed exact comparison and preserves the
remaining extra conditions.
\end{proof}

\begin{statement}{C.27}{Proved}{Descended parity is nonzero and gives a second independent character}
There is an actual multiplier-one full gentle element $g$ with
\[
                         \Delta(D(g))=1.
 \tag{C.27.1}
\]
Consequently
\[
 (\Delta,\Delta\circ D):
       \widehat{\mathrm{GT}}_{0,1}\twoheadrightarrow\mathbb F_2^2
 \tag{C.27.2}
\]
is onto. These are two linearly independent continuous quadratic
characters of the full gentle group, both trivial on the specified
Galois image. In particular $\Delta$ does not generate all such
characters, and the abelianization of the quotient by the closed
normal subgroup generated by Galois has a quotient $C_2^2$.
Moreover
\[
             \ker\Delta\not\subset C^\varphi,\qquad
                    C^\varphi\ne\ker\Delta.
 \tag{C.27.3}
\]
The inclusion $C^\varphi\subset\ker\Delta$ is not decided here.
The CM image of $D(\widehat{\mathrm{GT}}_{0,1})$ is the entire
product $H_3\times H_4$ of C.20.
\end{statement}
\begin{proof}
We specify a finite source, compute its actual centralization,
solve the peripheral obstruction there, and then use the full
central lifting criterion. This is more than survival of the
finite cycle equations.

Use the old marking
$x=(rs)^2$, $y=(sr)^2$ in $F=\ker(\Gamma\to S_3)$,
and put $\tau=sr$. For the six geometric outer permutations
of the peripheral triple, retain simultaneously the evaluations
\[
 \begin{split}
 \kappa_{4,\sigma}:F&\longrightarrow
                    \mathrm{SL}_2(\mathbb Z/4),\\
 \overline\kappa_{3,\sigma}:F&\longrightarrow
                    \mathrm{PSL}_2(\mathbb F_3)=A_4
 \end{split}
 \tag{C.27.4}
\]
of $\varphi\circ\sigma$. The signed first evaluation has
$x\mapsto S$, $y\mapsto U_0^{-1}$, with the matrices of C.23.
Let $N$ be the intersection of these twelve kernels and set
$G=\Gamma/N$. The same kernel-permutation argument as C.23
shows that $N$ is normal in the actual modular source.

Here is a short word for the required residue. Write
$c=[x,y]$ and use the commutator convention of C.23. Set
\[
 w=[c,ycy^{-1}]\,
       [c,(yx)c(yx)^{-1}]\,
       [c,(x^{-1}y^2)c(x^{-1}y^2)^{-1}].
 \tag{C.27.5}
\]
It belongs to $F''$ and has freely reduced length $58$.
Put $A=-U_0^2$ modulo four. In the order
$(S,\tau),(S,z),(\tau,S),(\tau,z),(z,S),(z,\tau)$
of the six ordered pairs, its signed modulo-four values are
\[
                           (A,I,A,A,I,A).
 \tag{C.27.6}
\]
All its projective modulo-three values are identity, since
$A_4$ is metabelian. For the old cycle
$\rho:(x,y,z)\mapsto(y,z,x)$, the first evaluation of
$\rho(w)w$ is identity modulo four.

The finite automorphism to be lifted is
\[
              \alpha(s)=s w^{-1},\qquad
              \alpha(r)=w r,\qquad \alpha(\tau)=\tau.
 \tag{C.27.7}
\]
The executed certificates establish the following exact finite
data:
\[
\begin{array}{c|l}
 |G|&6\,115\,295\,232\\
 |F/N|&1\,019\,215\,872\\
 (F/N)^{\rm ab}&C_{12}^2\\
 \operatorname{ord}(\bar x),\operatorname{ord}(\bar y)&12,\ 12\\
 G^{\rm ab}&C_6\\
 \alpha&\text{an automorphism, with pure }s\text{- and }r\text{-frames}.
\end{array}
 \tag{C.27.8}
\]
For reproducibility, \texttt{pbc\_modular\_envelope.py} constructs
the faithful induced degree-$288$ representation associated with
the signed modulo-four map. It uses the exact Schreier return
words for $x=(rs)^2$, $y=(sr)^2$, and checks the relations
$s^2=r^3=1$. The added projective modulo-three representation
has degree $72$. Their disjoint union is the faithful
degree-$360$ representation in \texttt{pbc\_modular\_joint.py}.
Thus the group in (C.27.8) is specified by actual permutations
of the two modular generators, not by an unproved presentation.
The word, every matrix in (C.27.6), and both finite peripheral
conjugacy assertions are checked by the programs.

We next calculate the exact source centralization
\[
                  E=\Gamma/[N,\Gamma].
 \tag{C.27.9}
\]
The details identifying the computed group with (C.27.9) matter.
The Polycyclic algorithm constructs the Schur extension $H$
of a finite polycyclic presentation of $G$; its defining
property is $H=F_m/[R_m,F_m]$ for that presentation
\cite[\S7.9-1 and \S7.9-2]{PBC_Polycyclic2026}.
Choose lifts $s_0,r_0$ of the two specified generators and let
$L=\langle s_0,r_0\rangle$. Since $H=L\ker(H\to G)$ and
the latter kernel is central, $L'=H'$. In particular $L$
contains the complete Schur multiplier in the central kernel.
The computation gives $L^{\rm ab}=\mathbb Z^2$.
The map from the free two-generator centralization to $L$
is therefore an isomorphism: it is an isomorphism on the
Hopf multiplier and on the free abelianization. Consequently
\[
                   E=L/\langle s_0^2,r_0^3\rangle
 \tag{C.27.10}
\]
is precisely the centralization for the source
$C_2\amalg C_3$, rather than an arbitrary Schur cover of $G$.
The two elements being killed are central.

The executed \texttt{pbc\_hopf\_joint.g} gives
\[
 \begin{gathered}
 |E|=293\,534\,171\,136,\qquad
 K=\ker(E\to G)\simeq C_2^4\times C_3,\\
 H_2(G,\widehat{\mathbb Z})\simeq K.
 \end{gathered}
 \tag{C.27.11}
\]
The second identification also follows from the actual Hopf
sequence: $H_2(\Gamma,\widehat{\mathbb Z})=0$ as in B.27,
and (C.27.8) makes $\Gamma^{\rm ab}\to G^{\rm ab}$ an
isomorphism. The saved certificate
\texttt{pbc\_hopf\_joint\_certificate.json} gives all $34$
relative generator orders, every power and conjugation
relation of this finite polycyclic group, and the coordinates
of $s,r$ and the conjugators used below. Together with the
two permutation generators, this retains the complete finite
object and its marked source maps.

Choose pure $s$- and $r$-frames for (C.27.7) in $F/N$,
and lift them to $E$. Conjugating $s,r$ by those lifts
initially gives a product $\tau\delta$, where $\delta\in K$
has order two. Changes of the two frames by their respective
finite centralizers alter $\delta$ by the images of the
central commutator maps
\[
 C_{F/N}(s)\longrightarrow K,\qquad
 C_{F/N}(r)\longrightarrow K.
 \tag{C.27.12}
\]
Their combined image $J$ has order $24$, and the computation
finds $\delta\in J$. More concretely, in the recorded
centralizer generating lists, multiplying the initial
$s$-frame by the lift of its seventh generator and leaving
the $r$-frame unchanged gives
\[
                 s_1^2=r_1^3=1,\qquad s_1r_1=\tau
 \tag{C.27.13}
\]
exactly in $E$. The output records the correction word
$[7]$, the ten $s$-centralizer generators, both corrected
frames, and both corrected images. It checks the product
equality and that $s\mapsto s_1$, $r\mapsto r_1$ is an
automorphism of $E$.
Independently of the last automorphism check, the actual
source map factors through (C.27.9), and Hopf naturality
identifies its action on $K$ with $H_2(\alpha)$.
It is therefore an isomorphism on kernel and quotient.

The modular form of B.11 now gives a full positive
peripheral automorphism $\Phi$ of $\Gamma$ covering
(C.27.7), with these finite frames retained.
Its cusp conjugator is identity in $G$; an inner correction
in $N$ therefore makes $\Phi(\tau)=\tau$ literally without
changing any retained value. The source torsion relations
give the full gentle cycle relations, not just their finite
images. Restriction to $F$ preserves the three positive
peripheral classes and centralizes the geometric $S_3$
in the outer group.

We check the remaining normalization explicitly. In the
temporary marking $x_{\rm new}=\tau^{-2}$,
$y_{\rm new}=s x_{\rm new}s$, put
$f_{\rm new}=s\Phi(s)$. Its two-cycle implies
$f_{\rm new}^{\rm ab}=(a,-a)$.
The finite value is $w^{-1}\in(F/N)'$, so (C.27.8)
gives $a\in12\widehat{\mathbb Z}$.
Conjugating $\Phi$ by $\tau^{-2a}=x_{\rm new}^a$
replaces this conjugator by
$y_{\rm new}^a f_{\rm new}x_{\rm new}^{-a}\in F'$.
Both $x_{\rm new}$ and $y_{\rm new}$ have order $12$
in $G$, so this correction is identity in $G$.
Returning to the old marking, as in C.23, the normalized
conjugator $f$ has precisely the finite residue $w$.
This proves existence of an actual full gentle pair
$g=(1,f)$ retaining (C.27.6) and every projective
modulo-three identity.

It remains to read its descended character without losing
the central signs not retained by the projective map.
There is a small complete signed check at three.
The six signed $\mathrm{SL}_2(\mathbb F_3)$ evaluations
have source image of order $18432$ and derived image of
order $512$. The program
\texttt{pbc\_descent\_cycle\_locus.py} exhausts all $512$
derived residues and checks both cycle equations.
Among the cycle residues whose six projective values
are identity, the two signed possibilities are
\[
 (I,I,I,I,I,I),\qquad
 (-I,-I,-I,I,-I,I).
 \tag{C.27.14}
\]
This is a necessary finite consequence of the full cycle
equations and of being an actual word; no lifting inference
is made from this enumeration.
In particular, for
$v=\varphi(f)$ and $u=\varphi(\rho(f)f)$, one has
\[
 v\bmod4=-U_0^2,\quad u\bmod4=I,\qquad
                 v\bmod3=u\bmod3=\epsilon I
                 \quad(\epsilon=\pm1).
 \tag{C.27.15}
\]

The actual pure descended frames of C.24 are
$\beta_D=v$ and $A_D=vu^{-1}$, since $u=1$ modulo two.
Their modulo-four diagonals are both $-1$, so their
canonical signed representatives are $-v$ and $-vu^{-1}$.
The uniquely normalized unipotent intertwiner is
$V_D=U_0^6$ modulo twelve: its exponent is two modulo
four and zero modulo three. Hence the two normalized
CM coordinates have reductions
\[
 \begin{array}{c|cc}
       &\bmod4&\bmod3\\\hline
 T_D=A_D^{-1}V_D&I&-I\\
 W_D=\beta_D^{-1}V_D&I&-\epsilon I .
 \end{array}
 \tag{C.27.16}
\]
The equianharmonic bit is one, whereas the Gaussian bit
is zero for either scalar $\pm I$. Thus
$\Delta(D(g))=1$, proving (C.27.1).

C.23 provides $e_U$ with character pair $(1,0)$, while
$g$ has pair $(\Delta(g),1)$. These two vectors generate
$\mathbb F_2^2$, proving (C.27.2) and the stated quotient
of the abelianization. Both characters vanish on Galois
because $D$ fixes that image. Multiplying $g$, if necessary,
by $e_U$ gives an element with $\Delta=0$ and
$\Delta D=1$. It cannot be fixed by $D$, which proves
(C.27.3). Finally $D(\widehat{\mathrm{GT}}_{0,1})$
contains the multiplier-one Galois image and now contains
an element in its other CM coset. The index-two image
calculation of C.20 proves the last assertion.
\end{proof}

\begin{statement}{C.28}{Proved}{Original parity remains free after two descended detectors are fixed}
There is an actual multiplier-one gentle element $e_2$ with
\[
 \bigl(\Delta(e_2),\Delta(D(e_2)),\Delta(D^2(e_2))\bigr)
                         =(1,0,0).
 \tag{C.28.1}
\]
In particular $\Delta$ does not belong to the span of
$\Delta D$ and $\Delta D^2$. This statement by itself
does not prove that the last two characters are independent.
\end{statement}
\begin{proof}
Let $q_1:\Gamma\twoheadrightarrow G_1$ be the exact
degree-$360$ modular quotient $G$ of C.27, and continue
to use the old marking. Let $N_1$ be the intersection
of the six signed level-twelve kernels in C.24, and put
\[
 N_2=\bigcap_{\sigma\in S_3}
                 \ker(q_1\varphi\sigma),\qquad
 U_2=F'F^3\cap N_1\cap N_2.
 \tag{C.28.2}
\]
The chosen six substitutions represent the geometric
$S_3$ up to inner automorphisms of $F$. Therefore each
intersection is normal in $\Gamma$, and $U_2$ is an
explicitly marked open normal subgroup.

We first verify the finite fact
\[
                         [5I]\in q(U_2),
 \tag{C.28.3}
\]
where $q:\Gamma\to Q_{12}$ is the original modular
quotient. The executed program
\texttt{pbc\_iterated\_\allowbreak ghost.py} constructs the following
three marked maps on the same two generators:
the six evaluations in $G_1$, the six signed evaluations
modulo three and modulo four, and the original signed
pure congruence map. The first map is represented on
$2160$ points and the second on $432$ points; the original
map uses $24+48$ points. All permutation arrays are saved
literally in \texttt{pbc\_iterated\_\allowbreak ghost.g}.
The original pure marking is
\[
       y\longmapsto U_0^{-2},\qquad
       x\longmapsto S U_0^{-2}S^{-1}.
 \tag{C.28.4}
\]
The resulting joint word image has order
$2^{148}3^{29}$. Its projection to the first two
detectors has order $2^{145}3^{29}$ and kernel of
order eight. Direct membership in its derived subgroup
places the permutation
\[
 \left(1,1;\ -I\pmod3,\ I\pmod4\right)
 \tag{C.28.5}
\]
in that derived subgroup. These exact orders and this
membership are retained in
\texttt{pbc\_iterated\_\allowbreak ghost\_\allowbreak certificate.json} and
in the executed GAP output. Membership in the derived
image means that (C.28.5) is the image of an actual
word in $F'$, not merely an independently specified
tuple in the product. Chinese remaindering gives
the central value $[5I]$ and proves (C.28.3).

The full construction of C.23 now applies verbatim
with $U_2$ in place of $U$. Indeed $q(U_2)$ is a
two-group because $U_2\subset F'F^3$.
Coprime centralizer lifting supplies a prescribed
$r$-centralizer element in the relative pro-two
component with old $q$-value $[5I]$ and identity
in $\Gamma/U_2$. On the relative product use the
identity automorphism with this $r$-frame and
identity $s$- and cusp frames. The full relative
lifting theorem of P.3 and C.23 then gives an
actual gentle element $e_2$ with
\[
 f\in R_{\rm nil}(U_2),\qquad
 q(A)=[5I],\quad q(\beta)=1,
 \qquad \Delta(e_2)=1.
 \tag{C.28.6}
\]
All the frames are retained; this is not an
existence inference from finite membership alone.

Write $v=\varphi(f)$ and
$u=\varphi(\rho(f)f)$. Normality of $U_2$ and
the two sets of kernels imply
\[
       q(v)=q(u)=1,
       \qquad q_1(v)=q_1(u)=1.
 \tag{C.28.7}
\]
Thus the actual pure frames of the first descent
are $v,vu^{-1}$; both have identity value in
$Q_{12}$ and in $G_1$. Their normalized cusp
intertwiner is identity at level twelve.
This proves $\Delta(D(e_2))=0$.

For the second descent one must also retain the
normalization of its input conjugator. The first
descended modular extension has both torsion
frames equal to one in $G_1$. Its unnormalized
pure conjugator consequently has value one there.
Its abelianization is of the form $(a,-a)$.
Since the pure abelianization of $G_1$ is
$C_{12}^2$, C.27 gives $a\in12\widehat{\mathbb Z}$.
Both pure inertia generators have order twelve
in $G_1$, so the correcting multiplication
$f_D\mapsto y^a f_Dx^{-a}$ is invisible in
$G_1$. The normalized conjugator of $D(e_2)$
therefore has identity value in $G_1$ as well.

Its six signed modulo-four evaluations are
identity, and its six projective modulo-three
evaluations are identity. The complete signed
cycle calculation (C.27.14) then gives, at the
second descent, $v_3=u_3=\epsilon I$ and
$v_4=u_4=I$. The pure $r$-frame has value $I$,
the pure $s$-frame has values $\epsilon I,I$
at three and four, and the normalized intertwiner
is $I$. Hence the equianharmonic bit is zero.
The Gaussian bit is zero for either scalar
$\epsilon I$; those two scalars are the squares
in its order-four norm-one group modulo three.
This proves $\Delta(D^2(e_2))=0$ and completes
(C.28.1). Evaluation on $e_2$ proves the last
claim of the statement.
\end{proof}

\begin{statement}{C.29}{Proved}{The degree-seven and degree-eleven resolvents, their fields and exact stabilizer domains}
For either exceptional index-seven modular cover, there
are three restriction--filling homomorphisms to
$\widehat{\mathrm{GT}}_0$, one for each retained
order-two factor. Their marked fields are
\[
 K_7=\mathbb Q(\eta),\quad \eta^2=-7,
 \qquad K_7\bigl(\sqrt{14-6\eta}\bigr),
 \tag{C.29.1}
\]
the first for one factor and the second for the other two
conjugate factors. For either exceptional index-eleven
cover, all six choices of a retained order-two and
order-three factor define such homomorphisms. Their
marked fields, each of degree twelve over $\mathbb Q$,
are
\[
 L_{a,b}=\mathbb Q(\omega,a,b),\qquad
 \omega^2-\omega+3=0,\quad Q_1(a)=P_1(b)=0,
 \tag{C.29.2}
\]
with the explicit polynomials below. Every homomorphism
is continuous, preserves the multiplier, and is identity
on the specified Galois group of its marked field.
Its intersection with the Galois image has exactly that
field as domain. Formula (C.29.8) below gives the exact
domain on the full gentle group using finite retained
frame data. These conclusions do not assert fixedness
of arbitrary pentagon elements.
\end{statement}
\begin{proof}
We first give the covers and fillings, then type and
verify their full-group domains. Before these calculations
we searched for the degree-seven and degree-eleven
polynomial dessins and opened the primary account
\cite[orbits 7.2 and 11.1, printed pp.~5--6 and
9--10]{PBC_AV2023}. Those are the classical Klein models.
We independently verify the polynomial identities and
correct two displayed formula errors below.

At seven a convenient normalized model is
\[
 \begin{gathered}
 b_7=\frac{11+\eta}{56},\qquad
 P(t)=t(t^2+t+b_7)^3,\qquad Q(t)=7t^2+4t+b_7,\\
 R=-\frac{351}{3294172}-\frac{27\eta}{470596},
 \qquad j(t)=1728P(t)/R.
 \end{gathered}
 \tag{C.29.3}
\]
The exact factorizations are
\[
 \begin{split}
 P'(t)&=(t^2+t+b_7)^2Q(t),\\
 P(t)-R&=\frac{Q(t)^2}{49}(t-a_0)
       \left(t^2+\frac{15+\eta}{14}t+
                         \frac{27(3+\eta)}{392}\right),\\
 a_0&=\frac{-11+\eta}{14}.
 \end{split}
 \tag{C.29.4}
\]
Thus the residual order-three point is $0$; the three
residual order-two points are $a_0$ and the two roots
$a_+,a_-$ of the quadratic in (C.29.4). Retaining
$a\in\{a_0,a_+,a_-\}$ gives the explicit filling
\[
                            j'_a=1728t/a.
 \tag{C.29.5}
\]
It sends the retained order-three point, order-two point
and unique cusp to $0,1728,\infty$ respectively.
The quadratic discriminant is $(14-6\eta)/49$.
Its numerator has norm $448$, which is not a rational
square; hence it is not a square in $K_7$.
This proves the fields in (C.29.1), once the absence
of hidden deck identifications is checked below.

For comparison with the opened primary model, put
$\omega=(1-\eta)/2$ and $z=-7\omega(2t+1)$.
Its polynomial factors are
$P_3=z^2+28$, $P_1=z+7\omega$,
$Q_2=z^2+6\omega z+4$ and
\[
 Q_1=\bigl(z^2-4z-4(1-\omega)(1-6\omega)\bigr)
                         \bigl(z-\omega(3+2\omega)\bigr).
\]
Exact subtraction gives
$P_3^3P_1-Q_2^2Q_1=-2^6 3^3\omega^{14}$.
The constant printed for orbit 7.2 has the opposite
sign. The change of variables above gives precisely
(C.29.3); in particular the displayed branch value
is $1728$ with the corrected sign.

At eleven put $K_{11}=\mathbb Q(\omega)$ with
$\omega^2-\omega+3=0$, and define
\[
 \begin{aligned}
 P_3&=(z-1-\omega)
            \bigl(z^2+(2+\omega)z-1+\omega\bigr),\\
 P_1&=z^2-3z-2\omega^2,\\
 Q_2&=(z-1+\omega)
       \bigl(z^3-(1+\omega)z^2-(1+2\omega)z-3(1-\omega)\bigr),\\
 Q_1&=z^3+4z^2+(6-5\omega)z+10-12\omega .
 \end{aligned}
 \tag{C.29.6}
\]
The exact identities are
\[
       j=-P_3^3P_1,\qquad
       j-1728=-Q_2^2Q_1,\qquad
       (P_3^3P_1)'=11P_3^2Q_2.
 \tag{C.29.7}
\]
All indicated factors are square-free, and the factors
whose roots must differ have nonzero resultants.
Thus the residual order-two points are the three
roots $a$ of $Q_1$, and the residual order-three
points are the two roots $b$ of $P_1$. For each
pair the filling is $j'_{a,b}=1728(z-b)/(a-b)$.
The second fraction displayed for orbit 11.1 in the
primary account omits its additive $1$; (C.29.7)
is the checked identity used here.

The discriminant of $P_1$ is $8\omega-15$, of
norm $297$, so $P_1$ is irreducible over $K_{11}$.
At the split prime five, with $\omega=2$, $Q_1$
reduces to $z^3+4z^2+z+1$; its values at
$0,1,2,3,4$ are $1,2,2,2,3$. Hence this cubic
is irreducible over the residue field and $Q_1$
is irreducible over $K_{11}$.
A cubic irreducible over a field remains irreducible
over a quadratic extension: otherwise it has a root
there, contradicting the degrees. Therefore
$[L_{a,b}:K_{11}]=6$ and $[L_{a,b}:\mathbb Q]=12$.
All symbolic assertions in this proof are executed
over the exact quadratic quotient rings in
\texttt{pbc\_resolvents\_\allowbreak remaining.py}.

We identify the finite covers precisely. The degree-seven
passport has exactly two permutation-pair classes:
fix an order-three permutation of type $3^2 1$,
enumerate the double transpositions whose product is
a seven-cycle, and quotient by its centralizer.
Each class generates a group of order $168$.
At eleven we use the primary orbit 11.1 identification
of the displayed model with monodromy
$\PSL_2(\mathbb F_{11})$; this is the source input,
not a conclusion from its passport alone.
In either case an epimorphism of the modular group
onto $\PSL_2(\mathbb F_p)$ whose cusp image has
order $p$ has the standard congruence kernel.
For this last assertion the finite program fixes
$S$, enumerates all order-three $R'$ with $SR'$
of order $p$, and verifies that they form precisely
one orbit under $C_{\PGL_2(\mathbb F_p)}(S)$.
It also checks the exceptional stabilizers and their
projective normalizers. Thus the models are the
index-$p$ congruence covers of C.17, with their two
possible conjugacy classes of stabilizer.

Here are exact finite domains, including the data
needed when several factors have the same finite
conjugacy class. Put $G_p=\PSL_2(\mathbb F_p)$,
let $H_p\simeq S_4,A_5$ be the chosen stabilizer,
and use left cosets $H_p\backslash G_p$ with
right multiplication. The retained points correspond
to fixed cosets $H_pg_a$ of $S$ and $H_pg_b$ of
$R$. An identification of the displayed cover with
this coset cover fixes these labels once and for all;
changing a lift or a base point changes only the
inner representatives described next.

For a full gentle element take a cusp-normalized
modular representative $\Phi$, its projective
intertwiner $V_p$ from C.12, and actual torsion
conjugators $B_2,B_3\in\Gamma$ such that
\[
 \begin{gathered}
 \Phi(s)=B_2sB_2^{-1},\qquad
 \Phi(r)=B_3r^\epsilon B_3^{-1},\\
                    \epsilon\in\{1,-1\},\quad
                    \epsilon\equiv\lambda\pmod3.
 \end{gathered}
\]
At a general multiplier we do not require $B_3$
to be pure. Its ambiguity is right multiplication
by a power of $r$, which is harmless on a fixed
coset. Write $q_p:\Gamma\to G_p$ for the
standard quotient. The exact domain is
\[
 \begin{gathered}
 (\lambda/p)=1,\\
 H_pg_aV_p^{-1}q_p(B_2)=H_pg_a,\qquad
 H_pg_bV_p^{-1}q_p(B_3)=H_pg_b .
 \end{gathered}
 \tag{C.29.8}
\]
In this formula the first condition puts $V_p$
in $G_p$, so both coset tests are defined.
The program supplies the explicit matrix representatives
and stabilizer matrices for all the fixed cosets.
At seven the order-three test is automatic, one
order-two point is fixed throughout this domain,
and the other two have the same additional
index-two stabilizer condition. At eleven the
order-two and order-three fixed sets have sizes
three and two, with stabilizers of orders four
and six in the respective normalizers of orders
twelve and twelve. Thus (C.29.8) retains both
necessary conditions and does not replace them
by a multiplier condition.

To prove the formula, the finite computation gives
$N_{\PGL_2(\mathbb F_p)}(H_p)=H_p$ for $p=7,11$.
The two $G_p$-classes of the stabilizer are exchanged
by the nonsquare determinant class. Hence a given
cover is preserved exactly when $(\lambda/p)=1$.
In that case $V_p\in G_p$; after an inner adjustment
of $\Phi$ its induced action on $G_p$ is identity.
Its action on a torsion fixed coset is then exactly
\[
       H_pg\longmapsto H_pgV_p^{-1}q_p(B_i).
\]
The order-two multiplier is automatically one modulo
two; the order-three multiplier $\epsilon$ merely
replaces a generator by its inverse and leaves its
fixed coset set unchanged. The stated right-power
ambiguity of $B_i$ leaves this formula unchanged.
This proves necessity and sufficiency of (C.29.8)
for preserving the selected geometric factor classes.

On this closed open stabilizer subgroup, restrict
to the index-$p$ modular subgroup and kill the
closed normal subgroup generated by all unretained
torsion factors. C.17 identifies the quotient with
$C_2\amalg C_3$ and sends the cusp to its standard
cusp. The subgroup being killed is preserved by
the action and its inverse. We obtain an actual
continuous automorphism of that quotient with
the original multiplier. Restricting to its
characteristic kernel $\Gamma(2)$ of its quotient
to $S_3$, a free subgroup of index six, gives the
claimed gentle element. The self-normalizing
property of $H_p$ removes the ambiguity in the
initial restriction, while all remaining changes
are inner. Composition and continuity therefore
give continuous homomorphisms on the domains
(C.29.8), exactly as in C.18.

Finally the displayed covers are defined over
$K_p=\mathbb Q(\sqrt{-p})$. Their arithmetic
stabilizer is exactly $G_{K_p}$ by the determinant
test just proved. Their deck group is trivial,
since $N_{G_p}(H_p)=H_p$. Consequently a selected
filling descends precisely when its retained
points do: there is no deck transformation that
can hide a Galois permutation of them. This proves
the exact fields (C.29.1)--(C.29.2) and the exact
intersection of each domain with Galois.
All covering, filling and coordinate maps are
defined over that field. Functoriality of the
arithmetic fundamental group then makes the
resolvent identity on its specified Galois image.
No implication for arbitrary pentagon pairs or
for the common fixed subgroup has been used.
\end{proof}

\begin{statement}{C.30}{Proved}{Countably infinite quadratic rank from beta defects}
Let $\mathcal S$ be the set of rational primes $p\equiv5\pmod {12}$.
For each $p\in\mathcal S$ there is an explicitly defined continuous
character
\[
 \vartheta_p:\widehat{\mathrm{GT}}_0\longrightarrow\mathbb F_2
 \tag{C.30.1}
\]
which vanishes on the full pentagon group $\widehat{\mathrm{GT}}$.
The joint character map is surjective even on
$\widehat{\mathrm{GT}}_{0,1}\cap\ker\Delta$:
\[
 \bigl(\vartheta_p\bigr)_{p\in\mathcal S}:
 \widehat{\mathrm{GT}}_{0,1}\cap\ker\Delta
 \twoheadrightarrow\prod_{p\in\mathcal S}\mathbb F_2.
 \tag{C.30.2}
\]
In particular, $\{\Delta\}\cup\{\vartheta_p:p\in\mathcal S\}$
is linearly independent.

The space of all continuous quadratic characters of
$\widehat{\mathrm{GT}}_0$ which vanish on the specified Galois image
has dimension exactly $\aleph_0$ over $\mathbb F_2$. The same is true
if ``Galois image'' is replaced by $\widehat{\mathrm{GT}}$.
More precisely, for either subgroup $H=G_{\mathbb Q}$ or
$H=\widehat{\mathrm{GT}}$, put
\[
 Q_H=\widehat{\mathrm{GT}}_0/\overline{\langle\!\langle H\rangle\!\rangle}.
\]
There are noncanonical isomorphisms
\[
 \begin{aligned}
 \Hom_{\mathrm{cts}}(Q_H,\mathbb F_2)
       &\simeq\bigoplus_{n\geq1}\mathbb F_2,\\
 Q_H^{\mathrm{ab}}/2Q_H^{\mathrm{ab}}
       &\simeq\prod_{n\geq1}C_2.
 \end{aligned}
 \tag{C.30.3}
\]
where the second isomorphism is topological. Thus the closed normal
closure of the pentagon group has infinite index in the gentle group.
We do not assert that the displayed $\vartheta_p$ form a complete
basis, and we do not determine the full profinite abelianization
$Q_H^{\mathrm{ab}}$.
C.33 supplies an additional order-$13$ quotient, and C.35
proves that it exists for both choices of $H$.
C.34 gives further odd-primary quotients for that arithmetic
normal closure.
\end{statement}
\begin{proof}
We retain the marked generators and normalized representatives of B.7.
Write $\Lambda=\widehat{\mathbb Z}[[\widehat{\mathbb Z}^{\,2}]]$,
with variables $X,Y$. For $g\in\widehat{\mathrm{GT}}_0$ of multiplier
$\lambda_g$, let $B_g\in\Lambda^\times$ be the normalized Alexander
coefficient, in the convention of
\cite[\S2.3, p.~101]{PBC_Furusho2003} and D.39. Thus, if $c_0$ denotes
the chosen generator of $M'$, the actual semilinear action is
\[
 \begin{aligned}
 g(c_0)&=S_{\lambda_g}(X)S_{\lambda_g}(Y)B_g(X,Y)c_0,\\
 S_a(T)&=\frac{T^a-1}{T-1}.
 \end{aligned}
 \tag{C.30.4}
\]
Here the divided power is the continuous Fox derivative; it is obtained
by continuity from the usual geometric sums. The identity
$S_{ab}(T)=S_a(T)S_b(T^a)$ also shows that $S_a(T)$ is a unit for
$a\in\widehat{\mathbb Z}^{\times}$. In the commutator-coordinate
notation used earlier, $B_g=1+(X-1)(Y-1)b_g$.

We first establish the crossed law with its full scope. Let $A_g$ be
the entire coefficient on the right of (C.30.4). Composition of the
actual metabelian automorphisms gives
$A_{gh}=A_gA_h(X^{\lambda_g},Y^{\lambda_g})$. Cancelling the two
$S$ factors by their displayed identity yields
\[
 B_{gh}(X,Y)=B_g(X,Y)
              B_h(X^{\lambda_g},Y^{\lambda_g}).
 \tag{C.30.5}
\]
This calculation applies to every gentle automorphism and every
multiplier; no arithmetic restriction has entered.

Fix $p\in\mathcal S$, choose a primitive cube root $\omega$, and put
$\mathcal O_p=\mathbb Z_p[\omega]$. Since $p\equiv2\pmod3$, this is
the unramified quadratic extension of $\mathbb Z_p$. Its involution
$\iota$ sends $\omega$ to $\omega^2$, and its residue field is
$\mathbb F_{p^2}$. All evaluations below are continuous evaluations
of the coefficient-$p$ factor of $\Lambda$ at roots of order dividing
six. Define the associativity defect
\[
 E_p(g)=
 \frac{B_g(\omega,\omega)B_g(\omega^2,-1)}
      {B_g(\omega,-1)B_g(\omega,-\omega)}
       \in\mathcal O_p^\times.
 \tag{C.30.6}
\]
Every denominator is a unit, since $B_g$ is a unit in the completed
coefficient ring. Equation (C.30.5), and the fact that every multiplier
is odd, imply
\[
 \begin{aligned}
 E_p(gh)&=E_p(g)\,\iota^{\epsilon(\lambda_g)}E_p(h),\\
 \epsilon(\lambda)&=
 \begin{cases}0&\lambda\equiv1\pmod3,\\1&\lambda\equiv2\pmod3.
 \end{cases}
 \end{aligned}
 \tag{C.30.7}
\]
Indeed substitution fixes $-1$ and either fixes or interchanges the two
primitive cube roots. It also commutes with the products of the
arguments appearing in (C.30.6).

For a pentagon element $g$, the beta factorization gives
\[
 B_g(a,b)B_g(ab,c)=B_g(b,c)B_g(a,bc).
 \tag{C.30.8}
\]
We use the exact full-profinite scope: Ihara's primary paper
\cite[introduction, p.~1; \S2, display (6), preprint p.~9]{DE_Ihara2000}
states the factorization for every element of
$\widehat{\mathrm{GT}}$, and gives faithful Fourier evaluation in
$W\otimes\mathbb Q^{\mathrm{ab}}$. The passage to the normalized
Alexander coefficient is the one in D.39. Only roots of order six
prime to $p$ are needed here. Consequently
\[
               E_p(g)=1\qquad(g\in\widehat{\mathrm{GT}}).
 \tag{C.30.9}
\]
In particular this holds on the specified Galois image.

We next prove that (C.30.6) always has norm one. When $\lambda_u=1$,
B.17 gives the integral identity $B_uB_u^*=1$, where the involution
$*$ sends $(X,Y)$ to $(X^{-1},Y^{-1})$. On each of the four evaluations
in (C.30.6), simultaneous inversion is precisely $\iota$: it sends
$\omega$ to $\omega^2$ and fixes $-1$. Therefore
$E_p(u)\iota E_p(u)=1$.
For arbitrary $g$, choose $\sigma\in G_{\mathbb Q}$ with
$\lambda_\sigma=\lambda_g$, using surjectivity of the cyclotomic
character. Put $u=g\sigma^{-1}$. Then $\lambda_u=1$, and (C.30.7)
and (C.30.9), applied to $g=u\sigma$, give $E_p(g)=E_p(u)$.
Thus
\[
 E_p(g)\in U_p:=\ker\bigl(
       N_{\mathcal O_p/\mathbb Z_p}:\mathcal O_p^\times
                       \longrightarrow\mathbb Z_p^\times\bigr)
                 \qquad(g\in\widehat{\mathrm{GT}}_0).
 \tag{C.30.10}
\]
The choice of $\sigma$ is used only to prove this membership; it is
absent from the definition (C.30.6).

Let $\overline E_p(g)$ denote reduction modulo $p$. It satisfies
$\overline E_p(g)^{p+1}=1$. Define $\vartheta_p(g)$ by
\[
       (-1)^{\vartheta_p(g)}
            =\overline E_p(g)^{(p+1)/2}\in\{1,-1\}.
 \tag{C.30.11}
\]
Exponentiation on this abelian norm-one group is multiplicative.
Moreover $\iota$ acts there by inversion, which fixes the resulting
sign. Equation (C.30.7) therefore proves that $\vartheta_p$ is an
ordinary continuous homomorphism to $\mathbb F_2$, with trivial
action on the target. Equations (C.30.9) and (C.30.11) prove its
vanishing on the pentagon group. Changing $\omega$ to $\omega^2$
inverts $E_p(g)$ and leaves (C.30.11) unchanged, so the character
is independent of that choice.

We now give actual full-group witnesses. By B.27 and C.19, for every
prime $p\geq5$ there is
\[
 u_p\in\widehat{\mathrm{GT}}_{0,1}\cap\ker\Delta
\]
whose normalized metabelian coefficient is $1-2e_p$ in the
coefficient-$p$ factor and is $1$ in every other coefficient factor.
Here $e_p$ is the idempotent supported on the two prime-to-$p$
characters which send $X,Y$ to the same primitive cube root.
These are full gentle lifts with retained frames, not merely
metabelian automorphisms. At the four pairs in (C.30.6), their
coefficient-$p$ values are respectively
\[
                         (-1,1,1,1).
 \tag{C.30.12}
\]
Hence $E_p(u_p)=-1$, whereas $E_q(u_p)=1$ for every other
$q\in\mathcal S$. Since $p\equiv1\pmod4$, the integer
$(p+1)/2$ is odd. Consequently
\[
          \vartheta_q(u_p)=\begin{cases}1&q=p,\\0&q\ne p.
                                  \end{cases}
 \tag{C.30.13}
\]
This establishes linear independence, including on the
multiplier-one subgroup with zero $\Delta$.
The set $\mathcal S$ is infinite: apply the Chebotarev theorem to
$\mathbb Q(\zeta_{12})/\mathbb Q$ and the element $5$ of
$(\mathbb Z/12)^\times$; the resulting prime set has density $1/4$
\cite[Chapter V, Theorem 3.23, p.~164]{PBC_MilneCFT}.
The image in (C.30.2) is closed by compactness and contains every
finite-support vector by (C.30.13), so it is the whole product.
If a finite linear relation involves $\Delta$ as well, evaluation
on the $u_p$ first kills all coefficients of the $\vartheta_p$;
C.23 then kills the remaining coefficient of $\Delta$.

For completeness, the upper bound on the quadratic rank is also
exact. A finitely generated free profinite group has only finitely
many open subgroups of each fixed index. Intersecting those of
bounded index gives a countable cofinal family of open characteristic
subgroups. Its continuous automorphism group, and the closed marked
gentle subgroup, are therefore compact and second countable.
There are only countably many clopen subsets of a compact space
with a countable clopen basis: every such subset is a finite union
of basis elements. It follows that there are at most countably
many continuous maps to $\mathbb F_2$. The independent family just
constructed proves that each character space in (C.30.3) has
countably infinite dimension.

Finally, characters to $\mathbb F_2$ separate points of an elementary
abelian profinite quotient, since they do so on every finite quotient.
Choose a countable vector-space basis of the character space.
Evaluation on that basis is injective on
$Q_H^{\mathrm{ab}}/2Q_H^{\mathrm{ab}}$. It is onto every finite
coordinate product by linear independence, hence onto the entire
product by compactness. This proves the topological assertion in
(C.30.3), without identifying our displayed independent family with
that chosen basis.

The exact finite certificate is
\texttt{pbc\_beta\_characters.py}, with its executed output and JSON.
For $p=5,17,29,41,53,89,101,113$ it constructs
$\mathbb F_p[\omega]$, checks every norm-one element and every pair
for the sign homomorphism and its conjugation invariance, and checks
all four values in (C.30.12), the defect $-1$, and the resulting
identity separation matrix. In the ring
$\mathbb F_p[(\mathbb Z/6)^2]$ it independently verifies
$e_p^2=e_p$ and $(1-2e_p)(1-2e_p)^*=1$, using the exact coefficients
\[
 (e_p)_{i,j}=\frac1{36}
 \begin{cases}2&i+j\equiv0\pmod3,\\-1&i+j\not\equiv0\pmod3,
 \end{cases}
 \qquad 0\leq i,j<6.
 \tag{C.30.14}
\]
The infinite-prime assertion and existence of the full witnesses
are supplied by the proof above, rather than inferred from those
finite samples.
\end{proof}


% Root and DE full proof audits passed; root independently replayed the full DAG and original GAP leaves.
\subsection{A third independent iterated descent character}

\begin{statement}{C.31}{Proved}{Three independent iterated parity characters}
There is an actual multiplier-one full gentle element $g_2$ such that
\[
                 \Delta(D(g_2))=0,\qquad
                 \Delta(D^2(g_2))=1.
 \tag{C.31.1}
\]
Consequently the continuous map
\[
 (\Delta,\Delta D,\Delta D^2):
       \widehat{\mathrm{GT}}_0\twoheadrightarrow C_2^3
 \tag{C.31.2}
\]
is onto, and is already onto on the multiplier-one subgroup.
These three characters are linearly independent and vanish
on the specified full Galois image. The arithmetic-normal quotient
therefore has an abelian quotient $C_2^3$.
This does not assert independence of all subsequent iterates, nor
does it determine the restriction of $\Delta$ to the pentagon group.
\end{statement}

\begin{proof}
We first solve one exact finite source-centralization problem, then
lift it to the full profinite modular group. Two normalizations are
essential: the first finite pure abelianization has exponent twelve
but its inertia has order twenty-four, so that normalization cannot
simply be omitted.

\paragraph{The actual marked finite source.}
Let $\Gamma=\widehat{C_2*C_3}=\langle s,r\mid s^2=r^3=1\rangle$,
$\tau=sr$, and $F=\ker(\Gamma\to S_3)$, with the old marking
$x=(rs)^2$, $y=\tau^2$. Write $G_1$ for the degree-$360$ group of
C.27, and $(\beta_1,A_1)$ for its actual corrected pure frames,
retained in \texttt{pbc\_hopf\_joint\_seed.json}. Put
\[
                  v_1=\beta_1,\qquad u_1=A_1^{-1}\beta_1.
 \tag{C.31.3}
\]
Retain the six evaluations of $q_1\varphi\sigma:F\to G_1$,
and also all six signed modulo-three and modulo-four evaluations
of $\varphi\sigma$, where $\sigma$ runs through the six geometric
peripheral permutations. Their actual Schreier-induced modular
representation has degree $2160+432=2592$. Denote its image by
$G=\Gamma/N$, its pure subgroup by $F_2$, and its generators by
$s,r$ again. The kernel is normal in $\Gamma$: the six substitutions
represent the geometric $S_3$ up to inner automorphisms, and the
induced representation retains exactly their simultaneous kernels.

In the ordered pairs
\[
           (s_1,\tau_1),(s_1,r_1^{-1}),
           (\tau_1,s_1),(\tau_1,r_1^{-1}),
           (r_1^{-1},s_1),(r_1^{-1},\tau_1),
\]
prescribe the six $G_1$ values
\[
 (v_1,u_1,v_1^{-1},u_1v_1^{-1},u_1^{-1},v_1u_1^{-1}),
 \tag{C.31.4}
\]
and prescribe identity in every first signed modulo-three and
modulo-four coordinate. The exact finite calculation places this
tuple in the derived subgroup of the actual joint word image.
Thus it is the value of an actual $F'$ word, not just a tuple in a
product. The word-image comparison with $F_2$ is an isomorphism;
let $f_2\in F_2'$ be the resulting element. The completed finite
calculations give
\[
\begin{array}{c|l}
 |G|&2^{146}3^{30}\\
 |F_2|&2^{145}3^{29}\\
 F_2^{\rm ab}&C_{12}^2\\
 \operatorname{ord}(x),\operatorname{ord}(y)&24,24\\
 G^{\rm ab}&C_6\\
 \alpha(s),\alpha(r)&s f_2^{-1},\ f_2r.
\end{array}
 \tag{C.31.5}
\]
Here $\alpha$ is an actual automorphism, fixing $\tau$ and admitting
pure $s$- and $r$-frames. The degree-$2592$ permutation generators,
all tuple coordinates, membership checks and marked comparison are
literal inputs of \texttt{pbc\_iterated\_joint.py} and its GAP
output. The finite portion of
\texttt{pbc\_hopf\_iterated\_joint\_sparse.g} completes the
permutation-to-polycyclic conversion and saves
\texttt{pbc\_iterated\_joint\_pc.json} before starting any Schur
calculation. The later memory failure of that old Schur attempt
does not enter this proof.

The completed triangular basis change
\texttt{pbc\_native\_adapt.py} retains every old-basis vector and
every new power and conjugation word. In the resulting $176$
generators $g_1,\ldots,g_{176}$, the marked elements are
\[
               r=g_1,\quad s=g_2,\quad
               B_s=g_7,\quad B_r=g_5.
 \tag{C.31.6}
\]
The fixed original source is
\texttt{pbc\_iterated\_shortcut\_v2\_original\_source.json},
with SHA256
\[
\texttt{95fb384593bfaf97d468bbc51f11cd797}
\quad\texttt{484a08cb3ec3b954961085b520ffc88}.
\]
It retains $71$ genuine pure $s$-centralizers and nine genuine pure
$r$-centralizers. A further finite-only permutation replay checks
all $176$ power and $15400$ conjugation clauses, before and after
the basis change, together with every marked frame and all these
pure centralizers, directly in the degree-$2592$ group. Its source
and output are \texttt{pbc\_iterated\_source\_permutation\_replay}.

\paragraph{The original integral central obstruction.}
Let $\mathcal F$ be the free discrete group on the $176$ marked
generators, $R=\ker(\mathcal F\to G)$, and
\[
              H=\mathcal F/[\mathcal F,R],\qquad
              K=R/[\mathcal F,R].
 \tag{C.31.7}
\]
The actual Hopf exact sequence identifies
$K\cap H'=H_2(G,\mathbb Z)$ with the torsion subgroup of the
finitely generated abelian group $K$; its remaining quotient is a
subgroup of the free abelian group $\mathcal F^{\rm ab}$.
All assertions obtained here as integral identities also hold on
profinite completion.

Use the canonical positive words $a,b$ for $\alpha(s),\alpha(r)$.
Each pair of words in the three differences
\[
       (B_s s)-(aB_s),\quad (B_r r)-(bB_r),\quad (ab)-(sr)
 \tag{C.31.8}
\]
has the same image in $G$. Interpreting each difference as the
central quotient of the two words and adding in $K$, their sum is
the actual product discrepancy
\[
             T=(B_s sB_s^{-1})(B_r rB_r^{-1})\tau^{-1}
                       \in K\cap H'.
 \tag{C.31.9}
\]
For a retained pure centralizer $c$ of $s$, put
$j_s(c)=csc^{-1}s^{-1}\in K\cap H'$. This is a homomorphism
on the finite centralizer: its values are central. It has exponent
dividing two, since $s^2$ is central in $H$. Likewise the
$r$-centralizer map $j_r$ has exponent dividing three.

C.42 proves that the three-primary part of the original $T$ is
zero. More precisely, its literal original-row identity in the
actual $26$-generator quotient is $T+3j_r(c_{74})=0$ after the
proved homology isomorphism, and $3j_r(c_{74})=0$ already in the
original centralization. We must solve the two-primary part with
the same original target.

\paragraph{A literal original-row identity in the $86$-generator quotient.}
Let $G_{86}$ be the actual quotient of C.43. C.42's integral
bound kills the original two-primary multiplier by $2^{14}$;
C.43's split surjection therefore gives the same bound on
$H_2(G_{86},\mathbb Z_2)$. For its $86$-generator free source,
the central relation module has $3741$ canonical power/conjugation
coordinates. If $K_{86}$ is its actual finitely generated abelian
central kernel, then
\[
 \operatorname{tors}(K_{86})_{(2)}\cap 2^{14}K_{86}=0.
 \tag{C.31.10}
\]
Indeed split $K_{86}$ abstractly into its free part and finite
torsion. Multiplication by $2^{14}$ kills the two-primary torsion,
and an element of the free part cannot be torsion. The odd-primary
part does not intersect the two-primary part.

The fresh calculation first recollects $9000$ consistency identities
using GAP's original positive-word source collector. The older
native-collected attempt failed its original-source comparison and
is not used. The successful fresh reducer gives an exact finite
DAG. Its complete expansion then replaces every auxiliary coordinate
or modulus leaf by its explicit Bezout combination of original
power/consistency rows. A separate literal GAP replay recollects all
$6446$ original leaves occurring in that expanded certificate.
The final identity is
\[
 T_{86}+\sum_i n_i j_s(c_i) +\sum_\ell m_\ell R_\ell
                      =2^{14}z\quad\text{in }\mathbb Z^{3741},
 \tag{C.31.11}
\]
where every $R_\ell$ is an actual source identity, every $c_i$ is
one of the original genuine pure centralizers, and $z$ is an
explicit integer vector. No independent power-word correction,
coordinate relation, modulus relation or $r$-frame correction
occurs in the final witness. The files
\texttt{pbc\_quotient86\_standard2\_literal\_original\_identity.json}
and the corresponding \texttt{\_replay.g/.out} retain every integer,
word and source comparison. The replay completed successfully on
all $6446$ leaves. The complete DAG and exact decoded-byte replay
are retained separately as reproducibility evidence.

Since every $j_s(c_i)$ has exponent two, only the parity of $n_i$
matters. In the original $176$-generator basis the nonzero corrections
are exactly
\[
\begin{split}
 \mathcal I=\{&6,16,27,35,39,41,43,49,\\
              &53,55,63,65,72,79,80,85\}.
\end{split}
 \tag{C.31.12}
\]
Here each $g_i$ for $i\in\mathcal I$ is one of the recorded pure
$s$-centralizers. The corrected quotient target is two-primary:
C.42 already kills the original target's three-primary part,
which remains zero after projection, and every added class is
two-torsion. Thus (C.31.10)--(C.31.11) prove that
\[
                   T+\sum_{i\in\mathcal I}j_s(g_i)
 \tag{C.31.13}
\]
maps to zero in $H_2(G_{86},\mathbb Z_2)$. This deduction uses
only genuine original identities; completeness of the searched
row set or of a guessed relation lattice is not assumed.

\paragraph{The two remaining original homology coordinates.}
C.43 identifies the kernel of
$H_2(G,\mathbb Z_2)\to H_2(G_{86},\mathbb Z_2)$ with
$(\bigwedge^2 A)_Q=C_2^2$, where
\[
       A=\ker(G\to Q)=C_2^{72},\qquad Q=G_{104},\qquad
       N_3=\ker(Q\to G_{86})=C_3^{18}.
\]
Here $A^{N_3}=A_{N_3}=0$, and the extension $G\to Q$ splits.
We need two actual central extensions of $G$ evaluating the two
coordinate forms; their mere existence would not determine the
value of (C.31.13).

The two invariant alternating forms $B_1,B_2$ of C.43 have ranks
$72$ and $36$. Their explicitly verified invariant quadratic
refinements are $q_1,q_2$. For either pair, take the bilinear form
$c$ with diagonal $q(e_i)$, upper triangular entries $B(e_i,e_j)$
and lower triangular entries zero. The finite group
\[
 E_q=A\times C_2,\qquad
 (v,z)(w,t)=(v+w,z+t+c(v,w))
 \tag{C.31.14}
\]
is a central extension of $A$. Its automorphisms fixing the
\emph{chosen} central $C_2$ kernel map onto $O(q)$ with kernel
$A^*$. For the rank-$36$ form the center is larger; no pointwise
fixing of that larger center is asserted.

Pull back this automorphism extension along the actual $Q$-action.
Odd-order averaging gives $(A^*)^{N_3}=0$ and
$H^2(Q,A^*)=0$, so it splits. Explicitly, for a matrix $M$ preserving
$q$, let $q_M$ have polar form $c(Mv,Mw)+c(v,w)$ and zero values on
the basis vectors. Its lifts are
\[
              (v,z)\longmapsto(Mv,z+q_M(v)+l(v)),\quad l\in A^*.
 \tag{C.31.15}
\]
The retained code first solves all order-three and commuting
relations of the $18$ actual $N_3$ generators, then uniquely lifts
the other quotient generators to normalize that complement.
There are $12312$ checked scalar equations of rank $1224$ in the
first step, and $1296$ equations of rank $72$ for each remaining
generator. All equations are derived from the original source
words and rechecked after solving them. Two automorphisms above
the same matrix differ by a linear form, so checking their
discrepancy on all basis vectors checks the whole identity.

The same normalizer construction gives an actual section of
$G\to Q$: a complement over $N_3$ has normalizer mapping onto
$Q$ by coprime complement conjugacy, and the kernel is
$A^{N_3}=0$. Every power and conjugation relation of this section
is additionally checked in the original finite $G$. Thus the
constructed actions give actual central covers $E_q\rtimes Q$ of
the actual $G$. This argument retains the $Q$ coordinate and does
not presume that its action on $A$ is faithful.

For clarity about the numerical convention, vectors are columns
and the right action of $g_i g_j$ is $M_jM_i$. Write a section lift
of $g_j$ as $g_j a_j$. On multiplying by the original $g_j$ in the
central cover, the coordinate update is
\[
 v'=M_jv+a_j,\qquad
 z'=z+q_{M_j}(v)+l_j(v)+c(M_jv,a_j).
 \tag{C.31.16}
\]
This evaluates all the positive words in (C.31.8) and every
centralizer commutator. A separate GAP calculation checks all
$83$ compared word pairs in the original finite group, not merely
their $A$-coordinates. The exact results are
\[
\begin{array}{c|c|c|c|c}
 \text{form}&\text{rank}&\text{three terms of }T&T&
             \text{all }s\text{- and }r\text{-corrections}\\\hline
 B_1&72&(0,0,0)&0&0\\
 B_2&36&(1,0,1)&0&0.
\end{array}
 \tag{C.31.17}
\]
The programs \texttt{pbc\_top\_two\_section},
\texttt{pbc\_top\_two\_cover\_actions}, and
\texttt{pbc\_top\_two\_cover\_values}, together with both original
GAP replays, retain the matrices, linear solutions, word values and
source hashes. Their independent reruns agree on every mathematical
field.

Restriction of a central extension of $A$ to $H_2(A,\mathbb Z)$
is its commutator pairing $B$. A change in quadratic diagonal lies
in the $\operatorname{Ext}^1$ part and does not affect this pairing.
The canonical inclusion in C.43 therefore identifies our two
cover evaluations on the kernel precisely with the two independent
forms $B_1,B_2$. An extension inflated from $Q$ evaluates zero on
that kernel by naturality. Consequently (C.31.17) detects every
remaining kernel element, and (C.31.13) is zero in the original
two-primary multiplier.

It is also zero at three: C.42 kills $T$ there, and each added
$s$-centralizer class is two-torsion. There are no other primary
parts, since $G$ has order a product of powers of two and three.
We have thus proved the single original integral identity
\[
                    T+\sum_{i\in\mathcal I}j_s(g_i)=0
                         \quad\text{in }K\cap H'.
 \tag{C.31.18}
\]
In particular the literal corrected frames are
\[
        \widetilde B_s=g_7\prod_{i\in\mathcal I}^{\rm increasing}g_i,
        \qquad \widetilde B_r=g_5.
 \tag{C.31.19}
\]
Their conjugates of $s,r$ have product exactly $sr$ in $H$.
The order of the displayed correction product is specified; the
centralizer commutator map makes its effect additive.

\paragraph{The modular source centralization and the full lift.}
Put $L=\langle g_2,g_1\rangle\subset H$. Its image is all of $G$,
so $H=LK$ and $L'=H'$ because $K$ is central. The images of
$g_2,g_1$ are independent basis vectors in $H^{\rm ab}=\mathbb Z^{176}$;
since they generate $L$, this proves $L^{\rm ab}=\mathbb Z^2$.
The actual free two-generator centralization maps onto $L$ and
induces the identity on the Hopf multiplier and an isomorphism on
free abelianizations. It is therefore an isomorphism. Hence
\[
              E=L/\langle g_2^2,g_1^3\rangle
                 =\Gamma/[N,\Gamma]
 \tag{C.31.20}
\]
after completion. The killed elements are central and have
independent nonzero abelian vectors, so their subgroup meets
$H'$ trivially. Since $G^{\rm ab}=C_6=\Gamma^{\rm ab}$,
the central kernel of $E\to G$ is exactly the finite Hopf
multiplier. This identifies the actual source centralization;
an arbitrary Schur cover is not being substituted for it.

Although the words in (C.31.19) were written in $H$, each can be
replaced by an element of $L$ with the same finite image. Their
difference is central and does not change any conjugation. Their
finite images are pure, so the chosen elements of $L$ are pure as
well. Equation (C.31.18) gives a homomorphism $\Gamma\to E$ that
conjugates $s,r$ by these frames and fixes $sr$ exactly. Its map on
the quotient is $\alpha$, so it sends $N$ into the central kernel
and kills $[N,\Gamma]$. It therefore factors through an endomorphism
of the actual $E$ in (C.31.20). Hopf naturality identifies its map on the
central kernel with $H_2(\alpha)$; it is consequently an
automorphism on kernel and quotient, hence on $E$.

Apply the modular central criterion proved by the finite trace and
refinement argument of P.3 and B.11. That proof uses only normality
of the retained source kernel, so it applies to this actual $N$.
It gives a full positive peripheral automorphism $\Phi$ of $\Gamma$
covering $\alpha$, retaining the two finite pure frames and the
identity cusp frame. An inner correction from $N$ makes
$\Phi(\tau)=\tau$ literally and changes none of the finite data.
The full torsion relations now give the full gentle two- and
three-cycle relations on restriction to $F$.

\paragraph{The first normalization and the common cusp factor.}
In the temporary marking $x_{\rm new}=\tau^{-2}$,
$y_{\rm new}=s x_{\rm new}s$, the conjugator
$s\Phi(s)$ has abelianization $(a,-a)$. Its finite value is
$f_2^{-1}\in F_2'$, so (C.31.5) implies
$a\in12\widehat{\mathbb Z}$. Normalize by the inner power
$x_{\rm new}^a$. Returning to the old marking replaces the finite
conjugator by
\[
                         f_2\longmapsto y^{-a}f_2x^a.
 \tag{C.31.21}
\]
This produces an actual normalized gentle pair $g_2=(1,f)$.
All first signed modulo-three and modulo-four values remain identity,
because all three corresponding peripheral matrices have orders
dividing twelve. The actual descended-frame formula of C.24
therefore gives $\Delta(D(g_2))=0$.

The change (C.31.21) need not be identity in $G$, whose inertia
has order twenty-four. Evaluate it in the next $G_1$ coordinates.
Since $s_1^a=r_1^a=1$, the first value and cycle-product value become
\[
       v'=\tau_1^{-a}\beta_1,\qquad
       u'=u_1=A_1^{-1}\beta_1.
 \tag{C.31.22}
\]
For the latter equality use the actual cycle $\rho$:
$\rho(y^{-a}f_2x^a)(y^{-a}f_2x^a)
=z^{-a}(\rho(f_2)f_2)x^a$, whose outer factors evaluate to one.
The retained actual pure frames of the first descended modular
automorphism are consequently
\[
              \beta_D=v'=\tau_1^{-a}\beta_1,\qquad
              A_D=v'(u')^{-1}=\tau_1^{-a}A_1.
 \tag{C.31.23}
\]
They have the \emph{same} left factor, a cusp-centralizing pure
power since $a$ is divisible by twelve. Remove it by a full inner
conjugation by $\tau^a$ before normalizing this descended class.
This does not change $D(g_2)$ as an outer class and restores exactly
the two C.27 finite frames, not just the same finite automorphism.
The certificate \texttt{pbc\_iterated\_normalization\_gate.py}
checks both possible values $a=0,12\pmod{24}$ against all six
actual permutation coordinates.

\paragraph{The second normalization and the signed character.}
With these restored frames the unnormalized pure conjugator of
$D(g_2)$ has the finite value $w$ of C.27 in the old marking.
It lies in $F_1'$, where $F_1$ is the pure subgroup of $G_1$.
Here $F_1^{\rm ab}=C_{12}^2$ and both inertia orders are twelve.
Thus the next abelian normalizing exponent is divisible by twelve,
and this second normalization is genuinely invisible in $G_1$.
The normalized actual pair representing $D(g_2)$ therefore retains
exactly C.27's six modulo-four residues and six projective
modulo-three identities.

Apply C.27's signed-three calculation to this actual full pair.
Its two allowed signed patterns have the same equianharmonic bit
one and Gaussian bit zero after the normalized intertwiner is
chosen. Equations (C.27.14)--(C.27.16), now applied to $D(g_2)$,
give $\Delta(D^2(g_2))=1$. This establishes (C.31.1) with both
normalizations and all retained pure frames accounted for.

Finally C.28 supplies the vector $(1,0,0)$, C.27 supplies a vector
$(b,1,c)$, and the present element supplies $(d,0,1)$. They span
$\mathbb F_2^3$, proving (C.31.2). Every character vanishes on
Galois because $D$ fixes the specified Galois image and $\Delta$
vanishes there. All stopped and failed calculations remain retained;
the positive assertion uses only the successful original-row,
actual-cover and full-lifting comparisons above.
\end{proof}


\begin{statement}{C.32}{Proved}{The full metabelian arithmetic-normal quotient}
Let
\[
 \rho_M:\widehat{\mathrm{GT}}_0\longrightarrow\Aut(M)
\]
be the homomorphism obtained by taking the actual normalized
representatives of B.7, and put
$G_M=\rho_M(\widehat{\mathrm{GT}}_0)$.
For either $H=G_{\mathbb Q}$ or
$H=\widehat{\mathrm{GT}}$, let $N_{M,H}$ be the closed normal
subgroup of $G_M$ generated by $\rho_M(H)$.
Then the entire quotient, without further abelianization or
reduction modulo two, is noncanonically isomorphic to
\[
                G_M/N_{M,H}\simeq\prod_{n\geq1}C_2
                \tag{C.32.1}
\]
as a topological group.
In particular every finite continuous quotient of this group is
an elementary abelian two-group. Equivalently, every continuous
homomorphism from $G_M$ to a finite group which kills $\rho_M(H)$
has elementary abelian two-group image.
This statement concerns the actual metabelian image of the full
gentle group. It does not determine the image of $\ker(\rho_M)$ in the
arithmetic-normal quotient of the full gentle group.
C.33 proves that this latter image has a nonzero quotient of order
$13$.
C.34--C.36 supply quotients of unbounded rank at every prime
at least thirteen.
C.37 also supplies nonzero quotients at $5,7,11$.
C.38 strengthens this to arbitrarily large finite $p$-power
quotients of arbitrarily large rank at every $p\geq5$.
\end{statement}
\begin{proof}
The normalization in B.7 is a group section, not merely a choice
of representatives as a set. Thus $\rho_M$ is a continuous
homomorphism into actual automorphisms, and its compact image
$G_M$ is closed. The multiplier of an element of $G_M$ is read
from its scalar action on
$M^{\mathrm{ab}}=\widehat{\mathbb Z}^{\,2}$.
Write $U_M$ for the kernel of this multiplier.

We first check both the group law and its faithfulness on $U_M$.
For a normalized element of multiplier one, B.7 writes
\[
 x\longmapsto x,\qquad
 y\longmapsto(c_0^r)^{-1}y c_0^r,\qquad
 B=1+(X-1)(Y-1)r .
 \tag{C.32.2}
\]
The product $(X-1)(Y-1)$ is a non-zero-divisor in
$\Lambda=\widehat{\mathbb Z}[[\widehat{\mathbb Z}^{\,2}]]$,
as proved in B.7. Hence $B$ uniquely determines $r$ and the
actual normalized automorphism. The all-multiplier crossed law
(C.30.5) reduces, on $U_M$, to ordinary multiplication
$B_{uv}=B_uB_v$. Consequently
\[
                   U_M\hookrightarrow\Lambda^\times,\qquad
                   u\longmapsto B_u
                   \tag{C.32.3}
\]
is an injective homomorphism. In particular $U_M$ is abelian.
We require only the image of full gentle automorphisms here;
we do not replace it by all abstract peripheral units of $M$.

Let $c=(-1,1)$ be the specified complex conjugation. It belongs
to the Galois image and therefore to either subgroup $H$ in the
statement. Its normalized action is
$x\mapsto x^{-1}$, $y\mapsto y^{-1}$.
The entire coefficient of its action on $c_0$ is
$S_{-1}(X)S_{-1}(Y)=X^{-1}Y^{-1}$, whereas its normalized
coefficient is
\[
                         B_c=1.
\]
Since $c^2=1$, applying (C.30.5) twice gives
\[
 B_{\rho_M(c)u\rho_M(c)^{-1}}(X,Y)
                    =B_u(X^{-1},Y^{-1})=B_u^*.
 \tag{C.32.4}
\]
Every $u\in U_M$ comes from a multiplier-one element of the
full gentle group: the multiplier of any preimage is already
determined by its image on $M^{\mathrm{ab}}$.
The necessary part of B.17 therefore applies and says
$B_uB_u^*=1$. By the faithful identification (C.32.3),
(C.32.4) is the equality of actual normalized automorphisms
\[
               \rho_M(c)u\rho_M(c)^{-1}=u^{-1}
                         \qquad(u\in U_M).
 \tag{C.32.5}
\]
This avoids any inference from a nonfaithful commutator
coefficient or from an unspecified inner normalization.

Put $Q=G_M/N_{M,H}$. The cyclotomic multiplier on $H$ is onto
$\widehat{\mathbb Z}^{\times}$, because it is already onto
on $G_{\mathbb Q}$. Given $g\in G_M$, choose $h\in\rho_M(H)$
with the same multiplier. Then $gh^{-1}\in U_M$ and $g$ and
$gh^{-1}$ have the same image in $Q$. Thus the continuous
map $U_M\to Q$ is onto. It follows from the abelianity of
$U_M$ that $Q$ is abelian. Moreover $\rho_M(c)$ maps to
the identity in $Q$. Equation (C.32.5) now gives
$\bar u=\bar u^{-1}$ for every $\bar u\in Q$.
The whole compact quotient $Q$ therefore has exponent two.

It remains to determine its topological rank. Each character
$\vartheta_p$ of C.30 depends only on the normalized coefficient
$B_g$, hence factors continuously through $\rho_M$.
For continuity of this factorization one may use that a
continuous surjection from a compact space to a Hausdorff space
is a quotient map. The characters vanish on $H$, so they descend
to $Q$. The actual full witnesses $u_p$ of C.30 retain
$\vartheta_q(u_p)=\delta_{pq}$ after passage to $G_M$.
They give infinitely many independent continuous characters
of $Q$.

On the other hand $G_M$, and therefore $Q$, is a quotient of
the compact second-countable group
$\widehat{\mathrm{GT}}_0$. The countability argument in C.30
shows that $\Hom_{\mathrm{cts}}(Q,\mathbb F_2)$ has at most
countable dimension. It consequently has dimension exactly
$\aleph_0$. A profinite abelian group of exponent two is
separated by its continuous characters to $\mathbb F_2$:
each nontrivial element survives in a finite quotient, and
that quotient is a finite-dimensional $\mathbb F_2$-vector
space. Choose a countable basis of the character space.
Its evaluation map embeds $Q$ in $\prod_{n\geq1}C_2$.
Every finite projection is onto, by linear independence;
compactness then makes the whole evaluation map onto.
This proves (C.32.1).

The two choices of $H$ give the same abstract topological
isomorphism type. No equality of their closed normal subgroups,
or isomorphism induced by the natural quotient map between
them, is asserted.
\end{proof}


% Root and DE independent audits passed on 2026-11-15.
\begin{statement}{C.33}{Proved}{An arithmetic-trivial character of order thirteen}
There is a continuous surjective homomorphism
\[
 \epsilon_{13}:\widehat{\mathrm{GT}}_0\longrightarrow\mathbb F_{13}
\]
which vanishes on the absolute Galois image and is nonzero on the
multiplier-one subgroup.  It does not factor through the actual
metabelian image of B.7.  Consequently the abelianization of the full
arithmetic-normal quotient has an odd-order quotient; it is neither
an elementary abelian two-group nor a pro-two group.
C.35 further proves that this same character vanishes on the
entire pentagon group.
\end{statement}
\begin{proof}
We first justify the actual local image and its integral filtration.
Let $F_{13}$ be the free pro-$13$ group on $x,y$.  Denote by
$\mathcal G$ the image of the full gentle group in its normalized
automorphism group, and by $\mathcal U$ the kernel of the local
multiplier $\lambda_{13}:\mathcal G\to\mathbb Z_{13}^{\times}$.
For $n\geq2$, let $\mathcal U_n$ consist of the elements whose
normalized conjugator belongs to $\gamma_nF_{13}$.

Put $z=-x-y$ in the integral free Lie algebra $L$ on $x,y$.
For a coefficient ring $R$, let $S_n(R)$ be the degree-$n$ solutions of
\[
\begin{split}
 q(x,y)+q(y,x)&=0,\\
 [x,q(x,y)]&=[x+y,q(-x-y,y)],\\
 q(x,y)+q(y,z)+q(z,x)&=0.
\end{split}                                                    \tag{C33.1}
\]
These are exactly the normalized special outer derivations fixed
by the geometric $S_3$ action.  To check the assertion, write
$D_q(x)=0$, $D_q(y)=[y,q]$.  The second equation says that
$D_q(z)=[z,q(x,y)-q(z,y)]$.  The transposition of $x,y$
changes $D_q$ by $\operatorname{ad}(q)$, and the cycle of $x,y,z$
does the same precisely when the second equation holds.
Its third iterate gives the third equation: the sum of the three
inner conjugators must be central, and the free Lie algebra has
zero center.  Conversely these computations give both required
outer invariances.  No degree-$n$ peripheral centralizer ambiguity
occurs for $n\geq2$.

The following equality concerns the actual completed groups:
\[
 \mathcal U_n/\mathcal U_{n+1}\simeq S_n(\mathbb Z_{13}).
                                                               \tag{C33.2}
\]
Necessity follows by taking the first nonzero Magnus terms of the
peripheral and $S_3$ equations.  For surjectivity, apply the general
integral argument inside P.2 to any $q$ in (C33.1), not just to its
particular degree-seven example.  Lift its three peripheral
conjugators to $F_{13}$.  At each subsequent degree $m$, the error
in their product relation can be changed by
$(u,v)\mapsto[u,x]+[v,y]$ in $L_m$; this map is integrally onto
by Jacobi induction on Lie words.  The corrections converge,
give an automorphism by the pro-$13$ basis theorem, and do not
change the original homogeneous derivation.  Coprime complement
conjugacy for the $S_3$ action then removes all discrepancies in
strictly higher degree.  Thus this is a complete pro-$13$
gentle automorphism.  The relative pronilpotent construction of
P.4 lifts that entire solution to the full gentle group, using
the identity at all other prime components and retaining the
cusp conjugator.  It has global multiplier one.  This proves
(C33.2) for $\mathcal U$ as defined, rather than for a group
of merely truncated solutions.

The exact integral computation gives the following ranks:
\[
\begin{array}{c|rrrrrrrrrrr}
 n&2&3&4&5&6&7&8&9&10&11&12\\ \hline
 \operatorname{rank}S_n&0&1&0&1&0&2&1&4&2&9&7\\
 \operatorname{rank}[S,S]_n&0&0&0&0&0&0&1&0&2&1&6 .
\end{array}                                                    \tag{C33.3}
\]
The two programs \texttt{pbc\_gentle\_even.py} and
\texttt{pbc\_gentle\_even\_fast.py} independently evaluate the
substitutions by associative expansion and recursive Lyndon
brackets.  Their complete mathematical outputs agree.
For every weight, every displayed integral kernel vector is
checked against every original equation row.  An independent
rank lower bound modulo $65521$ equals the complementary
dimension, proving the rational ranks.  The further certificate
\texttt{pbc\_gentle\_weight12\_integral.py} checks a unit minor
modulo $13$ for each kernel basis.  Since the full integral
kernel is saturated, these are bases over $\mathbb Z_{13}$.
It also computes the Ihara bracket
\[
 \{f,g\}=[f,g]+D_f(g)-D_g(f)
\]
and checks every bracket against all the defining rows.
In the retained seven-element basis of $S_{12}$, the primitive
integral functional
\[
             \ell=(-4,1,8,-4,-17,14,84)                         \tag{C33.4}
\]
annihilates every degree-twelve bracket.  Multiplying its reduction
by $3$ gives $(1,3,11,1,1,3,5)$, so the first basis vector has
value $1$.  Thus this is an integral witness, with no inference
from rational bracket membership modulo $13$.

Now put $U=\mathcal U/\mathcal U_{13}$.
The zero degree-two space in (C33.3) implies
$\mathcal U=\mathcal U_3$.  The commutator filtration gives
$[U_a,U_b]\subseteq U_{a+b}$, so $U$ has nilpotency class at most
four.  It is a finitely generated torsion-free pro-$13$ group:
its finite filtration has free finite-rank $\mathbb Z_{13}$
quotients.  Its Lazard Lie lattice
$\mathfrak u=\log U$ is defined integrally, preserves this
filtration and its associated graded brackets, and is functorial
for all automorphisms.
For a direct denominator check, use the action on the integral
Magnus algebra modulo augmentation degree $14$.
For $a\in U$, $a-1$ raises augmentation degree by at least three.
The operator logarithm and exponential therefore stop after
four terms; their denominators are units in $\mathbb Z_{13}$.
The Lazard correspondence applies to the actual closed subgroup
of these operators, since its class is less than $13$.
No logarithm of a free-group generator, and no class-thirteen
Lazard assertion about the free group, is used here.
See \cite[Theorem 1]{PBC_Garaialde2014} for the precise
finitely generated pro-prime correspondence, including its
subgroup and automorphism assertions.

Set $G=\mathcal G/\mathcal U_{13}$, so that
\[
       1\longrightarrow U\longrightarrow G
        \xrightarrow{\lambda_{13}}\mathbb Z_{13}^{\times}
        \longrightarrow1 .                                    \tag{C33.5}
\]
Its multiplier is onto because the Galois multiplier is onto.
In the local Galois image, the inverse image of the Teichmuller
group $\mu_{12}$ is an extension of $\mu_{12}$ by a pro-$13$
group.  The coprime complement theorem and compactness give an
element $t$ of order $12$ with primitive multiplier
$\omega\in\mu_{12}$.
Conjugation by $t$ acts on
$\operatorname{gr}_n\mathfrak u=S_n(\mathbb Z_{13})$
as $\omega^n$.  The ten possible weights $3,\ldots,12$
have distinct residues modulo $12$.  Hence the integral
idempotents
\[
 e_n=\frac1{12}\sum_{j=0}^{11}\omega^{-nj}\operatorname{Ad}(t)^j
\]
split the filtered lattice, and
$e_n\mathfrak u\to\operatorname{gr}_n\mathfrak u$ is an
isomorphism.  There can be no kernel or cokernel in this map:
their successive graded pieces would have to carry both
the eigencharacter $\omega^n$ and a distinct $\omega^m$.
The Lie bracket respects this splitting.  If $a+b\leq12$,
the bracket of the $a$ and $b$ summands is exactly the
degree-$a+b$ Ihara bracket; if $a+b>12$ it is zero.
This removes the possibility of hidden higher filtered
commutator relations.

Extend $\ell$ by zero on all eigensummands except weight $12$.
It annihilates $[\mathfrak u,\mathfrak u]$.  Its reduction,
normalized as after (C33.4), therefore gives a nonzero
continuous homomorphism
\[
                         \theta:U\longrightarrow\mathbb F_{13}.
                                                               \tag{C33.6}
\]
This character is invariant under the entire $G$.
Indeed conjugation on
$\mathfrak u/[\mathfrak u,\mathfrak u]$ factors through
$G/U=\mathbb Z_{13}^{\times}$ and hence commutes with $t$.
It preserves the weight-twelve eigensummand.  That summand
is the top filtered piece, where an element of multiplier
$a$ acts exactly by $a^{12}$, which is $1$ modulo $13$.
This proves invariance without assuming that a chosen wild
multiplier lift commutes with $t$ inside $G$.

There is no extension obstruction:
\[
 H^2_{\mathrm{cts}}(\mathbb Z_{13}^{\times},\mathbb F_{13})=0.
                                                               \tag{C33.7}
\]
The prime-to-$13$ factor $\mu_{12}$ has vanishing positive
cohomology by averaging; the procyclic factor
$1+13\mathbb Z_{13}\simeq\mathbb Z_{13}$ has
$13$-cohomological dimension one.  Equivalently, the
inflation maps in degree two along its finite cyclic
quotients multiply by $13$, so their direct limit is zero.
The continuous five-term sequence of (C33.5) now extends
$\theta$ to a continuous homomorphism
$\widetilde\theta:G\to\mathbb F_{13}$.

Finally remove its arithmetic restriction integrally.
The normalized pro-$13$ action on the tripod, with its
standard tangential basepoint, is unramified away $13$;
this is stated for the actual pro-prime automorphism
representation in \cite[\S8.5, immediately before
Theorem 8.7]{A_HainMatsumoto2003}, taking $N=1$.
Thus $\widetilde\theta|_{G_{\mathbb Q}}$ factors through
$G_{\mathbb Q,\{13\}}$.
By Kronecker--Weber, the maximal abelian extension of
$\mathbb Q$ unramified away $13$ is
$\mathbb Q(\mu_{13^\infty})$; see
\cite[Chapter I, Theorem 4.16 and its proof]{PBC_MilneCFT}.
Consequently
$\operatorname{Hom}_{\mathrm{cts}}
(G_{\mathbb Q,\{13\}},\mathbb F_{13})$ is one-dimensional,
generated by the nonzero cyclotomic character
\[
       c_{13}(\lambda)=\frac{\lambda_{13}^{12}-1}{13}\pmod {13}.
\]
This expression is an ordinary additive character on the
whole multiplier group: multiplication of its two numerators
has cross term divisible by $13^2$.
There is therefore $a\in\mathbb F_{13}$ such that
\[
 \epsilon_{13}=\widetilde\theta\circ\rho_{13}
                  -a\,c_{13}\circ\lambda
                                                               \tag{C33.8}
\]
vanishes on $G_{\mathbb Q}$.  The adjustment is zero on
global multiplier one.  The actual complete lift supplied
by (C33.2) for the first weight-twelve basis vector still
has value $1$, so (C33.8) is surjective.

If it factored through the actual metabelian image, it
would factor through that image's arithmetic-normal
quotient.  C.32 proves that quotient has exponent two,
which has no nonzero homomorphism to $\mathbb F_{13}$.
This contradiction proves the last assertion.
\end{proof}


% Root and DE independent audits passed on 2026-11-15.
\begin{statement}{C.34}{Proved}{Odd-primary arithmetic-trivial characters at all sufficiently large primes}
Let $V$ be the two-dimensional standard representation of $S_3$.
For $g\in S_3$, put
\[
 L_n(g)=\frac1n\sum_{d\mid n}\mu(d)
                   \operatorname{tr}(g^d\mid V)^{n/d}.
\]
Write $1,\tau,\rho$ for the identity, a transposition and a
three-cycle, and define, for $n\geq2$,
\[
\begin{split}
 a_n&=\frac{2L_n(1)-L_{n+1}(1)-2L_n(\rho)
                 -3L_{n+1}(\tau)-2L_{n+1}(\rho)}6,\\
 b_n&=\sum_{\substack{i+j=n\\2\leq i<j}}a_i a_j
              +\begin{cases}\binom{a_{n/2}}2,&2\mid n,\\0,&2\nmid n,\end{cases}\\
 d_n&=\max(0,a_n-b_n).
\end{split}                                                     \tag{C34.1}
\]
Let $Q$ be the full gentle group modulo the closed normal closure
of the absolute Galois image.  For every prime $p\geq5$,
\[
 \dim_{\mathbb F_p}\operatorname{Hom}_{\mathrm{cts}}(Q,\mathbb F_p)
                       \geq d_{p-1}.                            \tag{C34.2}
\]
Every displayed character can be chosen with its nonzero values
witnessed on global multiplier one.  In particular,
\[
\begin{array}{c|rrrrrr}
 p&13&17&19&23&29&31\\ \hline
 d_{p-1}&1&37&143&1805&77519&273536 .
\end{array}                                                     \tag{C34.3}
\]
The lower bound is positive and tends to infinity for every
sufficiently large prime $p$.  There is a continuous surjection
\[
 Q^{\mathrm{ab}}\longrightarrow
 \left(\prod_{j\geq1}C_2\right)
       \times\prod_{\substack{p\geq5\\p\ {\rm prime}}}
                         (\mathbb Z/p\mathbb Z)^{d_{p-1}} .
                                                               \tag{C34.4}
\]
This determines a large quotient, not the full abelianization.
C.36 makes the positivity threshold explicit: every prime
$p\geq13$ has a nonzero displayed factor.
For the general odd-primary family, vanishing on the pentagon
group is not asserted.  C.35 proves that vanishing for the
order-$13$ character.
\end{statement}
\begin{proof}
We give an independent character computation of the dimensions
used in C.33.  Work first over $\mathbb Q$.
The free Lie algebra on $V$ has degree-$n$ piece $L_n(V)$.
A homogeneous tangential derivation is specified by three
conjugators $q_x,q_y,q_z\in L_n(V)$, with $x+y+z=0$.
For $n\geq2$ these conjugators have no ambiguity, because
the homogeneous peripheral centralizer in that degree is zero.
The special condition is
\[
            [x,q_x]+[y,q_y]+[z,q_z]=0 .
\]
Inner derivations are the diagonal triples.  The permutation
representation on the three conjugator positions modulo its
diagonal is $V$.  Bracketing with $x,y$ is onto $L_{n+1}(V)$
by Jacobi induction.  There is therefore an exact sequence
of $S_3$-representations
\[
 0\longrightarrow \mathfrak s_n
 \longrightarrow V\otimes L_n(V)
 \longrightarrow L_{n+1}(V)\longrightarrow0,                    \tag{C34.5}
\]
where $\mathfrak s_n$ is the space of special outer derivations.
Its $S_3$ invariants are exactly $S_n(\mathbb Q)$ of (C33.1).

For completeness, the trace formula in the statement follows
directly from the tensor algebra and its ordered Lie-monomial
basis.  The trace series of the tensor algebra at $g$ is
$(1-\operatorname{tr}(g\mid V)t)^{-1}$.
The Poincare--Birkhoff--Witt decomposition into symmetric
powers of the $L_n(V)$, followed by logarithms, gives
\[
 \operatorname{tr}(g\mid V)^m
       =\sum_{n\mid m}n\,\operatorname{tr}(g^{m/n}\mid L_n(V)).
\]
Mobius inversion gives $L_n(g)=\operatorname{tr}(g\mid L_n(V))$.
Since the three values of $\operatorname{tr}(g\mid V)$ are
$2,0,-1$, taking invariants in (C34.5) gives exactly $a_n$.
Thus $a_n$ is the dimension of the actual degree-$n$
gentle Lie space.
There are at most $b_n$ independent degree-$n$ brackets:
for $i<j$ there are $a_i a_j$ possible pairs, and for $i=j$
there are $\binom{a_i}{2}$.  Nested brackets are already
included, since their two outer factors are elements of
the same graded Lie algebra.  We have not assumed that
all these pairs are independent.

Fix $p\geq5$ and apply C.33's actual-image argument with
$13$ replaced by $p$ and top conjugator degree $p-1$.
The complete integral peripheral correction and coprime
$S_3$ argument give
\[
       \mathcal U_n/\mathcal U_{n+1}=S_n(\mathbb Z_p)
              \qquad(2\leq n\leq p-1).                          \tag{C34.6}
\]
These are free finite-rank $\mathbb Z_p$ modules, of ranks $a_n$.
The zero degree-two space again gives $\mathcal U=\mathcal U_3$.
On the Magnus algebra modulo augmentation degree $p+1$,
the operator $u-1$ raises degree by at least three.
The resulting actual quotient
$U=\mathcal U/\mathcal U_p$ has nilpotency class at most
$\lfloor(p-1)/3\rfloor<p$ and is a finitely generated
torsion-free pro-$p$ group.  Lazard logarithms, their
filtration and the leading Ihara bracket are integral.
This uses the small class of the automorphism group,
not a class-$p$ correspondence for the underlying free group.

A tame order-$(p-1)$ element in the local Galois image
has distinct eigencharacters on every possible grade
$3,\ldots,p-1$.  Its averaging idempotents split the
actual filtered lattice as in C.33.  The invariant
eigensummand is precisely the top piece of rank $a_{p-1}$.
Its intersection with the Lie bracket ideal has rank
at most $b_{p-1}$.  Hence the lattice of integral
functionals on that top piece which annihilate brackets
has rank at least $d_{p-1}$.
This lattice is saturated in the dual of the top piece:
if $p\ell$ annihilates every bracket, so does $\ell$,
because $\mathbb Z_p$ is torsion-free.
Consequently a basis supplies at least $d_{p-1}$
linearly independent reductions modulo $p$.
This saturation argument is what permits reduction
modulo $p$; rational dimension alone would not suffice.

Extending each functional by zero on the other tame
eigensummands gives an ordinary homomorphism
$U\to\mathbb F_p$.  It is invariant under the whole
local gentle image: conjugation on the abelianization
factors through the multiplier and acts on the top
summand by $\lambda_p^{p-1}=1$ modulo $p$.
The vanishing of
$H^2_{\mathrm{cts}}(\mathbb Z_p^\times,\mathbb F_p)$
extends each character to the full local image.
Its Galois restriction is unramified away $p$.
Kronecker--Weber makes that restriction a scalar multiple
of
\[
               c_p(\lambda)=
                    \frac{\lambda_p^{p-1}-1}{p}\pmod p .
                                                               \tag{C34.7}
\]
Subtracting this scalar multiple makes the character
arithmetic-trivial.  The subtraction vanishes on
global multiplier one, where (C34.6) and P.4 provide
actual complete witnesses for the independent values.
Thus their independence is retained.  This proves
(C34.2), with no assumption about generation of the
arithmetic Lie algebra or its integral image.

We prove the assertion about sufficiently large primes.
Put $l_n=L_n(1)$.  The trace formula gives
\[
 l_n=\frac{2^n}{n}+O(2^{n/2}),\qquad
 a_n=\frac{2^n}{3n(n+1)}+O(2^{n/2})
          \sim\frac{2^n}{3n^2}.                                \tag{C34.8}
\]
For the nonidentity classes the largest exponential
terms have sizes $2^{n/2}$ and $2^{n/3}$, so they are
included in the displayed error.
In particular $r_n=a_n2^{-n}\sim(3n^2)^{-1}$.
The sum $S=\sum_{n\geq2}r_n$ is finite and small.
Indeed (C34.5) gives $a_n\leq2l_n-l_{n+1}$, and
\[
 \sum_{n=2}^{N}\frac{2l_n-l_{n+1}}{2^n}
              =\frac12-\frac{l_{N+1}}{2^N}
                 \longrightarrow\frac12.
                                                               \tag{C34.9}
\]
At $n=5$ this inequality is strict by two dimensions,
so in fact $S\leq7/16<1/2$.

The elementary convolution estimate for
$r_n\sim(3n^2)^{-1}$ is
\[
 \frac{\sum_i r_i r_{n-i}}{r_n}\longrightarrow2S.                \tag{C34.10}
\]
To verify it, first retain only $i\leq M$ and $n-i\leq M$.
Their ratio tends to $2\sum_{i\leq M}r_i$.
On the remaining half interval $M<i\leq n/2$,
$r_{n-i}/r_n$ is uniformly bounded, while
$\sum_{i>M}r_i=O(1/M)$.  The other half is identical;
then let $M$ tend to infinity.
The diagonal correction in $b_n$ is exponentially
smaller than $a_n$.  Therefore
\[
 \frac{b_n}{a_n}\longrightarrow S<\frac12,\qquad
 d_n\sim(1-S)\frac{2^n}{3n^2}.                                  \tag{C34.11}
\]
It follows that $d_{p-1}$ is positive and unbounded
for every sufficiently large prime.
The exact formula program
\texttt{pbc\_gentle\_dimension\_formula.py} verifies
(C34.3), agrees with all independently computed
dimensions through degree twelve, and records additional
counts through degree $300$.  Those finite computations
illustrate (C34.11); they are not the proof of its
eventual assertion.

For each prime choose the independent characters just
constructed.  They jointly surject onto the indicated
finite elementary abelian $p$-group.  Any finite product
over distinct primes is still a quotient: a subdirect
product of finite abelian groups of pairwise coprime
orders is their product.  The same argument includes
arbitrary finite subsets of C.30's independent
quadratic characters.  Compactness now gives
surjectivity onto the full product in (C34.4).
\end{proof}


% Root and DE independent proof/code audits passed on 2026-11-15.
\begin{statement}{C.35}{Proved}{The order-thirteen character vanishes on the full pentagon group}
The character $\epsilon_{13}$ constructed in C.33 can be
chosen exactly as there so that
\[
                    \epsilon_{13}|_{\widehat{\mathrm{GT}}}=0.
\]
Consequently the quotient of the full gentle group by the
closed normal subgroup generated by the full pentagon group
has a continuous quotient
\[
                    C_{13}\times\prod_{j\geq1}C_2.
\]
The order-thirteen character remains nonzero on global
multiplier one and does not factor through the metabelian
image.  This assertion gives no nonzero arithmetic-trivial
character of the pentagon group itself.
\end{statement}
\begin{proof}
Use the free-fibre representation of the infinitesimal
four-strand braid Lie algebra and the signed pentagon
defect $\mathcal D(q)$ from P.2.  In particular its
vanishing is necessary for the first nonzero term of
an actual pro-$13$ pentagon solution.
For $q\in S_{12}$, let $R_i(q)$ be the coefficient in the
following table:
\[
\begin{array}{c|c|l}
 i&\text{generator to which }\mathcal D(q)\text{ is applied}
       &\text{output word}\\ \hline
1&a&\mathtt{aaaaaaabcbcbc}\\
2&a&\mathtt{acbaacaaabcba}\\
3&b&\mathtt{aaaabacaacabb}\\
4&b&\mathtt{baaacacbaabac}\\
5&a&\mathtt{abbacbaaacaca}
\end{array}                                                     \tag{C35.1}
\]
In the same integral seven-element basis as C.33, exact
integer expansion gives
\[
\begin{split}
698922\,\ell={}&-192329R_1+472R_2-26842R_3\\
               &-107926R_4-11870R_5 ,
\end{split}                                                     \tag{C35.2}
\]
where $\ell=(-4,1,8,-4,-17,14,84)$.
Thus every degree-twelve linear pentagon solution is
annihilated by $\ell$.  The coefficient $698922$ is a
unit modulo $13$, although the argument below only
needs the identity over $\mathbb Q_{13}$.

Here are the exact verification details for (C35.2).
The independent program \texttt{pbc\_weight12\_pentagon.py}
computes a specified coefficient by contracting adjacent
output-letter pairs.  If a linear derivation $D$ takes
a letter $v$ to a two-letter polynomial, then
\[
 [w]D(P)=
 \sum_{j,v}[w_jw_{j+1}]D(v)\,
       [w_1\cdots w_{j-1}v w_{j+2}\cdots w_m]P.                 \tag{C35.3}
\]
Iterating this exact identity computes the coefficient
of each associative operator word in $q(D,E)$.
The program first compares nineteen independently
computed coefficients with the full degree-seven
pentagon expansion retained in P.2, including its
nonzero coefficient $-16$.  It then finds the five
rows modulo $65521$, recomputes each selected row
over $\mathbb Z$, and verifies (C35.2) on every one
of the seven integral basis vectors.
The entire five-by-seven integer matrix, its coefficient
labels, source hash and exact row combination are retained
in the JSON certificate.  The modular search alone is
not used as a proof of the identity.

It remains to compare this pure-degree certificate with
an arbitrary actual pentagon element, whose first nonzero
term need not have degree twelve.
Let $G,U,\mathfrak u,t$ be precisely the actual local
quotient, its multiplier-one subgroup, its Lazard lattice
and its tame Galois element from C.33.
Let $P$ be the image of the full pentagon group in $G$,
and put $P_1=P\cap U$.
It is a closed subgroup of the finitely generated
nilpotent pro-$13$ group $U$ and is finitely generated;
for example intersect its finite central filtration
with the finite-rank $\mathbb Z_{13}$ graded pieces.
The subgroup part of the Lazard correspondence identifies
$\mathfrak p=\log P_1$ with an actual Lie sublattice of
$\mathfrak u$.

Crucially, $t$ belongs to the Galois image and hence to
$P$.  It normalizes $P_1$ and preserves $\mathfrak p$.
Therefore the integral projector
\[
                     e_{12}=\frac1{12}
                         \sum_{j=0}^{11}\operatorname{Ad}(t)^j
                                                               \tag{C35.4}
\]
preserves $\mathfrak p$ as well.
For any $h\in P_1$, the element
$v=e_{12}\log h$ belongs to
$\mathfrak p\cap\mathfrak u_{12}$.
Its exponential is a genuine element of $P_1$,
and hence is the image of an actual full pentagon pair.
All its local conjugator terms of degrees below twelve
are zero.  Its degree-twelve term $q$ therefore satisfies
the linear pentagon, with no contribution from lower
terms or their higher corrections.
By (C35.2), $\ell(q)=0$.
The character $\theta$ on $U$ in C.33 is precisely the
normalized reduction of $\ell(e_{12}\log h)$.
We have proved
\[
                             \theta(P_1)=0.                    \tag{C35.5}
\]

Let $\widetilde\theta:G\to\mathbb F_{13}$ be the
continuous extension used in C.33.  Equation (C35.5)
shows that its restriction to $P$ factors through
the local multiplier.
For an even more direct comparison, given $g\in P$
choose a Galois element $\sigma\in P$ with
$\lambda_{13}(\sigma)=\lambda_{13}(g)$.
Then $g\sigma^{-1}\in P_1$, so
$\widetilde\theta(g)=\widetilde\theta(\sigma)$.
The cyclotomic correction in (C33.8) was chosen to
remove this value on every Galois element.
It consequently removes the value on every $g\in P$.
This is the same character $\epsilon_{13}$,
with the same multiplier-one witness.

Finally combine $\epsilon_{13}$ with the independent
quadratic characters of C.30.  Each finite joint
projection is onto by coprime orders and the stated
quadratic independence.  Compactness gives the displayed
product quotient.  The character is zero on the entire
pentagon group, so it makes no assertion about that
group's own arithmetic-trivial character space.
\end{proof}


% Root independent proof/code audit passed on 2026-11-15.
\begin{statement}{C.36}{Proved}{A nonzero arithmetic-trivial character at every prime at least thirteen}
The lower bound $d_{p-1}$ in C.34 is positive for every
prime $p\geq13$.  Hence the full gentle arithmetic-normal
quotient has a nonzero continuous character to
$\mathbb F_p$ for every such prime, and its abelianization
has the simultaneous product quotient of C.34 with
nontrivial factors at all these primes.
No absence assertion is made for smaller odd primes.
\end{statement}
\begin{proof}
We use the exact quantities $a_n,b_n,d_n$ of C.34.
The independent program
\texttt{pbc\_gentle\_uniform\_prime\_bound.py}
uses a Mobius sieve and the proved character formula
to compute all relevant integers.
It agrees with the previous implementation through
degree $300$ and verifies
\[
                 a_n>b_n\quad
                   (n=12,14,\ldots,998).                        \tag{C36.1}
\]
All $494$ exact rows are retained.  The following
analytic estimate proves every remaining degree,
rather than extrapolating this finite check.

Put $r_n=a_n2^{-n}$.  The trace formula implies, for
$n\geq2$,
\[
 \left|r_n-\frac1{3n(n+1)}\right|\leq3\,2^{-n/2}.                \tag{C36.2}
\]
Indeed the nonleading divisors in $L_n(1)$ contribute
at most $2^{n/2}$; the transposition trace has absolute
value at most $2^{n/2}$, and the three-cycle trace
has absolute value at most $2^{n/3}+1$.
Substitution in (C34.1), including the terms at $n+1$,
gives (C36.2).
For every $n\geq100$,
\[
 3\,2^{-n/2}\leq\frac1{100n^2}.
                                                               \tag{C36.3}
\]
The inequality holds at $n=100$, and
$n^2 2^{-n/2}$ decreases thereafter, since
$(101/100)^4<2$.
Combining these estimates gives the convenient rational
bounds
\[
             \frac8{25n^2}\leq r_n\leq\frac7{20n^2}
                         \qquad(n\geq100).                     \tag{C36.4}
\]
For the lower bound use
$100/303-1/100>8/25$; for the upper bound use
$1/3+1/100<7/20$.

By definition of the unordered bracket count,
\[
 \frac{b_n}{a_n}
    \leq\sum_{2\leq i\leq n/2}
                       r_i\,\frac{r_{n-i}}{r_n}.               \tag{C36.5}
\]
This slightly overcounts the middle term when $n$ is even,
which is harmless.
If $n\geq1000$ and $i\leq100$, (C36.4) gives
\[
 \frac{r_{n-i}}{r_n}
   \leq\frac{35}{32}\left(\frac{n}{n-i}\right)^2
   \leq\frac{875}{648}.
\]
The sum of these endpoint terms is at most
$(875/648)S$, where C.34 proves
$S=\sum r_i\leq7/16$.
For $100<i\leq n/2$, the same ratio is at most $35/8$,
and
\[
            \sum_{i>100}r_i
               \leq\frac7{20}\sum_{i>100}i^{-2}
               \leq\frac7{2000}.
\]
Consequently, for every $n\geq1000$,
\[
 \frac{b_n}{a_n}
 \leq \frac{875}{648}\frac7{16}
           +\frac{35}{8}\frac7{2000}
       =\frac{78547}{129600}<1.                                \tag{C36.6}
\]
Equations (C36.1) and (C36.6) prove $d_n>0$ at every
even $n\geq12$.  Set $n=p-1$ and apply the already
proved actual-character construction of C.34.
\end{proof}


% Root and FGH independent proof/code audits passed on 2026-11-15.
\begin{statement}{C.37}{Proved}{The remaining primes five, seven and eleven}
Let $Q$ be the quotient of the full gentle group by the
closed normal subgroup generated by its absolute Galois image.
Then
\[
\begin{array}{c|rrr}
p&5&7&11\\ \hline
\dim_{\mathbb F_p}\operatorname{Hom}_{\mathrm{cts}}
             (Q,\mathbb F_p)&\geq1&\geq1&\geq513 .
\end{array}                                                     \tag{C37.1}
\]
Every displayed independent family has its independent
values witnessed on global multiplier one.
Together with C.34--C.36, this gives a nonzero continuous
arithmetic-trivial character of the full gentle group
to $\mathbb F_p$ for every prime $p\geq5$.
The general odd-primary families in this statement are
not asserted to vanish on the pentagon group.
\end{statement}
\begin{proof}
We allow repeated tame eigencharacters, which were avoided
in C.34 by truncating at degree $p-1$.
For the three primes choose
\[
                 (p,m)=(5,12),(7,12),(11,20).
                                                               \tag{C37.2}
\]
Use the actual local image $\mathcal G$, its
multiplier-one subgroup $\mathcal U$, and its conjugator
filtration from C.33, now at the indicated prime.
The completed P.2 correction and P.4 lifting give
$\mathcal U_n/\mathcal U_{n+1}=S_n(\mathbb Z_p)$
in every degree used here.
Set $U=\mathcal U/\mathcal U_{m+1}$ and
$G=\mathcal G/\mathcal U_{m+1}$.
Their integral Magnus operators raise degree by at
least three, so $U$ has nilpotency class at most
$\lfloor m/3\rfloor$, respectively $4,4,6$.
Each is strictly less than $p$, and the actual
Lazard Lie lattice $\mathfrak u=\log U$ is integral.

Let $t$ be a tame order-$(p-1)$ element in the local
Galois image, with primitive multiplier $\omega$.
Its averaging projector onto eigenvalue one is
\[
                 e=\frac1{p-1}
                       \sum_{j=0}^{p-2}\operatorname{Ad}(t)^j .
\]
Although the same eigencharacter can occur in several
degrees, we can choose an integral eigenbasis adapted
to the filtration.  Indeed averaging splits each
eigenspace, and each successive filtered quotient
is a free $\mathbb Z_p$ module; lift a basis of the
quotient and project to its required eigencharacter.
Thus each basis vector has a specified first degree
$n$, of multiplicity $a_n$, and eigencharacter
$\omega^n$.
It follows that
\[
 \operatorname{rank}e\mathfrak u
        =\sum_{\substack{2\leq n\leq m\\p-1\mid n}}a_n.
                                                               \tag{C37.3}
\]

The following count concerns actual filtered brackets,
not only their leading terms.
For two such basis vectors with first degrees $i,j$,
their bracket has eigencharacter $\omega^{i+j}$
and filtration at least $i+j$.  It can contribute to
$e[\mathfrak u,\mathfrak u]$ only if
$p-1\mid i+j\leq m$.
There is at most one resulting vector for each pair
of basis vectors.  Antisymmetry removes the diagonal
pair.  Hence
\[
 \operatorname{rank}e[\mathfrak u,\mathfrak u]
       \leq\sum_{\substack{2\leq n\leq m\\p-1\mid n}}b_n.
                                                               \tag{C37.4}
\]
Higher terms of a lower-degree bracket can change
that one vector, but cannot increase this count.
The Lie derived ideal is the linear span of all
such pair brackets, so this also counts every
possible hidden filtered relation.

Here are all the necessary dimensions:
\[
\begin{array}{c|c|r|r|r}
p&\text{nonzero invariant degrees}
 &\operatorname{rank}e\mathfrak u
 &\text{bracket bound}&\text{difference}\\ \hline
5&8,12&8&7&1\\
7&12&7&6&1\\
11&10,20&820&307&513 .
\end{array}                                                     \tag{C37.5}
\]
The omitted invariant degrees $4,6$ have zero space.
At degree twenty, the exact character formula of C.34
gives $a_{20}=818$ and $b_{20}=305$.
The program \texttt{pbc\_small\_prime\_characters.py}
checks these rows and the class bounds, retains the
input hashes, and independently recomputes the
lower-degree bracket ranks used next.

The lower invariant degrees cause no nontrivial
weight extension in the rational abelianization.
For $p=5$, the degree-eight space has dimension one
and is the span of the bracket of degrees three and
five.  For $p=7$ there is no nonzero lower invariant
degree.  For $p=11$, the degree-ten space has
dimension two and is spanned by the brackets of
degrees three and seven.
These are the exact bracket ranks in (C33.3).
Lift their factors to the appropriate tame eigenspaces
of the actual lattice.
The resulting actual brackets span each indicated
lower invariant graded piece over $\mathbb Q_p$.
For any invariant lattice vector, subtract a rational
combination of those actual brackets to cancel its
first lower component.  Continue if necessary.
There is no other lower invariant graded piece.
Consequently the natural map
\[
 \mathfrak u_m\otimes\mathbb Q_p
 \longrightarrow
      e\bigl(\mathfrak u/[\mathfrak u,\mathfrak u]\bigr)
                      \otimes\mathbb Q_p
                                                               \tag{C37.6}
\]
is onto.

Let $R$ be the integral lattice of functionals
$\ell:\mathfrak u\to\mathbb Z_p$ which annihilate
$[\mathfrak u,\mathfrak u]$ and are invariant under
$t$.  Equations (C37.3)--(C37.5) give ranks at least
$1,1,513$.
This is a saturated submodule of the integral dual
of $\mathfrak u$: if $p\ell$ satisfies both defining
linear conditions, so does $\ell$.
Thus an integral basis of $R$ remains linearly
independent after reduction modulo $p$.

Each such reduction is invariant under the whole $G$.
Conjugation on the Lie abelianization factors through
$G/U=\mathbb Z_p^\times$ and hence commutes with $t$.
By (C37.6), its invariant part over $\mathbb Q_p$
is generated by the top degree, where multiplier
$a$ acts exactly by $a^m$.
Therefore, for $\ell\in R$ and $g\in G$,
\[
          \ell(\operatorname{Ad}(g)v)=
                         \lambda_p(g)^m\ell(v).                 \tag{C37.7}
\]
This equality first holds over $\mathbb Q_p$ and
hence holds in the torsion-free integral target
$\mathbb Z_p$ as well.
In each row of (C37.2), $p-1$ divides $m$.
Reduction modulo $p$ makes the factor in (C37.7)
equal to one.
The reductions thus give the required number of
independent invariant ordinary homomorphisms
$U\to\mathbb F_p$.

The rest of the full-group argument is exactly the
already proved one in C.33: the continuous five-term
sequence and
$H^2_{\mathrm{cts}}(\mathbb Z_p^\times,\mathbb F_p)=0$
extend each character to $G$.
Its arithmetic restriction is unramified away $p$;
Kronecker--Weber removes it by a scalar multiple of
$c_p(\lambda)$ from (C34.7).
The adjustment vanishes on global multiplier one.
All elements of the actual local multiplier-one
image have full multiplier-one lifts by P.4, so
the independent values of the reductions remain
actual full-group witnesses.
This proves (C37.1) and the last assertion.
\end{proof}


% Root and DE independent full proof audits passed on 2026-11-15.
\begin{statement}{C.38}{Proved}{Unbounded prime-power arithmetic-trivial characters}
Let \(Q\) be the full gentle group modulo the closed normal
subgroup generated by its absolute Galois image.
For every prime \(p\geq5\), every integer \(k\geq1\), and every
even integer \(m\geq12\) divisible by \((p-1)p^{k-1}\), there is
a continuous surjection
\[
                  Q^{\mathrm{ab}}\longrightarrow
                       (\mathbb Z/p^k\mathbb Z)^{d_m},
                                                               \tag{C38.1}
\]
where \(d_m=a_m-b_m>0\) is the dimension lower bound of
C.34--C.36.  The independent values are witnessed on global
multiplier one.  The integers \(d_m\) tend to infinity along
the indicated progression.

Consequently, for every \(p\geq5\),
\[
 \dim_{\mathbb F_p}\operatorname{Hom}_{\mathrm{cts}}
                         (Q,\mathbb F_p)=\aleph_0,
                                                               \tag{C38.2}
\]
and there is a simultaneous continuous surjection
\[
 Q^{\mathrm{ab}}\longrightarrow
       \left(\prod_{j\geq1}C_2\right)
       \times
       \prod_{\substack{p\geq5\\p\ {\rm prime}}}
             \prod_{k,j\geq1}\mathbb Z/p^k\mathbb Z .
                                                               \tag{C38.3}
\]
In particular its \(p\)-primary quotients have unbounded
exponent for every \(p\geq5\).  Neither a
\(\mathbb Z_p\) quotient nor the complete abelianization is
asserted.  The general characters in this statement are not
asserted to vanish on the pentagon group.
\end{statement}
\begin{proof}
Fix \(p\geq5\) and a finite top degree \(m\).
Write \(\mathcal G\) for the actual local gentle image,
\(\mathcal U\) for its local multiplier kernel, and
\(\mathcal U_n\) for the normalized conjugator filtration.
The complete integral lifting argument of C.33 applies
in every finite degree, and gives
\[
       \mathcal U_n/\mathcal U_{n+1}=S_n(\mathbb Z_p).
                                                               \tag{C38.4}
\]
These are the actual full-group images, not just solutions
of the equations modulo a nilpotent term.
Set
\[
       U=\mathcal U/\mathcal U_{m+1},\qquad
       G=\mathcal G/\mathcal U_{m+1}.
\]
The action on the integral Magnus algebra modulo
augmentation degree \(m+2\) realizes \(U\) faithfully as
a compact group of unipotent matrices over \(\mathbb Z_p\).
For \(u\in U\), the operator \(u-1\) raises degree by at
least three.  This follows from \(S_2=0\), as in C.33.
No bound on the nilpotency class relative to \(p\) is imposed.

\paragraph{The actual rational Lie algebra.}
The finite matrix logarithm and exponential are inverse
polynomials over \(\mathbb Q_p\).  Define
\[
             L=\operatorname{span}_{\mathbb Q_p}
                               \{\log u:u\in U\}.
                                                               \tag{C38.5}
\]
This is a Lie subalgebra of the finite-dimensional nilpotent
matrix algebra.  Here is a direct verification that avoids
an integral correspondence in class at least \(p\).
For \(X=\log u\) and \(Y=\log v\), the finite polynomial
\(\log(u^a v^b)=\operatorname{BCH}(aX,bY)\) belongs to \(L\)
whenever \(a,b\) are nonnegative integers.  Evaluating on a
sufficiently large finite integer grid and inverting the
two rational Vandermonde matrices shows that every
coefficient of this polynomial belongs to \(L\).
Its coefficient of \(ab\) is \([X,Y]/2\).
Thus \([X,Y]\in L\); bilinearity proves the assertion.

The group \(U\) contains a neighborhood of the identity
in \(\exp L\).  Choose a basis \(X_1,\ldots,X_d\) of \(L\)
consisting of logarithms of actual elements
\(u_1,\ldots,u_d\in U\).
For every \(t_i\in\mathbb Z_p\), the continuous power
\(u_i^{t_i}\) is in \(U\), and equals \(\exp(t_iX_i)\):
the equality holds on the dense subset of integers and
both sides are continuous.
In these basis coordinates the polynomial map
\[
 (t_1,\ldots,t_d)\longmapsto
       \log\bigl(u_1^{t_1}\cdots u_d^{t_d}\bigr)
                                                               \tag{C38.6}
\]
has the form \(t+R(t)\), where \(R\) has no term of
degree zero or one.
For sufficiently large \(r\), the map \(R\) sends
\(p^r\mathbb Z_p^d\) into \(p^{r+1}\mathbb Z_p^d\) and is
Lipschitz there with constant at most \(p^{-1}\).
For each \(y\) in that ball, the contraction
\(t\mapsto y-R(t)\) therefore gives a solution of
\(t+R(t)=y\).
It follows that
\[
       \exp\left(p^r\sum_i\mathbb Z_pX_i\right)\subseteq U .
                                                               \tag{C38.7}
\]
This is an explicit rational argument; no claim that
\(\log U\) itself is an integral Lie lattice is needed.

Let \(L_n\) be the intersection of \(L\) with the
operators that raise augmentation degree by at least \(n\).
This is the same filtration as the normalized conjugator
filtration: a first degree-\(n\) conjugator term \(f_n\)
changes \(Y\) first by \([Y,f_n]\), and the centralizer
of \(Y\) in the characteristic-zero free Lie algebra
has no homogeneous part in degree \(n\geq2\).
The leading terms of logarithm and exponential agree
with their linear terms.
Equation (C38.7) shows that every vector of \(L_n\),
after multiplying by a sufficiently large power of \(p\),
is the logarithm of an element of \(U_n\).
Conversely logarithms of \(U_n\) lie in \(L_n\).
Taking leading terms and using (C38.4) therefore gives
\[
           L_n/L_{n+1}=S_n(\mathbb Q_p).
                                                               \tag{C38.8}
\]
The induced bracket is the actual leading Ihara bracket.
This verifies the rational Lie filtration directly from
the completed group image.

\paragraph{A grading supplied by an actual wild multiplier.}
Choose an element \(g_0\) in the actual local Galois image
whose multiplier is \(1+p\), and put
\(A=\operatorname{Ad}(g_0)|_L\).
It preserves the filtration and acts on its degree-\(n\)
quotient by the scalar \((1+p)^n\).
These scalars, for \(3\leq n\leq m\), are pairwise
distinct in \(\mathbb Q_p\).
The filtration therefore implies
\[
           \prod_{n=3}^{m}\bigl(A-(1+p)^n\bigr)=0.
                                                               \tag{C38.9}
\]
For example, apply the factors in ascending filtration
order; each annihilates the current successive quotient
and moves the image to the next term.
The polynomial has no repeated root, so \(A\) is
semisimple and yields a rational splitting
\[
      L=\bigoplus_{n=3}^{m}L^{(n)},\qquad
      L^{(n)}\xrightarrow{\ \sim\ }S_n(\mathbb Q_p).
                                                               \tag{C38.10}
\]
Since \(A\) is a Lie automorphism, the bracket of the
\(i\) and \(j\) eigenspaces lies in the \(i+j\)
eigenspace when \(i+j\leq m\), and is zero otherwise.
Its leading term identifies it with the Ihara bracket
under (C38.10).
Thus the entire rational Lie algebra has the gentle
grading.  This removes hidden higher filtered bracket
relations even when many tame weights coincide and the
nilpotency class exceeds \(p\).

The weight-\(m\) part of \(L/[L,L]\) has dimension
\[
        h_m\geq a_m-b_m=d_m.
                                                               \tag{C38.11}
\]
Indeed the number of independent pair brackets in that
degree is at most \(b_m\), exactly as in C.34.
Choose \(h_m\) independent rational linear functionals
on this part and extend them by zero on the other
eigenspaces.  Together they give a surjection
\(\ell:L\to\mathbb Q_p^{h_m}\) annihilating \([L,L]\).
Consequently
\[
                      u\longmapsto\ell(\log u)
                                                               \tag{C38.12}
\]
is an ordinary continuous group homomorphism:
every non-linear term of the finite BCH polynomial
is a Lie commutator and is annihilated.

Conjugation on \(L/[L,L]\) factors through
\(G/U=\mathbb Z_p^\times\).
Indeed inner conjugation by \(u=\exp X\) induces
\(\exp(\operatorname{ad}X)\), which is the identity
on the Lie abelianization.
These conjugation operators commute with \(A\) and
therefore preserve its distinct weight spaces.
On the top weight \(m\), they act exactly as on
the top filtration quotient, namely by
\(\lambda_p(g)^m\).
Thus
\[
       \ell(\operatorname{Ad}(g)X)
                    =\lambda_p(g)^m\ell(X)
             \quad(g\in G,\ X\in L).
                                                               \tag{C38.13}
\]

\paragraph{Integral values and finite prime-power characters.}
The image of (C38.12) is a compact additive subgroup
\(J\subset\mathbb Q_p^{h_m}\).
It is a \(\mathbb Z_p\)-module: it is closed and stable
under integer multiplication.
It spans \(\mathbb Q_p^{h_m}\), since \(\log U\)
spans \(L\) and \(\ell\) is onto.
It is consequently a full free \(\mathbb Z_p\)-lattice.
Choose a basis of \(J\).  In this basis (C38.12) becomes
a continuous surjection
\[
                \Phi:U\longrightarrow\mathbb Z_p^{h_m},
       \qquad
       \Phi(gug^{-1})=\lambda_p(g)^m\Phi(u).
                                                               \tag{C38.14}
\]
This passage uses the actual compact image, so no
unsupported integral BCH assertion or reduction of a
possibly nonprimitive rational functional is involved.

Suppose \((p-1)p^{k-1}\mid m\).
For every \(\lambda\in\mathbb Z_p^\times\),
\(\lambda^m\equiv1\pmod{p^k}\).
Reduction of (C38.14) is therefore an invariant
surjection \(U\to(\mathbb Z/p^k\mathbb Z)^{h_m}\).
The continuous five-term sequence extends each coordinate
to \(G\), because
\[
 H^2_{\mathrm{cts}}(\mathbb Z_p^\times,\mathbb Z/p^k\mathbb Z)=0.
                                                               \tag{C38.15}
\]
As before, this follows from averaging on \(\mu_{p-1}\)
and the \(p\)-cohomological dimension one of
\(1+p\mathbb Z_p\).

The restriction of each extended character to
\(G_{\mathbb Q}\) is unramified away \(p\), by the
actual pro-\(p\) tangential action cited in C.33.
Kronecker--Weber identifies every such character with
a scalar multiple of the cyclotomic character
\[
 c_{p,k}(\lambda)=
        \frac{\log(\lambda_p^{p-1})}{p}\pmod{p^k}.
                                                               \tag{C38.16}
\]
Here the logarithm is taken on \(1+p\mathbb Z_p\).
It is an ordinary additive character and is onto:
the map \(\log/p:1+p\mathbb Z_p\to\mathbb Z_p\)
is an isomorphism, while \(p-1\) is a unit.
Subtracting these scalar multiples removes the
arithmetic restrictions simultaneously.
The correction is zero on global multiplier one.
Every element of the actual local multiplier-one image
has a full global multiplier-one lift by P.4.
Thus the joint surjectivity in (C38.14), reduced
modulo \(p^k\), remains valid for the corrected
full-group characters.
This proves (C38.1), since \(h_m\geq d_m\).
C.34 gives \(d_m\to\infty\) along the progression,
and C.36 gives positivity for every even \(m\geq12\).

\paragraph{Choosing all finite orders simultaneously.}
Let \(K\) be the global multiplier-one subgroup.
For fixed \(p,k\), let \(V_{p,k}\) be the subspace of
\(\operatorname{Hom}_{\mathrm{cts}}(K,\mathbb F_p)\)
consisting of reductions of the restrictions of
arithmetic-trivial full-group characters to
\(\mathbb Z/p^k\mathbb Z\).
The joint surjections just proved show that
\(\dim_{\mathbb F_p}V_{p,k}=\infty\).
Enumerate all pairs \((k,j)\) of positive integers.
At each step choose an actual such character whose
mod-\(p\) restriction lies outside the span of all
previously chosen restrictions.
This is possible because the relevant \(V_{p,k}\)
is infinite-dimensional and only finitely many
characters have been chosen.

For any finite collection of the resulting characters,
the image \(H\) in its product target \(T\) satisfies
\(H+pT=T\), by independence of their mod-\(p\)
restrictions to \(K\).
The finite \(\mathbb Z_p\)-module \(T/H\) therefore
equals \(p(T/H)\), so it is zero.
The finite product map is onto already on \(K\).
Compactness now gives a surjection onto
\(\prod_{k,j\geq1}\mathbb Z/p^k\mathbb Z\).
Combine these maps for distinct primes and C.30's
independent quadratic family: every finite target
product has pairwise coprime primary factors and is
onto.  Another use of compactness proves (C38.3).

Finally the full gentle group is a closed subgroup
of the automorphism group of a finitely generated
free profinite group, hence is countably based.
It has only countably many clopen subsets, and thus
only countably many continuous maps to any fixed
finite group.  This gives the upper bound
\(\aleph_0\) in (C38.2); the preceding construction
gives the matching lower bound.
The product of finite cyclic groups in (C38.3)
is not being replaced by a \(\mathbb Z_p\) factor:
unbounded finite orders alone do not provide a
compatible infinite tower of characters.
\end{proof}


% Root and DE independent full proof audits passed on 2026-11-15.
\begin{statement}{C.39}{Proved}{Pentagon-trivial characters at five, seven and thirteen}
For each \(p\in\{5,7,13\}\) there is a continuous surjective
homomorphism
\[
              \epsilon_p:\widehat{\mathrm{GT}}_0
                                      \longrightarrow\mathbb F_p
\]
which vanishes on the full pentagon group and is nonzero
on global multiplier one.
In particular, if \(Q_{\rm GT}\) is the quotient of the full
gentle group by the closed normal subgroup generated by the
full pentagon group, then
\[
          Q_{\rm GT}^{\rm ab}\longrightarrow
             C_5\times C_7\times C_{13}
                       \times\prod_{j\geq1}C_2
                                                               \tag{C39.1}
\]
is a continuous surjection.
This gives characters of the gentle group that vanish on
the pentagon group, not nonzero characters of the pentagon
group itself.
\end{statement}
\begin{proof}
Use C.38 at top conjugator degree \(m=12\), and write
\(L=\bigoplus_{n=3}^{12}L^{(n)}\) for the rational grading
of the actual local multiplier-one image \(U\).
The exact calculation of C.33 shows
\[
       \dim S_{12}=7,\qquad
       \dim[S,S]_{12}=6.
\]
Its nonzero rational annihilator is generated by the
functional \(\ell=(-4,1,8,-4,-17,14,84)\) in the
retained basis.
Extend this functional to \(L\) by zero in every degree
other than twelve.
The argument of C.38 normalizes the actual compact image
of \(u\mapsto\ell(\log u)\) and gives an onto homomorphism
\[
       \Phi:U\longrightarrow\mathbb Z_p,\qquad
       \Phi(gug^{-1})=\lambda_p(g)^{12}\Phi(u).
                                                               \tag{C39.2}
\]
The scalar used to normalize this image is nonzero and
will not affect any vanishing assertion.

We must show vanishing on the actual pentagon image, not
only on homogeneous solutions of its first equation.
Let \(P\) be the image of the full pentagon group in
\(G=\mathcal G/\mathcal U_{13}\), and put \(P_1=P\cap U\).
It is a closed compact subgroup of the actual matrix group.
Define
\[
          L_P=\operatorname{span}_{\mathbb Q_p}
                                      \{\log h:h\in P_1\}.
\]
The finite BCH/Vandermonde proof in C.38 shows that \(L_P\)
is a Lie subalgebra.  The same explicit contraction proof,
using a basis of logarithms of elements of \(P_1\), shows
that \(P_1\) contains \(\exp\) of a sufficiently small full
lattice in \(L_P\).

Choose the actual Galois element \(g_0\) of local multiplier
\(1+p\) used in C.38.
It lies in \(P\), so \(L_P\) is stable under
\(A=\operatorname{Ad}(g_0)\).
The rational polynomial projector \(e_{12}\) for the
eigenvalue \((1+p)^{12}\) therefore preserves \(L_P\).
For any \(h\in P_1\), put
\[
                         X=e_{12}\log h .
\]
It belongs to \(L_P\cap L_{12}\).
For sufficiently large \(r\), the open exponential
neighborhood just proved gives
\(\exp(p^rX)\in P_1\).
This is an actual top-filtered element of the pentagon
image; all its lower conjugator terms vanish.
Its first pentagon equation is consequently the linear
weight-twelve pentagon equation, with homogeneous
conjugator \(p^rX\) under the identification with \(S_{12}\).

The exact rational row identity proved and executed in
C.35 gives
\[
 698922\,\ell
   =-192329R_1+472R_2-26842R_3-107926R_4-11870R_5.
                                                               \tag{C39.3}
\]
Every \(R_i\) vanishes on that actual leading pentagon
term.  Hence \(\ell(p^rX)=0\), and therefore
\(\ell(X)=0\) in \(\mathbb Q_p\).
This deduction is valid also when \(p\) divides the
integer in (C39.3): the identity is applied over
\(\mathbb Q_p\), before normalizing the actual compact
image or reducing it.
Since \(\ell\) is zero on all other eigenspaces,
\(\ell(\log h)=\ell(X)=0\).
Thus
\[
                              \Phi(P_1)=0.                     \tag{C39.4}
\]

For \(p=5,7,13\), the integer \(p-1\) divides twelve.
Reduction of (C39.2) modulo \(p\) is therefore invariant
under all of \(G\).
Extend it to a character of \(G\) using the vanishing
of \(H^2(\mathbb Z_p^\times,\mathbb F_p)\), and subtract
its cyclotomic Galois restriction as in C.38.
The resulting full gentle character is onto, since
the correction is zero on global multiplier one
and all local multiplier-one values have actual
complete full-group lifts.

It vanishes on \(P\).
Indeed for each \(g\in P\), choose a Galois element
\(\sigma\in P\) with the same local multiplier.
Then \(g\sigma^{-1}\in P_1\).
The corrected character is zero on \(\sigma\)
by construction and on \(P_1\) by (C39.4), so it
is zero on \(g\).
This proves the stated assertions for all three primes.
Their orders are pairwise coprime, and are also coprime
to the quadratic family of C.30.
Every finite product is consequently onto; compactness
gives (C39.1).
\end{proof}


% C.40: root and DE independent full proof/code audits passed.
\subsection{Independent local coordinates before and after descent}

\begin{statement}{C.40}{Proved}{The simultaneous gentle image at every prime at least five}
For a prime $p\geq5$, let $G_{0,p}$ denote the group of actual
normalized pro-$p$ gentle pairs, with conjugator in the pro-$p$
commutator subgroup, equivalently their faithful pure outer classes:
the two cycle identities hold and the induced
map of the free pro-$p$ group is an automorphism. Let
$\lambda_p:G_{0,p}\to\mathbb Z_p^\times$ be its multiplier.
Then the image of the actual full group under
\[
 \widehat{\mathrm{GT}}_0\longrightarrow
 \prod_{p\geq5}(G_{0,p}\times G_{0,p}),\qquad
 g\longmapsto (g_p,D(g)_p)_{p\geq5}
                                                        \tag{C.40.1}
\]
is exactly
\[
 \prod_{p\geq5}
       (G_{0,p}\mathbin\times_{\mathbb Z_p^\times}G_{0,p}).
                                                        \tag{C.40.2}
\]
In particular, the multiplier-one image is the whole product
of the two multiplier-one local gentle groups, simultaneously
at all these primes. This determines the gentle-group image
asked for in Task C.5. It does not determine the simultaneous
image of the full pentagon group, or assert any inclusion of
that group in $C^\varphi$.
\end{statement}

\begin{proof}
We first prescribe multiplier-one elements at a finite set of
primes $p\geq5$. The construction retains both local actions
in one actual quotient of the modular source, then checks its
entire source-centralization. No assertion that a relative
quotient is an absolute pronilpotent quotient is used.

\paragraph{The six actual fillings and their finite base.}
Put
\[
 \begin{gathered}
 \Gamma=\widehat{C_2*C_3}=\langle s,r\mid s^2=r^3=1\rangle,
 \qquad F=\ker(\Gamma\to S_3),\\
 x=(rs)^2,\qquad y=(sr)^2,\qquad z=(xy)^{-1}.
 \end{gathered}
                                                        \tag{C.40.3}
\]
Write $\sigma=\operatorname{Int}(s)$ and
$\rho=\operatorname{Int}(r)$. Direct reduction in the free
product gives the exact formulas
\[
 \sigma(x)=y,\quad\sigma(y)=x,\qquad
 \rho(x)=z,\quad\rho(y)=x.                              \tag{C.40.4}
\]
Thus $\sigma^2=\rho^3=1$ as actual automorphisms of $F$.
Let $\Gamma'$ be a separately marked copy of $\Gamma$, and
let the filling $\kappa:F\twoheadrightarrow\Gamma'$ be
\[
       \kappa(x)=s',\qquad \kappa(y)=s'r',\qquad
                         \kappa(z)=r'^{-1}.
                                                        \tag{C.40.5}
\]
For $k=0,1,2$ use the six maps
$\kappa_k=\kappa\rho^k$ and
$\kappa'_k=\kappa\sigma\rho^k$. These are actual maps,
not a list of unrelated outer isomorphism types. The order-two
automorphism
\[
                     t_k=\rho^{-k}\sigma\rho^k
                                                        \tag{C.40.6}
\]
interchanges $\kappa_k$ and $\kappa'_k$ literally.

In the standard two-dimensional representation of the target
$S_3$, choose
\[
 S=\begin{pmatrix}0&1\\1&0\end{pmatrix},\quad
 R=\begin{pmatrix}-1&-1\\1&0\end{pmatrix},\quad
 E=SR,\quad T=RS.
                                                        \tag{C.40.7}
\]
Here $T^2=1$, $TST=E$, and $TET=S$. Thus the finite
maps belonging to a sibling pair have the same kernel and
are related by $\operatorname{Int}(T)$. The three distinct
finite maps give the image $Q_0\subset S_3^3$ generated by
\[
             X=(S,R^{-1},E),\qquad Y=(E,S,R^{-1}).
                                                        \tag{C.40.8}
\]
This group is
\[
                  Q_0=C_3^3\rtimes C_2^2,
                       \qquad |Q_0|=108.              \tag{C.40.9}
\]
For example, the two sign vectors are $(1,0,1)$ and
$(1,1,0)$. Their commutator has a nonzero entry in each
of the three $C_3$ coordinates. The three distinct
nontrivial characters of $C_2^2$ on these coordinates
separate them over $\mathbb F_3$, so conjugates of that
commutator span $C_3^3$. This proves (C.40.9).

Let $K_0=\ker(F\to Q_0)$. The six kernels form the
orbit under conjugation by $\Gamma$, hence $K_0$ is
normal in $\Gamma$. Consequently
\[
       \widetilde Q_0=\Gamma/K_0,
                 \qquad |\widetilde Q_0|=648.          \tag{C.40.10}
\]
Only this finite extension of $Q_0$ by the old $S_3$
is asserted. No splitting inside the original modular
group is assumed. For an exact finite verification,
\texttt{pbc\_joint\_local\_descent.py} enumerates the
$108$ triples, verifies every generator edge for the
actions $X\leftrightarrow Y$ and
$(X,Y)\mapsto((XY)^{-1},X)$, and obtains a group of
$648$ actual permutations. Its inner subgroup has order
$108$ and $Q_0$ has trivial centre. Thus this also gives
a faithful explicit model of the finite extension.

\paragraph{The joint local quotient has seven independent free factors.}
Fix $p\geq5$ and put
\[
 E_p=\Gamma'/\ker(K'_2\to K_2'{}^{(p)}),\qquad
 P_p=K_2'{}^{(p)}.
                                                        \tag{C.40.11}
\]
The group $P_p$ is free pro-$p$ of rank two and is the
normal kernel of $E_p\to S_3$. It has a complement
because six is prime to $p$.

For one sibling pair, the two original filling kernels
together impose $x^2=y^2=z^3=1$. This is the finite
triangle group of type $(2,2,3)$, namely $S_3$.
The relative pro-$p$ kernels contain these original
kernels and are contained in the common finite
$S_3$ kernel. Their product is therefore precisely
that finite kernel. Goursat's lemma gives the whole
fibre product $E_p\times_{S_3}E_p$, with the finite
identification specified by $T$ in (C.40.7).
The image of $K_0$ in its pro-$p$ kernel $P_p^2$ is
also onto: the quotient by that image is a quotient
of $F/K_0$, so it is both pro-$p$ and of order prime
to $p$, and hence is trivial.

Let $J_p$ be the image of all six fillings in $E_p^6$.
Its finite image is the actual $Q_0$ just specified.
The image of $K_0$ is a closed normal pro-$p$ subgroup
of $P_p^6$. On the Frattini quotient of this product,
the three sibling blocks are three pairwise
nonisomorphic absolutely irreducible standard
$\mathbb F_p[Q_0]$ modules, each with multiplicity two.
They are distinct because the three $C_3$ coordinates
act on different blocks. Maschke's theorem and the
surjectivity to each whole sibling pair show that the
Frattini image is the entire direct sum. The Frattini
criterion for a finitely generated pro-$p$ group now
gives
\[
                1\longrightarrow P_p^6
                  \longrightarrow J_p
                  \longrightarrow Q_0\longrightarrow1.
                                                        \tag{C.40.12}
\]
Each standard module has zero $Q_0$ coinvariants.
The prime-to-$p$ quotient therefore gives
$H_1(J_p,\mathbb F_p)=0$. In particular $J_p$ has no
nontrivial pro-$p$ quotient: any such quotient would
have a nontrivial finite $p$-group quotient and then
a quotient of order $p$.

The joint image of $F$ in $F^{(p)}\times J_p$ is thus
the full product, again by Goursat. The kernels of
the six fillings form a deck orbit, and the ordinary
pro-$p$ kernel is characteristic. Their intersection
is normal in $\Gamma$. The resulting actual quotient
$Q_p$ of $\Gamma$ has the exact sequence
\[
 1\longrightarrow\Pi_p\longrightarrow Q_p
        \longrightarrow\widetilde Q_0\longrightarrow1,
 \qquad
 \Pi_p=P_{0,p}\times
       \prod_{k=0}^2(P_{k,p}\times P'_{k,p}),            \tag{C.40.13}
\]
where all seven factors are free pro-$p$ groups of
rank two. Here $P_{0,p}$ is the ordinary quotient of
the old $F$. The restriction to $F$ retains its map
to that factor and all six original fillings.

For finitely many distinct primes, the joint image
of $K_0$ in these kernels is the product of the
$\Pi_p$: finite quotients of different prime orders
have no common quotient, and compactness passes from
finite quotients to the pro-prime groups. We therefore
obtain one actual quotient $Q$ with finite base
$\widetilde Q_0$ and normal kernel
\[
                         \Pi=\prod_{p\in\mathcal S}\Pi_p.
                                                        \tag{C.40.14}
\]

\paragraph{The central homology vanishes at the retained primes.}
We compute the continuous homology with compact
$\mathbb Z_p$ coefficients. Each free pro-$p$ factor
has homological dimension one. The completed Kunneth
formula gives
\[
 H_2(\Pi_p,\mathbb Z_p)=
   \bigoplus_{i<j}
       H_1(P_{i,p},\mathbb Z_p)\otimes_{\mathbb Z_p}
       H_1(P_{j,p},\mathbb Z_p).                        \tag{C.40.15}
\]
The tensor factors are finite free lattices; their
ordinary and completed tensor products agree. There
are no single-factor $H_2$ terms and no Tor terms.
Finite-group averaging applies because $648$ is a
unit in $\mathbb Z_p$.

The old $Q_0$ acts trivially on $H_1(P_{0,p})$.
Its three other standard isotypic types are distinct.
Consequently taking $Q_0$ invariants in (C.40.15)
leaves exactly the three sibling invariant lines.
Here is the integral calculation of their actual
deck action. In the bases (C.40.7), put
\[
 B=\begin{pmatrix}2&-1\\-1&2\end{pmatrix},\qquad
 C=BT^{\mathsf t}=
               \begin{pmatrix}-1&-1\\-1&2\end{pmatrix}.
                                                        \tag{C.40.16}
\]
For $A=S,R$ one has $ABA^{\mathsf t}=B$.
The two actions in the sibling pair are $A$ and
$TAT^{-1}$, so its invariant tensor is $C$. It is
primitive and symmetric. These matrices describe
the actual standard representation: in the original
column basis, $\rho$ acts by
$\left(\begin{smallmatrix}-1&1\\-1&0\end{smallmatrix}\right)$;
the basis matrix
$\left(\begin{smallmatrix}1&2\\2&1\end{smallmatrix}\right)$,
of determinant $-3$, changes it to $R$ and leaves $S$
unchanged. It is invertible over every $\mathbb Z_p$
under consideration.

The actual involution $t_k$ in (C.40.6) interchanges
the two sibling coordinates without an additional
linear map. Interchanging two degree-one homology
factors acts on their cross term as minus transpose.
It therefore acts by $-1$ on that invariant line.
Since this pair is fixed as a set, no other line is
sent into it; an invariant vector must have zero
coefficient on that line. Applying this for each
$k$ kills all three lines. Thus the full
$\widetilde Q_0$ invariants, and equivalently its
coinvariants, vanish. Hochschild--Serre gives
\[
                H_2(Q_p,\mathbb Z_p)=0,
                    \qquad H_1(Q_p,\mathbb Z_p)=0.      \tag{C.40.17}
\]
The latter assertion also follows directly from the
standard representations, including the ordinary
factor on which the old $S_3$ acts in its standard
representation.

At a coefficient prime $q\ne p$, a pro-$p$ group
has zero positive compact $\mathbb Z_q$ homology:
the averaging maps on its finite quotients use only
orders invertible in $\mathbb Z_q$, and are compatible
with passage to the limit. For (C.40.14) the same
argument removes every prime except the coefficient
prime. Therefore, for every prime $q$ and $i=1,2$,
\[
 H_i(Q,\mathbb Z_q)\xrightarrow{\sim}
                 H_i(\widetilde Q_0,\mathbb Z_q).       \tag{C.40.18}
\]
At a retained prime both sides are zero; at every
other prime the normal kernel is acyclic.

\paragraph{The entire source-centralization is a pullback.}
Let $N=\ker(\Gamma\to Q)$ and form the two actual
centralizations
\[
 H=\Gamma/\overline{[\Gamma,N]},\quad K=N/\overline{[\Gamma,N]},
 \qquad
 H^0=\Gamma/\overline{[\Gamma,K_0]},\quad
 K^0=K_0/\overline{[\Gamma,K_0]}.
                                                        \tag{C.40.19}
\]
The compact abelian kernel has its usual prime
components $K_q$. As in B.27, the finite cyclic
periodic resolution and the free-product homology
sequence give $H_2(\Gamma,\mathbb Z_q)=0$ for every
$q$. The actual five-term sequence is therefore
\[
 0\longrightarrow H_2(Q,\mathbb Z_q)
  \longrightarrow K_q\longrightarrow H_1(\Gamma,\mathbb Z_q)
  \longrightarrow H_1(Q,\mathbb Z_q)\longrightarrow0.
                                                        \tag{C.40.20}
\]
Use the same sequence for $\widetilde Q_0$ and $K^0$.
Naturality and (C.40.18) show that $K_q\to K_q^0$ is
an isomorphism for every $q$. This is deliberately
the full five-term argument: $N\subset\Gamma'$ is
not assumed, and no identification with a universal
central extension of an arbitrary quotient is used.
Comparing kernels and quotients gives the actual
Cartesian square
\[
                    H\simeq Q\mathbin\times_{\widetilde Q_0}H^0.
                                                        \tag{C.40.21}
\]
In fact $K^0$ is finite and has only two- and
three-primary parts, by (C.40.20) for the finite
$\widetilde Q_0$ and $\Gamma^{\rm ab}=C_6$.

It follows that every automorphism $a$ of $Q$ that
is identity on $\widetilde Q_0$ has a lift
$(a,\operatorname{id})$ to $H$. If its three positive
peripheral conjugators lie in $\Pi$, then their pairs
$(b_s,1),(b_r,1),(b_\tau,1)$ are actual elements of
this fibre product and are peripheral conjugators
for that lifted automorphism. This verifies the
central refinement, including the peripheral data.

We recall precisely why this yields a full lift.
For any finite source quotient $\Gamma\twoheadrightarrow E$,
let $A$ be the image of $N$ and $D=E/A$. In the
finite scalar-cover construction of P.3, every
auxiliary source map to a central extension of $D$
sends $N$ centrally. It thus factors through $H$.
Precomposing by the actual peripheral automorphism
of $H$ supplies the new epimorphism and the three
specified conjugators in every auxiliary cover.
The trace count and subgroup-by-subgroup Mobius
inversion in P.3 apply because they require only
normality of $A$. They give an epimorphism to $E$
with the specified retained values and frames.
Adjoin any finite retained quotient of $Q$ to $E$
when necessary. Compactness of the three conjugator
variables, imposing generation in every finite
quotient, then gives a surjective continuous
endomorphism of the actual modular source. Its
finite generation makes it an automorphism. This
is the modular version of B.11 with arbitrary
specified $N$, and proves a full lift retaining
the action on $Q$ and the three frame values.

\paragraph{Constructing arbitrary prescribed two local actions.}
Now prescribe $\alpha_p,\beta_p\in G_{0,p}$ of
multiplier one for every $p\in\mathcal S$.
An actual local gentle class extends to the relative
modular group $E_p$ with identity quotient action on
$S_3$, by the two cycle formulas of P.4, or by the
normalizer proof of C.1. It preserves its three
positive orbifold peripheral classes. After an
inner correction in $P_p$, it fixes a chosen finite
$S_3$ complement pointwise: two such complements
are conjugate by the normal pro-$p$ group. Its
restriction to $P_p$ therefore commutes with that
complement, and represents the same local gentle
outer class.

Choose a complement $\widetilde Q_0$ in $Q$ to
$\Pi$. Existence and conjugacy follow from the
coprime complement theorem, first at finite
quotients and then by compactness; all primes of
the kernel are at least five. On $P_{0,p}$ its
action factors through the old $S_3$: the kernel
$Q_0$ would otherwise act through a finite
prime-to-$p$ subgroup of the inner pro-$p$
automorphism group, which must be trivial.
Choose the representative of $\alpha_p$ commuting
with this complement action.

The six other factors form a transitive orbit of
$\widetilde Q_0$, with stabilizer $Q_0$. On a
chosen factor its stabilizer action factors through
its actual finite target $S_3$. Indeed the image
of this finite stabilizer in its relative target
$E_p$ is a finite complement; its kernel in $P_p$
is trivial. Choose a representative of $\beta_p$
commuting with that action, and transport it to
the other five factors using the actual complement
action. Stabilizer commutation makes the transport
independent of the chosen transporting element.
All transported maps preserve the corresponding
positive cusp classes; the transport comes from
the six actual marked fillings, up to their inner
target changes.

The product of these seven automorphisms, for all
$p\in\mathcal S$, commutes with $\widetilde Q_0$.
Extending it by identity on this complement gives
an actual automorphism $a$ of $Q$. Its ordinary
local outer coordinate is $\alpha_p$, and its
selected filled coordinate is $\beta_p$.

We check all three required modular conjugators
in the kernel $\Pi$. The elements $a(s)$ and $s$
have order two above the same finite base element,
and $a(r)$ and $r$ have order three. Coprime
complement conjugacy supplies $\Pi$-conjugators.
For the third element put $\tau=sr$, so
$\tau^2=y$. The ordinary coordinate of $a(y)$ is
conjugate to $y$ by a pro-$p$ element because
$\alpha_p$ is peripheral. In each filled coordinate,
the image of $y$ is one of the three orbifold
peripheral elements, and the corresponding extended
$\beta_p$ gives a pro-$p$ conjugator. These independent
coordinate conjugators assemble to an element of
$\Pi$. Thus $a(\tau^2)$ is $\Pi$-conjugate to
$\tau^2$.

After this inner adjustment, suppose
$a(\tau^2)=\tau^2$. Write $a(\tau)=v\tau$.
Then $v\in C_\Pi(\tau^2)$, and
\[
                  v\operatorname{Int}(\tau)(v)=1.
                                                        \tag{C.40.22}
\]
The automorphism $\operatorname{Int}(\tau)$ has
order dividing two on this centralizer, since
its square is conjugation by $\tau^2$.
Coprime complement conjugacy gives
$H^1(C_2,C_\Pi(\tau^2))=1$; it can equivalently
be applied in each pro-$p$ factor and assembled.
Equation (C.40.22) is therefore a coboundary,
and an additional $\Pi$-conjugator carries
$a(\tau)$ to $\tau$. This proves positive peripheral
conjugacy of the third element with a pure retained
frame. The same argument, with one target factor,
justifies the pro-$p$ frame for its infinite
orbifold cusp used in the preceding paragraph.

By (C.40.21), these three frames lift as
$(b_i,1)$ to the entire source-centralization.
The finite refinement and compactness argument
therefore gives an actual positive peripheral
automorphism $\Phi$ of the full $\Gamma$, retaining
$a$ and those frames. Its quotient action on the
old $S_3$ is identity. Restricting to $F$ gives
a pure peripheral outer class centralizing the
geometric $S_3$, hence a full gentle element of
multiplier one. Equivalently, fixing the modular
cusp by an inner correction gives the exact full
two- and three-cycle equations as in P.4; standard
normalization gives the conjugator in $F'$.
These are full equations, not only equations in
the retained local quotients.

The ordinary retained quotient is $F^{(p)}$, so
the first local outer class is exactly $\alpha_p$.
The selected filling is the original $\kappa$;
its relative target action restricts to
$\beta_p$ on its characteristic pure subgroup.
This is precisely the operation $D$ of C.1.
Inner normalization of the full representative
does not change either outer coordinate. We have
therefore realized every prescribed finite tuple
$(\alpha_p,\beta_p)_{p\in\mathcal S}$.

\paragraph{All primes and all common multipliers.}
The full gentle group is compact, and the local
coordinate maps and $D$ are continuous. Surjectivity
for every finite set of primes therefore gives
surjectivity onto the whole multiplier-one product.
For a tuple of arbitrary common multipliers
$(\lambda_p)_{p\geq5}$, choose a global unit with
those components and an actual Galois element
having that cyclotomic character. The cyclotomic
character is onto. Its two coordinates agree
because $D$ fixes the specified Galois image.
Divide each desired pair by these common Galois
coordinates, apply the multiplier-one result,
then multiply by the same Galois element.
This proves (C.40.2). The reverse inclusion is
the multiplier preservation already proved in C.1.
\end{proof}


% C.41: root and FGH independent full proof/scope audits passed.
\subsection{The closed normal quotient by the descent-fixed subgroup}

\begin{statement}{C.41}{Proved}{Odd prime-power quotients after killing the fixed subgroup}
Put $G=\widehat{\mathrm{GT}}_0$ and let
$N_\varphi=\overline{\langle\!\langle C^\varphi\rangle\!\rangle}$
be the closed normal subgroup generated by the actual fixed subgroup.
There is a continuous surjection
\[
 G/N_\varphi\twoheadrightarrow
 C_2\times\prod_{\substack{p\geq5\ {
m prime}\\ k,j\geq1}}
                         \mathbb Z/p^k\mathbb Z.       \tag{C.41.1}
\]
For every $p\geq5$, its continuous $\mathbb F_p$ character
space has dimension exactly $\aleph_0$. In particular,
$C^\varphi$ has index $2^{\aleph_0}$ in $G$ and Haar
measure zero. No statement about the prime three part,
the complete abelianization, or a $\mathbb Z_p$ quotient
is implied. The inclusion $C^\varphi\subset\ker\Delta$
and the pentagon-to-fixedness question remain unresolved.
\end{statement}

\begin{proof}
We first make a locality point in C.38 explicit. Its
characters at a fixed prime $p$ are built on the actual
finite Magnus image of the local gentle group, together
with its local multiplier. The invariant homomorphism
of the multiplier kernel is extended across
$\mathbb Z_p^\times$, and the arithmetic correction
is a multiple of
\[
                  \log(\lambda_p^{p-1})/p\pmod{p^k}.
\]
Thus every character in that construction factors
through the single normalized local gentle group
$G_{0,p}$ of C.40. No other local coordinate or
global representative is used in its definition.

For fixed $p,k$, the construction gives arbitrarily
many jointly independent reductions on the actual
local multiplier kernel $U_p$. One may therefore
perform the final mixed-order selection of C.38
entirely inside these local characters: at each
step choose a character whose reduction on $U_p$
lies outside the span of all earlier reductions.
The relevant spaces have infinite dimension, so
this remains possible for every pair $(k,j)$.
For any finite collection, independence modulo $p$
gives $H+pT=T$ for its image $H$ in the finite
target $T$. Finite Nakayama gives $H=T$.
Compactness then gives actual continuous local
characters
\[
 \theta_{p,k,j}:G_{0,p}\longrightarrow\mathbb Z/p^k\mathbb Z
                                                        \tag{C.41.2}
\]
whose joint restriction to $U_p$ maps onto the
whole product over $k,j$. Their global inflations
kill the specified Galois image, as in C.38.

For each of these characters define on the full group
\[
 \eta_{p,k,j}(g)=
       \theta_{p,k,j}(g_p)-\theta_{p,k,j}(D(g)_p).
                                                        \tag{C.41.3}
\]
Both terms are continuous additive characters because
$D$ is a homomorphism. Equation (C.41.3) vanishes
pointwise on $C^\varphi=\operatorname{Fix}(D)$, hence
on its closed normal closure.

To check joint surjectivity, prescribe any finite
collection of target values. For each coefficient
prime involved, choose $\alpha_p\in U_p$ giving
those values under (C.41.2), and prescribe the
second local action to be the identity. C.40
supplies an actual full multiplier-one element
$g$ with $g_p=\alpha_p$ and $D(g)_p=1$ at all
these primes. Its values in (C.41.3) are precisely
the prescribed values. Compactness proves
surjectivity onto the entire odd-primary product
in (C.41.1).

There is also the quadratic character
\[
                       \eta_2=\Delta+\Delta\circ D.
                                                        \tag{C.41.4}
\]
It vanishes on the same fixed subgroup and its
closed normal closure. The actual element of
C.23 has $(\Delta(g),\Delta(D(g)))=(1,0)$,
so (C.41.4) is onto. For every finite target
subproduct, its quadratic factor and its odd
primary factors have coprime orders. Surjectivity
to each primary factor therefore gives
surjectivity to their product. A final compactness
argument proves (C.41.1).

The displayed quotient supplies infinitely many
independent mod-$p$ characters for every $p\geq5$.
The full group is countably based, as used in
C.38, so it has only countably many continuous
maps to a fixed finite group. This gives the
matching upper bound $\aleph_0$ for the dimension.

The product in (C.41.1) has cardinality
$2^{\aleph_0}$. Since $C^\varphi\subset N_\varphi$,
the coset space $G/C^\varphi$ maps onto this
quotient and has at least that cardinality.
The countably based compact group $G$ has at
most that cardinality. Its fixed subgroup
therefore has exactly the asserted index.
For Haar measure, arbitrarily large finite
quotients in (C.41.1) put $C^\varphi$ in
subgroups of arbitrarily large finite index.
Its measure is consequently zero.

This argument uses difference characters. It
does not identify $\Delta$ itself on the fixed
subgroup, and it does not say that the new
difference characters vanish on the full
pentagon group. Those implications would
require additional information about $D$ on
that subgroup.
\end{proof}


% C.42 candidate for integration after root's final proof read.
\begin{statement}{C.42}{Proved}{The three-primary part of the iterated finite lifting obstruction vanishes}
Let $G$ be the actual finite group of order $2^{146}3^{30}$ in the
retained iterated-descent calculation, with its original $176$-generator
pc presentation.  The complete presentation, the two modular generators
$s,r$, their prescribed images $a,b$, and the left peripheral frames
$B_s,B_r$ are specified in
\texttt{pbc\_iterated\_shortcut\_v2\_original\_source.json}.
Write $G_n$ for the actual quotient obtained by killing pc generators
with index greater than $n$, for the normal cuts used below.
Then the actual quotient map induces an isomorphism
\[
 H_2(G,\mathbb Z_3)\xrightarrow{\ \sim\ }H_2(G_{26},\mathbb Z_3),
 \qquad 9H_2(G,\mathbb Z_3)=0.
 \tag{C.42.1}
\]
Together with the previously retained two-primary bound, this gives the
integral annihilator
\[
                 2^{14}3^2H_2(G,\mathbb Z)=0.
 \tag{C.42.2}
\]

Let $L$ be the free group on the original pc generators, let
$R=\ker(L\to G)$, and put $K=R/[L,R]$.  The central framed-product
defect $T\in K$ is the sum of the three central differences
\[
 (B_ss)-(aB_s),\qquad (B_rr)-(bB_r),\qquad (ab)-(sr).
 \tag{C.42.3}
\]
Its three-primary component in the Hopf subgroup is zero:
\[
                T_{(3)}=0\quad\hbox{in }H_2(G,\mathbb Z)_{(3)}.
 \tag{C.42.4}
\]
No boundary correction is needed for this component.
This is a proved component of the iterated finite lifting problem.
The two-primary defect and the full gentle lift are treated separately
in C.31, which combines this component with the exact comparison of C.43.
\end{statement}
\begin{proof}
The source is identified throughout by its SHA256 hash
\texttt{95fb3845\allowbreak{}93bfaf97\allowbreak{}d468bbc5\allowbreak{}1f11cd79\allowbreak{}7484a08c\allowbreak{}b3ec3b95\allowbreak{}4961085b\allowbreak{}520ffc88},
retained with the executed certificates.  Its provenance is the actual
finite word image, rather than a group selected only by its order or
composition factors.  The new checker reads all $176$ original power
relations and all $15\,400$ original conjugation relations.  It verifies
the normal suffixes and obtains the following actual quotient maps:
\[
\begin{array}{c|c}
 \text{map}&\text{kernel}\\\hline
 G_{176}\to G_{104}&(C_2)^{72}\\
 G_{104}\to G_{86}&(C_3)^{18}\\
 G_{86}\to G_{26}&\text{a }2\text{-group of order }2^{60}\\
 G_{26}\to G_{17}&(C_3)^9\\
 G_{17}\to G_4&\text{a }2\text{-group of order }2^{13}.
\end{array}
 \tag{C.42.5}
\]
The two unspecified $2$-groups need not be abelian.  Each elementary
abelian assertion is checked from every power and internal conjugation
relation in its actual quotient.  Quotienting by a finite normal
$2$-group induces an isomorphism on all homology with $\mathbb Z_3$
coefficients, by the homology Hochschild--Serre spectral sequence.

\paragraph{The actual module calculations.}
For $A_d=\ker(G_h\to G_l)$ with $(l,h,d)=(17,26,9)$ or
$(86,104,18)$, take the original pc generators with indices
$l+1,\ldots,h$ as the $\mathbb F_3$-basis.  The action matrices are
read directly from their original conjugation words.  For $d=9$, the
normal $2^{13}$ subgroup of $G_{17}$ has coinvariants zero on $A_9$;
for $d=18$, the normal $2^{60}$ subgroup of $G_{86}$ has coinvariants
zero on $A_{18}$.  In both cases this is a full-rank calculation of
the stacked original matrices $M_g-I$, with ranks $9$ and $18$.
Since the normal subgroup has order prime to three, averaging and
Hochschild--Serre give
\[
                   H_i(G_l,A_d)=0\qquad(i\ge0).
 \tag{C.42.6}
\]

For the exterior-square module, the action rows are the actual
two-by-two minors of those same matrices.  The complete rank data are
\[
\begin{array}{c|r|r|r|r}
 d&\dim\Lambda^2A_d&\text{rows}&\text{rank}
                   &\dim(\Lambda^2A_d)_{G_l}\\\hline
 9&36&612&35&1\\
 18&153&13\,158&153&0.
\end{array}
 \tag{C.42.7}
\]
Every row is checked.  The certificate additionally retains the selected
independent original rows and verifies their ranks with a second
plain-Python elimination.  In particular, the zero in the last entry
is an exact finite certificate, not an inferred representation type.

For an elementary abelian three-kernel $A$ with quotient $Q$, the
total-degree-two terms of the homology spectral sequence are
\[
 H_2(Q,\mathbb Z_3),\qquad H_1(Q,A),\qquad
                         H_0(Q,\Lambda^2A).
 \tag{C.42.8}
\]
Here $H_2(A,\mathbb Z_3)=\Lambda^2A$ follows from the integral
abelian-group identity.  For $A_{18}$ the last two terms vanish.
The only possible outgoing differential from the first term is
$d_2$ to $A_{18,Q}$, which vanishes by (C.42.6); there is no incoming
differential and no later outgoing differential.  Thus the actual map
$G_{104}\to G_{86}$ induces an isomorphism on $H_2(-,\mathbb Z_3)$.
Combining this with the two normal $2$-kernel maps proves the
isomorphism in (C.42.1).

For $A_9$, the middle term in (C.42.8) vanishes and the last is killed
by three.  Its spectral-sequence filtration therefore multiplies a
bound for $H_2(G_{17},\mathbb Z_3)$ by at most three.  No splitting
of that filtration is assumed.

\paragraph{An elementary bound for the actual order-$54$ quotient.}
The four-generator quotient has relative orders $(3,2,3,3)$ and
trivial power words.  Its actual subgroup $A=\langle g_3,g_4\rangle$
is $C_3^2$, with quotient $C_6$.  On $A$, the actions of $g_1,g_2$
are respectively
\[
             U=\begin{pmatrix}1&0\\1&1\end{pmatrix},
             \qquad -I.
 \tag{C.42.9}
\]
These matrices and the quotient presentation are checked directly in
the original relations.  A generator of $C_6$ acts by $-U$, so its
action minus the identity is invertible over $\mathbb F_3$.
The cyclic resolution gives $H_1(C_6,A)=0$, while
$H_2(C_6,\mathbb Z_3)=0$.  Consequently $H_2(G_4,\mathbb Z_3)$ is
a subquotient of $(\Lambda^2A)_{C_6}$ and is killed by three.
The normal $2^{13}$ comparison and the $A_9$ filtration now give
$9H_2(G_{26},\mathbb Z_3)=0$.  This proves (C.42.1).
For completeness, the retained actual derived-series calculation has
successive abelian-factor exponents
\[
                         6,\ 3,\ 4,\ 6,\ 4,\ 3,\ 2.
\]
For an abelian normal kernel of exponent $r$, the last two terms of
(C.42.8), with that kernel in place of $A$, are killed by $r$.
The total-degree-two filtration therefore multiplies an annihilator
of the quotient's multiplier by at most $r^2$.  Induction through
the derived series gives the annihilator
$(6\cdot3\cdot4\cdot6\cdot4\cdot3\cdot2)^2=2^{14}3^8$.
The complete derived generators and their actual abelian invariants
are retained in \texttt{pbc\_derived\_h2\_bound.json}.
In particular the two-primary bound is $2^{14}$.  Combining it
with the sharper three-primary calculation, and using that finite-group
positive homology has no other primary parts here, proves (C.42.2).

An independently executed GAP calculation of the actual order-$54$
presentation returns its full multiplier $C_3$.  The elementary
argument above suffices for the bound, so it does not depend on
that software result.  The relevant documented multiplier and
source-centralization conventions are
\cite[\S7.9.2,\S7.9.4]{PBC_Polycyclic2026}.

\paragraph{A literal original-row identity for the defect.}
The quotient has $26$ pc generators and $351$ formal central tails.
All three terms in (C.42.3), and all original finite centralizer
words, are recollected in this actual quotient presentation.
Raw tail vectors from the $176$-generator collector are not
truncated: only the group words are projected before recollection.
In the resulting original tail coordinates the target is
\[
                             T_{26}=e_{12}+e_{24}.
 \tag{C.42.10}
\]
The exact search modulo nine uses only the nine genuine
$r$-centralizer commutators.  Its $16\,278$-node arithmetic DAG
is independently expanded into original leaves.  Every reduced
coordinate bound $g=\gcd(d,9)$ is justified by explicit integers
$a,b$ with $ad+9b=g$, where $d e_i$ is itself a retained integer
combination of original power-consistency identities.

The final displayed certificate removes every coordinate and
modulus leaf before forming its integer sum.  It then applies the
explicit original power identities just described.  The exported
ledger contains exactly $440$ original identity rows and one
genuine commutator row, and verifies the literal equality
\[
       T_{26}+3c_{74}+\sum_{j=1}^{440} n_j R_j
                              =9z\quad\text{in }\mathbb Z^{351}.
 \tag{C.42.11}
\]
Every $n_j$, every original row $R_j$, and all $351$ coordinates
of $z$ are retained in
\texttt{pbc\_\allowbreak{}quotient26\_\allowbreak{}literal\_\allowbreak{}original\_\allowbreak{}identity.json}.
There are no coordinate or modulus rows in this final witness.
The standard GAP collector independently replays all $442$ rows
appearing in it, including the target and the commutator, in the
original unprocessed presentation.  A separate earlier replay
checks the full set of $443$ used leaves before consolidation.
The source hashes and both complete successful outputs are retained.

The index $74$ is the first of the nine original $r$-centralizer
corrections; the two power-word positions and every $s$-correction
have coefficient zero.  In the actual source centralization,
$c_{74}$ is central and has exponent dividing three, because it
is the commutator of a lift of an element centralizing $r$ with a
lift of $r$, and $r^3$ is central.  Thus $3c_{74}=0$ integrally.
Equation (C.42.11) proves $T_{26}\in9K_{26}$.

Finally the image of (C.42.3) in the free abelianization is zero:
the three differences telescope.  The Hopf exact sequence therefore
places $T$ and $T_{26}$ in the respective finite torsion subgroups
$H_2(G,\mathbb Z)$ and $H_2(G_{26},\mathbb Z)$ of the centralization
kernels.  Since the three-primary subgroup is killed by nine, its
map to $K_{26}/9K_{26}$ is injective.  Indeed a finitely generated
abelian group splits abstractly as a free group plus its finite
torsion subgroup, and nine kills precisely the relevant component.
Hence $(T_{26})_{(3)}=0$.  Hopf naturality identifies this with the
image of $T_{(3)}$ under the actual isomorphism (C.42.1), proving
(C.42.4).  The remaining two-primary component is not addressed
by this last inference.
\end{proof}


% C.43 candidate for canonical integration after final root read.
\begin{statement}{C.43}{Proved}{Exactly two binary classes remain above the smaller finite quotient}
Use the actual finite source and normal quotients of C.42.  Put
\[
 G=G_{176},\qquad Q=G_{104},\qquad
 A=\ker(G\to Q)\simeq(C_2)^{72}.
\]
The actual finite extension $1\to A\to G\to Q\to1$ splits.
The natural quotient map to $G_{86}$ induces a split exact sequence
\[
 0\longrightarrow(C_2)^2\longrightarrow H_2(G,\mathbb Z_2)
   \longrightarrow H_2(G_{86},\mathbb Z_2)\longrightarrow0.
 \tag{C.43.1}
\]
Its kernel is the actual image of
$H_0(Q,\Lambda^2A)$ and has exactly the displayed order.
The splitting depends on a choice of group section; the natural map
and its kernel do not.

Consequently a verified two-primary lifting calculation in $G_{86}$
leaves two binary values to determine in the original source.
Neither value is computed by this statement. C.31 supplies both
actual cover evaluations and the full iterated-descent lift.
\end{statement}
\begin{proof}
The original pc generators $g_{105},\ldots,g_{176}$ give the stated
basis of $A$.  The actual normal subgroup
\[
       N=\ker(Q\to G_{86})=\langle g_{87},\ldots,g_{104}\rangle
                            \simeq(C_3)^{18}
 \tag{C.43.2}
\]
has odd order.  These normality and elementary abelian assertions
are checked from every original relation, as in C.42.

\paragraph{An exact rank certificate on the original module.}
For every one of the $104$ quotient generators, read its action on
$A$ directly from the original conjugation words.  The matrices
are invertible over $\mathbb F_2$.  The stacked rows $M_g-I$ for
the $18$ generators in (C.43.2) have rank $72$, proving
\[
                              A_N=0.
 \tag{C.43.3}
\]
The exterior square has dimension $\binom{72}{2}=2556$.
Stacking the exterior action minus the identity for all $104$
quotient generators gives $265\,824$ actual rows and rank $2554$.
Hence
\[
                       (\Lambda^2A)_Q\simeq(C_2)^2.
 \tag{C.43.4}
\]

Both rank inequalities in (C.43.4) have small exact certificates.
The full program retains $2554$ independent original action rows.
A separate checker reconstructs them from the original presentation
and eliminates them with lowest-coordinate pivots, whereas the first
calculation used highest-coordinate pivots.  This independently
proves rank at least $2554$.  The checker also reconstructs two
explicit independent alternating forms on $A$ and verifies their
invariance under every one of the $104$ original matrices.  They
give rank at most $2554$.  Thus no unprocessed part of the full
row system is presumed to vanish.

The full data are in
\texttt{pbc\_two\_top\_module.json}; the fast independent
certificate is
\texttt{pbc\_two\_top\_certificate.py} and its retained output.
The second checker rederives all action matrices from the source
hash of C.42, rather than trusting only the first program's rank.
The independent audit also checks the literal normal suffixes and
their elementary kernels in the original group.

\paragraph{The actual extension splits.}
Because $|N|$ is odd, averaging is defined on the exponent-two
module $A$.  It identifies invariants with coinvariants and makes
both functors exact.  Equation (C.43.3) therefore gives
\[
 A^N=0,\qquad H_i(N,A)=H^i(N,A)=0\quad(i>0).
\]
Apply both homology and cohomology Hochschild--Serre to the actual
normal subgroup $N$ in $Q$.  Every second-page coefficient vanishes,
including degree zero, so
\[
                    H_i(Q,A)=H^i(Q,A)=0\qquad(i\ge0).
 \tag{C.43.5}
\]
In particular $H^2(Q,A)=0$.  Classification of extensions by this
cohomology group gives an actual homomorphism section $Q\to G$
\cite[II, Example~1.18(b), Remark~1.35(b)]{PBC_MilneCFT}.
This proves the group-level splitting; no integral Lie-algebra or
associated-graded splitting is substituted for it.

\paragraph{All possible differentials in degree two.}
For $A$ normal in $G$, use the homology spectral sequence
\[
 E^2_{a,b}=H_a(Q,H_b(A,\mathbb Z_2))
                              \ \Longrightarrow H_{a+b}(G,\mathbb Z_2).
 \tag{C.43.6}
\]
As the groups are finite, these coefficients retain exactly the
two-primary part of their integral homology.  The integral identity
$H_2(A,\mathbb Z)=\Lambda^2A$ identifies the left edge in total
degree two.  By (C.43.4)--(C.43.5),
\[
 E^2_{0,2}=(C_2)^2,\qquad E^2_{1,1}=0,\qquad
                         E^2_{2,0}=H_2(Q,\mathbb Z_2).
 \tag{C.43.7}
\]
There are only two possible incoming differentials to $E_{0,2}$.
The $d_2$ source is $H_2(Q,A)=0$.  The $d_3$ source is in the
bottom row $E_{3,0}$, and this differential is zero because the
actual section makes every bottom-row class permanent: projection
$H_i(G,\mathbb Z_2)\to H_i(Q,\mathbb Z_2)$ is split surjective.
There are no later incoming differentials and no outgoing
differentials from the left edge.

The same section protects $E_{2,0}$.  Thus the filtration gives
\[
 0\longrightarrow(C_2)^2\longrightarrow H_2(G,\mathbb Z_2)
        \longrightarrow H_2(Q,\mathbb Z_2)\longrightarrow0,
\]
and the section splits this exact sequence as well.  This proves
that the kernel is exactly $(C_2)^2$, rather than merely a possible
quotient of that group.  Finally $Q\to G_{86}$ has the odd-order
kernel $N$, so its homology spectral sequence gives a natural
isomorphism on $H_2(-,\mathbb Z_2)$.  Composing proves (C.43.1).

\paragraph{What the finite calculation does not decide.}
The two explicit invariant alternating forms both vanish on
$a\wedge s(a)$ for every basis vector $a$ of $A$, where
$s=g_2$ is the actual involution.  The map
$a\mapsto B(a,s(a))$ is additive: its cross terms cancel because
$s$ is an involution preserving $B$.  Thus these values vanish for
every $a\in A$.  All norms $a+s(a)$ lie in the span of the $36$
retained top-kernel centralizer generators.  This particular norm
calculation therefore does not provide the two missing values.

A separate modular search in $G_{86}$ has been retained, including
its complete proof records.  Its first independent original-source
replay rejects the raw target normal form.  It supplies no verified
quotient identity and is not used anywhere in this proof.  Even a
successful quotient replay would still require the two values in
(C.43.1) before a conclusion about the original finite obstruction.
C.31 uses a new standard-GAP-row calculation that passes that replay,
and evaluates both coordinates in actual central covers; the failed
attempt above remains outside its mathematical argument.
\end{proof}
